\chapter{分割与场景理解}
\label{ch:10}

\paragraph*{本章阅读锚点}
本章的锚点是语义绑定。原始 3DGS 只存几何和颜色；理解方法把语言、分割、实例或特征嵌入高斯，使场景可以被查询、选择和推理。


3D Gaussian Splatting 用数百万个各向异性高斯基元精确重建了场景的外观和几何，但这些高斯本身没有语义——它们不知道自己属于地面、墙壁还是椅子。这种"看得见但看不懂"的困境，正是 3DGS 场景理解需要解决的核心问题。赋予高斯基元语义信息面临三重挑战。第一，高斯基元是离散的体积元素，不存在天然的对象边界，一个横跨前景和背景的大高斯会让分割边界模糊不清。第二，语义信息通常来自 2D 基础模型（如 SAM、CLIP），但 2D 掩码在不同视角间往往不一致，直接提升到 3D 会产生噪声累积。第三，3DGS 的光栅化过程只优化光度一致性，没有任何语义约束，语义场需要与外观场解耦或联合优化。

本章系统梳理 3DGS 场景理解的多个方向。第~\ref{sec:10_1}~节覆盖 3D 分割与语义理解。第~\ref{sec:10_3}~节聚焦开放词汇场景理解与特征蒸馏。第~\ref{sec:10_2}~节讨论目标检测、场景图与场景补全。第~\ref{sec:10_7}~节介绍变化检测、异常检测等其他场景理解方法。

%=====================================================================
\begin{figure}[htbp]
\centering
\resizebox{\textwidth}{!}{\begin{tikzpicture}[
  route/.style={rectangle, rounded corners=5pt, draw=#1, fill=#1!10,
    text width=3.2cm, minimum height=0.9cm, align=center, font=\small},
  method/.style={rectangle, rounded corners=3pt, draw=#1, fill=#1!6,
    font=\scriptsize, text width=3.0cm, align=center, minimum height=0.6cm},
  arr/.style={-{Stealth[length=4pt]}, techgray!60},
  >=Stealth
]

  \node[route=techgray, text width=4cm, font=\small\bfseries] (root) at (30.0, 5.6) {分割与场景理解方法分类};
  \node[route=techblue] (b0) at (11.0, 3.1) {3D分割};
  \node[method=techblue] (b0_l0) at (1.5, 0.6) {特征蒸馏\\{\tiny Gaussian Grouping, SAGA, FlashSplat, COB-GS, GradiSeg}};
  \node[method=techblue] (b0_l1) at (5.1, 0.6) {开放词汇\\{\tiny LangSplat, OpenGaFF, SCOUP, ReferSplat, PanoGS}};
  \node[method=techblue] (b0_l2) at (9.0, 0.6) {交互式/点击\\{\tiny Click-Gaussian, GaussianCut, ArtisanGS, B$^3$-Seg}};
  \node[method=techblue] (b0_l3) at (13.0, 0.6) {全景分割\\{\tiny PLGS, PanopticSplatting, OnlinePG, PanSt3R}};
  \node[method=techblue] (b0_l4) at (16.6, 0.6) {前馈泛化\\{\tiny RT-GS2, SemGS, SegSplat, OVGaussian}};
  \node[route=techgreen] (b1) at (26.2, 3.1) {目标检测与场景图};
  \node[method=techgreen] (b1_l0) at (20.5, 0.6) {3D目标检测\\{\tiny 3DGS-DET, Gaussian-Det, GVSynergy-Det}};
  \node[method=techgreen] (b1_l1) at (24.7, 0.6) {场景图生成\\{\tiny GaussianGraph, DynamicGSG}};
  \node[method=techgreen] (b1_l2) at (28.6, 0.6) {场景补全\\{\tiny VG3S, GaussRender, SplatSSC, GS-Occ3D}};
  \node[route=techorange] (b2) at (34.2, 3.1) {特征蒸馏};
  \node[method=techorange] (b2_l0) at (32.3, 0.6) {CLIP蒸馏\\{\tiny Feature 3DGS, LangSplat, CLIP-GS, FAST-Splat}};
  \node[method=techorange] (b2_l1) at (36.2, 0.6) {DINO蒸馏\\{\tiny LUDVIG, econSG, SLAG}};
  \node[method=techorange] (b2_l2) at (40.1, 0.6) {SAM融合\\{\tiny CUS-GS, CF3, FHGS, Quantile Rendering}};
  \node[route=techpurple] (b3) at (43.8, 3.1) {其他};
  \node[method=techpurple] (b3_l0) at (43.8, 0.6) {变化检测\\{\tiny GD-DIFF, 3DGS-CD}};
  \node[method=techpurple] (b3_l1) at (47.4, 0.6) {异常检测\\{\tiny SplatPose, SplatPose+, GS-MoE}};
  \draw[arr] (root) -- (b0);
  \draw[arr] (root) -- (b1);
  \draw[arr] (root) -- (b2);
  \draw[arr] (root) -- (b3);
  \draw[arr] (b0) -- (b0_l0);
  \draw[arr] (b0) -- (b0_l1);
  \draw[arr] (b0) -- (b0_l2);
  \draw[arr] (b0) -- (b0_l3);
  \draw[arr] (b0) -- (b0_l4);
  \draw[arr] (b1) -- (b1_l0);
  \draw[arr] (b1) -- (b1_l1);
  \draw[arr] (b1) -- (b1_l2);
  \draw[arr] (b2) -- (b2_l0);
  \draw[arr] (b2) -- (b2_l1);
  \draw[arr] (b2) -- (b2_l2);
  \draw[arr] (b3) -- (b3_l0);
  \draw[arr] (b3) -- (b3_l1);
\end{tikzpicture}}
\caption{分割与场景理解方法分类}
\label{fig:ch10_segmentation_taxonomy}
\end{figure}

\section{3D分割与语义理解}
\label{sec:10_1}

3D 分割旨在为场景中的每个高斯基元赋予语义标签或实例归属，是下游编辑、检测和推理的基础。本节涵盖约 66 篇论文，按技术路线组织为四个子主题：基于 2D-3D 特征蒸馏的分割、开放词汇分割、交互式/点击式分割和全景分割。

\subsection{基于特征蒸馏的3D分割}
\label{ss:10_1}

本节讨论3条技术路线：（1）基于2D-3D特征蒸馏的分割；（2）对比学习与身份编码路线；（3）特征蒸馏的改进路线。

\subsubsection{基于2D-3D特征蒸馏的分割}

本路线介绍将 2D 基础模型（SAM、CLIP、DINO（DINOv2，一种基于自监督视觉Transformer的特征提取模型））的分割结果通过可微渲染提升到 3D 高斯基元以实现语义分割的基本范式。

3DGS 场景缺乏语义标注，而 2D 基础模型（SAM、CLIP、DINO）已具备强大的分割和语义理解能力。最直觉的方案是将 2D 分割结果"提升"（lift）到 3D：从多个视角获取 2D 掩码，通过可微渲染将掩码信息反向传播到高斯基元上，为每个高斯分配语义标签或特征向量。这一思路简洁有效，但在实践中需要克服几个关键难题。

这一范式最早由 LERF（2023.03）~\cite{2303_09553} 在 NeRF 框架下提出，通过将 CLIP 特征蒸馏到 NeRF 的密度场中实现语言驱动的 3D 查询。LangSplat（2023.12）~\cite{2312_16081} 将这一思想迁移到 3DGS，为每个高斯基元附加语言特征，通过 3D 语言高斯光栅化实现开放词汇的 3D 查询。但 LERF 和 LangSplat 的分割精度有限，它们的特征场是连续的，难以生成锐利的对象边界。后续工作围绕三个核心问题展开：如何学习具有区分性的身份编码？如何消除分割边界模糊？如何从逐场景优化走向可泛化？

\subsubsection{对比学习与身份编码路线}

本路线介绍通过对比学习为高斯基元学习身份编码以实现实例分割的方法。

Gaussian Grouping（2023.12）~\cite{2312_00732} 提出了一个简洁而有效的框架，奠定了 2D-3D 特征蒸馏分割的基本范式。其核心思想是为每个高斯基元附加一个紧凑的 Identity Encoding（身份编码），通过可微渲染和多视角 SAM 掩码的对比监督学习实例归属。

\textbf{动机。}3DGS 缺乏语义结构，一个横跨前景和背景的大高斯会让分割边界模糊不清；通过为每个高斯学习身份编码并用 2D 掩码的对比监督驱动优化，可以将语义信息注入离散高斯基元。

具体地，给定 $N$ 个高斯基元，每个高斯附加一个 $D$ 维身份编码 $\mathbf{e}_i \in \mathbb{R}^D$（$D$ 通常为 16 或 32）。渲染时，身份编码通过标准的 $\alpha$-混合方式合成：

\begin{equation}
\label{eq:10_identity_render}
\hat{\mathbf{e}}(\mathbf{p}) = \sum_{i=1}^{N} \mathbf{e}_i \, \alpha_i \, T_i, \quad T_i = \prod_{j=1}^{i-1}(1 - \alpha_j)
\end{equation}

其中 $\alpha_i$ 和 $T_i$ 分别为第 $i$ 个高斯的不透明度和透射率，与颜色渲染共享相同的排序和混合逻辑。等式~\ref{eq:10_identity_render}~的关键洞察是：身份编码的渲染与颜色渲染完全同构，因此可以复用 3DGS 的高效光栅化管线，无需设计专门的语义渲染器。

监督信号来自 SAM 在每个视角生成的 2D 掩码：同一掩码内的高斯应具有相似的身份编码，不同掩码的高斯应相互区分。这一目标通过对比损失实现。此外，Gaussian Grouping 引入 3D 空间一致性正则化，鼓励空间上邻近的高斯共享身份编码，从而弥补 2D 监督的视角不一致性。基于 Gaussian Grouping 的分组结果，该方法还支持多种下游编辑应用，包括 3D 对象移除、修复、上色和场景重组。

与 Gaussian Grouping 同期，SAGA（2023.12，Scale-Aware Gaussian Attachment）~\cite{2312_00860} 提出了另一种身份编码学习策略。

\textbf{动机。}Gaussian Grouping 的身份编码粒度固定，无法同时处理从部件级到整体级的分割需求；SAGA 引入尺度感知的 soft scale gate 机制，使分割结果在不同粒度层级上保持一致性。SAGA 为每个高斯附加一个 scale-gated affinity feature，引入尺度感知的对比训练，通过一个 soft scale gate 机制调整特征通道的响应幅度，使分割结果在不同粒度层级（从部件到整体）上保持一致性。SAGA 可在 4ms 内完成一次 3D 提示分割，是首批实现 3DGS 实时交互分割的方法之一。

\subsubsection{特征蒸馏的改进路线}

本路线介绍在身份编码、边界处理、训练效率等方面对基础特征蒸馏范式的系统改进。

后续工作沿多个方向改进了 Gaussian Grouping 奠定的特征蒸馏范式。

\paragraph{UniC-Lift：单阶段统一优化} Gaussian Grouping 采用"先对比学习优化嵌入、后 HDBSCAN 聚类提取实例"的两阶段范式。UniC-Lift（2025.12）~\cite{2512_24763} 将两阶段流程统一为单阶段优化：引入可学习的 Embedding-to-Label 过程直接将特征嵌入解码为实例标签，避免了 HDBSCAN（Hierarchical Density-Based Spatial Clustering of Applications with Noise，一种基于密度的层次聚类算法）聚类 $O(n^2)$ 的计算开销和超参数敏感的聚类后处理。

\textbf{动机。}Gaussian Grouping 的两阶段范式（先优化嵌入再 HDBSCAN 聚类）中，HDBSCAN 的 $O(n^2)$ 开销和超参数敏感的聚类后处理成为效率和鲁棒性瓶颈；UniC-Lift 通过可学习的软标签向量直接在训练中形成类别可分的嵌入空间，无需后处理。其关键设计是为每个高斯分配一个可学习的软标签向量，通过交叉熵损失直接监督：同一 SAM 掩码内的高斯被推向相同标签，不同掩码的高斯被推离，使嵌入空间在训练过程中就形成类别可分的结构，无需后处理即可获得实例分配。该方法还在对象边界处应用 hard-mining 策略，通过在光栅化后的特征嵌入上计算 triplet loss 来稳定边界区域的训练。

\paragraph{FreeGS：免监督语义-实例联合框架} UniC-Lift 仍依赖 2D 掩码标签进行监督。FreeGS（2024.11）~\cite{2411_19551} 绕开这一依赖，提出无需 2D 标签的免监督框架：引入 IDentity-coupled Semantic Field（IDSF，身份耦合语义场），同时捕获每个高斯的语义表示和视角一致的实例索引，通过两步交替优化策略实现语义和实例的相互增强。

\textbf{动机。}UniC-Lift 仍依赖 2D 掩码标签进行监督，而 2D 掩码本身存在跨视角不一致问题；FreeGS 完全绕开 2D 标签依赖，通过语义和实例的相互增强实现免监督分割。语义帮助提取 3D 一致实例，实例反过来正则化 2D 语义注入。

\paragraph{AG$^2$aussian：锚点图结构替代全局聚类} 前述方法均在平坦的特征空间中学习实例编码。AG$^2$aussian（2025.08，Anchor-Graph Gaussian）~\cite{2508_01740} 从表示结构角度改进，用 anchor-graph 结构组织语义特征。

\textbf{动机。}Gaussian Grouping 在平坦的特征空间中进行全局聚类，忽略了高斯之间的空间邻接关系；AG$^2$aussian 用锚点图结构显式编码空间邻接，通过图传播实现更紧凑的实例选择。锚点作为语义簇中心取代 Gaussian Grouping 中的全局聚类，图结构编码高斯之间的空间邻接关系，通过图传播实现紧凑的实例级高斯选择，支持交互式点击查询、开放词汇文本驱动查询和物体移除编辑等多种应用。

\paragraph{LBG：无需训练的直接融合} 前述方法均需要逐场景优化。LBG（2025.02，Lift-Bind-Gaussian）~\cite{2502_00173} 将特征蒸馏范式推向极致效率：无需逐场景训练，直接将 SAM 的 2D 掩码和 CLIP/DINOv2（Deep Image Network Operators v2，一种基于自监督视觉Transformer的强鲁棒特征提取模型）的特征融合到 3DGS 上，跳过特征优化和聚类环节，比现有方法快一个数量级且高度模块化，可在任何已有 3DGS 重建上直接运行。

\textbf{动机。}前述方法均需要逐场景优化，限制了实际部署效率；LBG 通过直接融合跳过所有优化步骤，以极低的时间成本实现可接受的分割质量。

\paragraph{FlashSplat：闭合形式标签分配} 分割边界模糊是 3DGS 分割的顽疾，大尺寸高斯横跨对象边缘导致锯齿状分割。FlashSplat（2024.09）~\cite{2409_08270} 从数学上重新审视了这个问题，将分割标签分配从迭代优化转化为闭合形式求解。

\textbf{动机。}Gaussian Grouping 的迭代优化需要数千步收敛且可能陷入局部最优；FlashSplat 的核心洞察是：在已重建的 3DGS 中，$\alpha$-混合 权重均为已知常数，因此分割标签分配可转化为线性整数规划问题，通过加权多数投票直接获得全局最优解。在 FlashSplat 之前，Gaussian Grouping 等方法通过梯度下降迭代优化每个高斯的标签，需要数千步才能收敛，且可能陷入局部最优。FlashSplat 的核心洞察是：给定已重建的 3DGS 场景，$\alpha$ 值 $\alpha_i$ 和透射率 $T_i$ 均为已知常数，因此 2D 掩码渲染可写成高斯标签 $P_i \in \{0,1\}$ 的\textbf{线性函数}。直觉上，像素被标记为前景的渲染值恰好是所有前景高斯在该像素处的 $\alpha$-混合贡献之和。据此，FlashSplat 将分割建模为整数线性规划（ILP，Integer Linear Programming）：以平均绝对误差为目标函数 $\min_{P_i} F = \sum_v \sum_{jk} | \sum_i P_i \alpha_i T_i - M_{jk}^v |$。由于 $\alpha_i \geq 0$ 且 $\sum_i \alpha_i T_i \leq 1$，该 ILP 可通过加权多数投票直接求解。对每个高斯 $G_i$，统计其在所有掩码中属于前景的加权贡献 $A_1^i = \sum_{v,j,k} \alpha_i T_i \cdot \mathbb{I}(M_{jk}^v = 1)$ 和背景的贡献 $A_0^i$，取 $\arg\max_n A_n^i$ 即为全局最优分配。此外，FlashSplat 引入可调的背景偏置 $\gamma \in [-1,1]$（L1 归一化后执行 $P_i = \arg\max\{\hat{A}_0, \hat{A}_1\}$，其中 $\hat{A}_0 = \hat{A}_0 + \gamma$），以缓解 SAM 掩码噪声导致的边界伪影。整个优化在约 30 秒内完成（约 $50\times$ 加速），且解是全局最优的。

\paragraph{COB-GS：边界自适应高斯分裂} FlashSplat 证明了标签分配可以一步到位，但并未解决边界高斯本身的几何歧义：横跨对象边缘的大尺寸高斯即使被正确标记，渲染出的掩码边界仍然模糊。COB-GS（2025.03，Clean Object Boundaries via Gaussian Splitting）~\cite{2503_19443} 正是从这一局限出发，将优化焦点从"标签分配"推进到"高斯几何改造"。

\textbf{动机。}FlashSplat 虽能一步完成标签分配，但横跨对象边缘的大尺寸高斯即使被正确标记，渲染出的掩码边界仍然模糊；COB-GS 将语义梯度异常高的"边界歧义高斯"沿语义边界分裂为两个子高斯，使高斯分布贴合对象轮廓。COB-GS 为每个高斯引入连续掩码标签 $m_i \in (0,1)$，通过 $\alpha$-混合渲染后与 SAM2 生成的 2D 掩码做监督训练。其核心操作是 boundary-adaptive Gaussian splitting（边界自适应高斯分裂）：在训练过程中统计每个高斯的掩码标签梯度 $\nabla m_i$，梯度异常高的高斯即为语义上不确定的"边界歧义高斯"，它们同时被前景和背景掩码施加了矛盾信号。COB-GS 将这些歧义高斯沿语义边界分裂为两个子高斯，分别独立学习前景和背景的掩码标签，使高斯分布贴合对象轮廓。与 FlashSplat 直接删除边界高斯不同，COB-GS 同时联合优化语义和纹理（外观优化修复分裂后的几何退化），在 NVOS 数据集上以相当的运行时间实现了显著更高的视觉质量。

\paragraph{BEA-GS：双路损失解决遮挡} COB-GS 的语义边界对齐损失仍然经过光栅化器传播梯度，对于严重遮挡的高斯，由于透射率 $T_i$ 极小导致梯度信号衰减，调整效果有限。BEA-GS（2026.05，Boundary-Enhanced Alignment Gaussian Splatting）~\cite{2605_09662} 针对这一问题设计了两组专用损失函数来分别处理可见和不可见高斯。

\textbf{动机。}COB-GS 的边界对齐损失经光栅化器传播梯度，对于被严重遮挡的高斯，透射率极小导致梯度信号衰减；BEA-GS 的核心创新是为不可见高斯设计绕过光栅化器的直接 3D 空间损失。可见高斯的语义边界对齐损失通过标准光栅化管线传播，直接优化对象表面的高斯几何使其贴合语义边界；而非可见高斯的几何调整损失\textbf{绕过光栅化器}，直接在 3D 空间中对被遮挡的高斯参数施加梯度，使其在对象被提取后（遮挡消失）仍能呈现正确的几何边界。这种双路损失设计使 BEA-GS 在对象提取任务中实现了近乎完美的边界精度，在 4 个数据集上超越 12 种基线方法。

\paragraph{GradiSeg：训练过程中的身份梯度引导} GradiSeg（2024.12）~\cite{2412_00392} 从另一个角度切入边界问题，将 FlashSplat 和 COB-GS 的后处理范式前移到训练过程中。

\textbf{动机。}FlashSplat 和 COB-GS 在训练完成后才处理边界歧义，无法在训练过程中利用边界信息改进特征学习；GradiSeg 通过监控身份编码的累积梯度在训练中期主动分裂歧义高斯，从根源上消除优化冲突。GradiSeg 为每个高斯附加一个 16 维的 Identity Encoding 向量，用于编码语义身份并通过分类网络渲染为分割结果。其关键创新在于两个互补模块。Identity Gradient Guided Densification（IGD，身份梯度引导致密化）在训练中期（12K--15K 迭代）监控身份编码的累积梯度，梯度异常高的高斯被判定为边界歧义区域，它们同时受到不同语义方向的监督信号，随后将其分裂为两个子高斯分别学习各自的语义身份，从根源上消除优化冲突。Local Adaptive KNN（LA-KNN，局部自适应 K 近邻）则在训练后期（15K--30K 迭代）接替全局 KNN，通过计算高斯位置梯度来建立方向感知的局部邻域，使身份编码只在\textbf{同侧}邻居之间传播，防止特征跨边界扩散污染。相比 COB-GS 的语义梯度统计，GradiSeg 的 IGD 直接作用于身份编码梯度，语义信号更为直接；LA-KNN 则弥补了 COB-GS 分裂操作后缺乏局部特征一致性约束的不足。

\paragraph{GaussianTrimmer：即插即用后处理修剪} 上述方法均在训练或优化过程中改造高斯结构，而 GaussianTrimmer（2026.01）~\cite{2601_12683} 则另辟蹊径，提出即插即用的后处理边界修剪方案，可直接应用于任何已有分割方法的结果之上。

\textbf{动机。}COB-GS 和 GradiSeg 需要在训练循环中修改高斯结构，无法应用于已完成的分割结果；GaussianTrimmer 作为纯后处理方法无需重训练，以极低的时间成本改善分割边界质量，具有极强的实用性。其核心流程分为两步。首先在对象周围生成均匀分布的虚拟相机阵列，确保从各角度充分覆盖对象边界；然后对每张虚拟相机视角执行 2D 分割，根据分割结果在基元级别修剪跨边界的过大高斯。与 COB-GS 和 GradiSeg 需要在训练循环中修改高斯不同，GaussianTrimmer 作为纯后处理方法无需重训练，以极低的时间成本即可改善分割边界质量，具有极强的实用性。

\paragraph{多层语义与层次化理解} 真实场景的语义具有层次结构，"椅子"包含"椅背"和"椅腿"，而传统方法如 Gaussian Grouping 只提供扁平的单层分割。TreeGaussian（2026.04）~\cite{2604_03309} 引入树状结构建模这种层次关系。首先利用 SAM 生成的多级掩码构建多视角 2D 对象树（object tree），通过细化重叠和间隙掩码建立精确的层次父子关系。然后设计两阶段级联对比学习策略：第一阶段在粗粒度掩码上优化全局实例特征以区分不同对象，第二阶段在细粒度掩码上精化局部特征以区分同一对象内的部件，逐步从全局到局部精化特征表示。最后设计 Consistent Segmentation Detection（CSD，一致性分割检测）机制动态平衡不同视角间的对比监督权重，确保跨视角分割模式的一致性，并辅以基于图的去噪模块移除聚类中的噪声点。InstanceGaussian（2024.11）~\cite{2411_19235} 从表示密度角度指出 Gaussian Grouping 将外观和语义绑定在同一高斯集上造成的三个关键挑战：外观与语义的密度不平衡（纹理建模需要密集高斯但语义只需少量，导致语义特征被冗余高斯稀释）、外观与语义的不一致（纯外观高斯常误标对象边界）、以及自顶向下实例分割的类别偏差。为此提出 Semantic-Scaffold-GS 表示，借鉴 Scaffold-GS 的锚点-子高斯层次结构但为语义引入独立的锚点网络，使语义高斯密度与外观高斯密度解耦：语义锚点覆盖更稀疏但每个锚点承载更丰富的语义信息，有效平衡外观和语义密度。配合渐进式外观-语义联合训练策略（外观先收敛提供稳定几何，语义随后优化避免与外观梯度冲突）和自底向上的类别无关实例聚合（通过最远点采样初始化簇中心，再用连通分量分析合并空间邻接的高斯为完整实例），在 ScanNet 上显著超越 OpenGaussian 和 Gaussian Grouping。

\paragraph{NG-GS：NeRF-高斯混合弥补语义空洞} 上述方法均依赖 3DGS 的离散点基表示，当高斯分布稀疏或边界区域覆盖不足时，语义特征场会出现不连续的空洞。NG-GS（2026.04，NeRF-Gaussian Splatting）~\cite{2604_14706} 引入轻量 NeRF 模块弥补这一缺陷。

\textbf{动机。}纯 3DGS 的离散点基表示在高斯分布稀疏或边界覆盖不足时，语义特征场会出现不连续的空洞；NG-GS 引入轻量 NeRF 模块通过 RBF 插值构建空间连续的特征场，弥补离散高斯之间的语义间隙。首先通过掩码方差分析自动识别边界歧义高斯，统计每个高斯在多视角掩码中被标记为前景和背景的频率差异，差异大的即为边界高斯。然后以这些边界高斯为锚点，用径向基函数（RBF，Radial Basis Function）插值在其周围构建空间连续的特征场，使离散高斯之间获得平滑的语义过渡。同时通过多分辨率哈希编码（MRHE，Multi-Resolution Hash Encoding）为 RBF 提供多尺度空间输入，增强表示容量。最后通过联合优化策略对齐 3DGS 和 NeRF 的输出，利用对齐损失和空间连续性损失确保分割兼顾高频细节和跨视角平滑。

\paragraph{NRGS：方差感知特征修正} NRGS（2026.04，Noise-Reduced Gaussian Splatting）~\cite{2604_22439} 从正则化角度解决同一问题：跨视角 2D 语义特征的不一致性导致直接提升到 3D 后语义场噪声严重。

\textbf{动机。}LangSplat 的迭代训练蒸馏耗时且可能过拟合，OccamLGS 的加权平均虽快但无法区分可靠与不可靠视角的信号；NRGS 采用训练无关的特征聚合加轻量 MLP 修正，以方差加权优先从多视角信号可靠的区域学习。不同于 LangSplat 需要迭代训练蒸馏特征、OccamLGS 仅做加权平均，NRGS 采用训练无关的特征聚合方式先将 2D 特征直接映射到 3D 高斯，然后引入轻量 variance-aware conditional MLP 利用高斯的固有属性（位置、协方差、颜色、不透明度）作为条件输入来修正 3D 空间中的语义误差。该 MLP 以方差加权为目标函数，优先从多视角信号可靠（方差低）的区域学习，抑制来自不一致视角的噪声信号，无需额外预处理且计算开销极小。

\paragraph{Semantic Foam：Voronoi 网格语义表示} Semantic Foam（2026.04）~\cite{2604_26262} 则从根本上更换了场景表示——将 Radiant Foam 的 Voronoi 网格表示从纯几何重建扩展到语义领域。

\textbf{动机。}3DGS 的点基表示中高斯之间存在遮挡和跨视角不一致伪影，难以在相邻高斯间施加空间正则化；Voronoi 网格的凸多面体单元天然提供空间分区，可直接在相邻单元间施加正则化约束。Voronoi 图将空间划分为不重叠的凸多面体单元，每个单元关联一个原始顶点，天然提供空间分区。Semantic Foam 为每个 Voronoi 单元赋予对象身份特征和语言对齐语义，通过体渲染积累这些特征来产生 2D 语义输出。关键优势在于：单元级的语义特征场可直接在相邻单元间施加空间正则化约束（相邻单元更可能属于同一对象），避免了 3DGS 点基表示中高斯之间的遮挡和跨视角不一致伪影。同时 Voronoi 结构的常数时间光线遍历保持了高效渲染。在 Mip-NeRF 360、LERF-Mask 和 LLFF 共 16 个场景上，Semantic Foam 在对象级分割上超越 Gaussian Grouping 和 SAGA。

\paragraph{高斯点编码与3D识别} 3D Gaussian Point Encoders（2025.11）~\cite{2511_04797} 提出了一种不同于传统分割的视角：将学习到的 3D 高斯混合作为显式的逐点嵌入用于 3D 识别任务。

\textbf{动机。}传统 3D 识别方法（如 PointNet）需要在原始点云上运行重量级网络；该方法将高斯混合的参数作为紧凑的中间表示，基于自然梯度和 PointNet 蒸馏优化高斯基，构建的编码器比同精度 PointNet 快 2.7 倍。该方法基于自然梯度和 PointNet 蒸馏优化高斯基，构建的编码器比同精度 PointNet 快 2.7 倍、内存减少 46\%、FLOPs 减少 88\%，且轻量到可在纯 CPU 设备上实现高帧率。这一工作将 3DGS 从场景重建扩展到通用 3D 特征学习。

\paragraph{RT-GS2：前馈可泛化语义分割} 上述方法大多需要逐场景优化——每个新场景都要从头训练特征场，限制了实际部署。RT-GS2（2024.05）~\cite{2405_18033} 首次提出可泛化的语义分割方法，从逐场景优化转向前馈推理。

\textbf{动机。}前述方法大多需要逐场景优化，每个新场景都要从头训练特征场，限制了实际部署；RT-GS2 训练一次即可应用于新场景无需重训练，将分割速度从分钟级提升至 27 FPS。以整个 3DGS 表示作为输入，先通过自监督方式提取视角无关的 3D 高斯特征 $F^k = \{f^k_1, \dots, f^k_N\} \in \mathbb{R}^{N \times D}$，再通过 $\alpha$-混合渲染到特定视角得到 2D 特征图，最后设计 View-Dependent/View-Independent（VDVI）特征融合模块——将渲染的 3D 特征（视角无关但缺乏上下文）与 2D 图像特征（视角相关但包含丰富上下文）自适应融合——增强跨视角语义一致性。这种设计使模型在训练后可直接应用于新场景无需重训练，在 Replica 上平均交并比（mean Intersection over Union, mIoU）提升 8.01\% 且达到 27 FPS（901 倍加速）。

\paragraph{SemGS：双高斯前馈语义重建} SemGS（2026.03）~\cite{2603_02548} 从双表示角度设计前馈架构：每个输入图像像素对应两个互补高斯——颜色高斯用于外观建模、语义高斯用于语义推理——共享 3D 位置和不透明度但各自维护独立的协方差和属性，通过共享浅层 CNN 的双分支提取器分别预测，引入相机感知注意力建模视角间几何关系，支持从稀疏视角输入进行前馈语义重建。

\paragraph{SegSplat：语义记忆库前馈分割} SegSplat（2025.11）~\cite{2511_18386} 构建紧凑的语义记忆库（semantic memory bank）：先用 SAM+CLIP 从输入视角提取语义特征并聚类为代表性语义概念，形成码本；然后基于 DepthSplat 前馈生成 3D 高斯基元，每个基元关联一个语义索引指向码本中的条目；在单次前向传播中同时预测几何属性和离散语义索引，渲染时通过码本查找即可产生语义特征图，无需逐场景语义优化即可实现开放集分割。

\paragraph{OVGaussian：跨场景泛化语义分割} OVGaussian（2025.01）~\cite{2501_00326} 则将可泛化性推进到跨场景级别：构建了大规模 3DGS 语义数据集 SegGaussian，提出 Generalizable Semantic Rasterization（GSR）使 3D 高斯能跨场景一致地表示语义信息，以及 Cross-modal Consistency Learning（CCL）框架对齐 3D 和 2D 语义——用 3D 神经网络学习每个高斯的语义属性并渲染为多视角一致的语义图，首次实现跨场景、跨域的开放词汇 3DGS 语义分割。

\paragraph{无监督与弱监督分割} PointGS（2026.05，Point-Gaussian Segmentation）~\cite{2605_11520} 利用 3DGS 作为中间表示桥接离散 3D 点云和连续 2D 图像的域间隙。

\textbf{动机。}3D 点云稀疏且缺乏外观信息，2D 图像密集但缺乏几何一致性，两者之间存在难以跨越的域间隙；3DGS 天然兼具几何连续性和外观信息，可作为桥接中间表示实现语义跨域迁移。其桥接机制分三步：先将稀疏点云重建为密集高斯空间，利用高斯的连续可微性填补空间空隙并建立多视角对应关系；然后通过多视角 SAM 掩码和亲和特征对比学习将语义蒸馏到高斯基元——同一 SAM 掩码内的高斯亲和特征被拉近、不同掩码的被推远；最后通过两步 ICP 配准（先粗配准再精配准）将高斯语义迁移回原始点云，在 ScanNet-V2 和 S3DIS 上分别取得 +0.9\% 和 +2.8\% mIoU 提升。DCSEG（2024.12，Decoupled Clustering-based Segmentation）~\cite{2412_10972} 则从流水线架构角度解决语义分割与 3D 重建的耦合问题。

\textbf{动机。}现有方法将 3D 重建与语义分割耦合在统一框架中，限制了模块替换的灵活性；DCSEG 将流水线完全解耦为类别无关的掩码提议和类别感知的标签分配两个阶段，各模块可独立替换，提供即插即用接口。提出完全解耦的两阶段流水线：第一阶段为类别无关的掩码提议（mask proposal），通过对比学习聚类高斯形成粗 3D 分割，仅依赖 SAM 的类别无关掩码和 3D 空间邻近关系；第二阶段为类别感知的标签分配，将提议的 3D 掩码投影到各视角并与 CLIP 等视觉-语言模型的 2D 开放词汇掩码匹配，赋予语义标签。这种解耦设计提供即插即用接口——3D 重建模块（如 3DGS）和语义特征提取模块（如 SAM+CLIP）均可独立替换，同时在多对象实例分割上无额外开销，克服了现有方法无法同时实现类别感知、3D 信息利用、实例区分和解耦重建的局限性。

\paragraph{OMEGAS：2DGS 目标分割与遮挡补全} 许多分割方法的最终目标是支持场景编辑。OMEGAS（2024.04，Object-level Mesh Extraction via Gaussian Splatting）~\cite{2404_15891} 利用 2DGS 进行 3D 目标分割，生成多视角一致的目标掩码和初步模型，再用扩散先验补充未见的遮挡部分，实现从大场景中精确重建特定对象。

\textbf{动机。}场景级重建中特定对象的几何往往不完整（被遮挡或视角不足），直接分割提取质量有限；OMEGAS 先分割目标再利用扩散先验补全遮挡区域，实现从大场景中精确提取完整对象。

\paragraph{LabelGS：遮挡感知的高效分割} OMEGAS 未显式处理遮挡导致的过拟合。LabelGS（2025.08）~\cite{2508_19699} 引入跨视角一致语义掩码和遮挡分析模型（Occlusion Analysis Model），避免优化中的遮挡过拟合，配合随机区域采样策略将训练效率提升 22 倍。

\textbf{动机。}OMEGAS 未显式处理遮挡导致的过拟合——被遮挡区域的 2D 掩码缺失会使分割模型在遮挡边界处过度拟合可见信号；LabelGS 通过遮挡分析模型显式建模可见性，避免在不可靠区域过度优化。

\paragraph{Gradient-Driven 3D Segmentation：梯度投票分割} LabelGS 依赖语义监督，而 Gradient-Driven 3D Segmentation（2024.09）~\cite{2409_11681} 提出基于投票的梯度方法：利用 2D 掩码过滤梯度作为投票信号实现 3D 分割，并发现推理时梯度也可用于高斯剪枝（最高 21\% 压缩）。

\paragraph{室内资产检测：无人机 3DGS 分割} 面向无人机采集的特殊场景，室内资产检测方法（2026.04）~\cite{2604_05316} 利用 360\textdegree{} 无人机图像重建 3DGS 场景，引入 3D 对象码本联合利用掩码语义和空间信息指导多视角掩码关联，F1 分数比基线提升 65\%。

\subsection{开放词汇与交互式分割}
\label{ss:10_2}

本节讨论3条技术路线：（1）开放词汇分割；（2）交互式与点击式分割；（3）全景分割。

\subsubsection{开放词汇分割}

本路线介绍允许用户用自然语言描述任意对象并在 3D 场景中定位分割该对象的开放词汇分割方法。与固定类别集的传统分割不同，开放词汇分割需要理解任意自然语言描述并将其与 3D 空间中的高斯基元关联。

\paragraph{LangSplat：语言特征场先驱} LangSplat（2023.12）~\cite{2312_16081} 是这一方向的先驱，为每个高斯附加 CLIP 特征并通过 3D 语言高斯光栅化实现开放词汇查询。

\textbf{动机。}LERF 通过在 NeRF 密集体中存储 CLIP 特征建立语言场，但体渲染的计算开销严重制约查询效率；LangSplat 利用 3DGS 显式点云表示的高效渲染优势，首次在 3DGS 上构建端到端的开放词汇查询语言场。

\paragraph{Ilov3Splat：像素到实例级语言理解} LangSplat 的语言特征在像素级操作，缺乏实例感知。Ilov3Splat（2026.05）~\cite{2605_04506} 将像素级语言特征提升到实例级开放词汇理解。

\textbf{动机。}LangSplat 为每个高斯直接存储高维特征向量，内存开销巨大；Ilov3Splat 利用多分辨率哈希嵌入高效编码 CLIP 特征，通过轻量 MLP 按需解码，避免直接存储高维向量。利用多分辨率哈希嵌入（hash grid）高效编码语言对齐的 CLIP 特征——以每个高斯位置为查询点，通过轻量 MLP 从哈希网格中解码语义嵌入，避免了 LangSplat 为每个高斯直接存储高维特征向量的内存开销——通过 SAM 掩码上的对比损失同时训练语言对齐特征和实例判别特征，在推理时通过两阶段 3D 聚类（先 DBSCAN 粗聚类再精细分组）检索相关高斯组。

\paragraph{OpenGaFF：几何条件化连续特征场} Ilov3Splat 仍将语义作为独立特征附加。OpenGaFF（2026.05）~\cite{2605_06088} 提出 Gaussian Feature Field，将语义建模为高斯几何和外观的连续函数，而非 LangSplat 的独立特征附加：通过以高斯的几何结构（位置、协方差）为条件增强语义预测，使空间相邻且结构相似的高斯获得一致的语义，提升相似结构间的空间一致性。该方法还引入数据驱动的自适应结构化码本作为共享语义词元集——不同于 LangSplat 的场景特定自编码器，OpenGaFF 的码本从所有 SAM+CLIP 提取的特征中初始化，覆盖场景中实际存在的语义类别——配合码本引导注意力机制检索语言特征，降低对象内特征方差。

\paragraph{SCOUP：稀疏码本加速} OpenGaFF 引入结构化码本但训练成本仍高。SCOUP（2026.05）~\cite{2605_13600} 从效率角度出发，提出稀疏码本提升策略：将语言表示学习与 3DGS 优化解耦——核心洞察是场景的 CLIP 特征高度冗余，仅占据完整 CLIP 空间的一个小子空间，因此用稀疏编码替代密集特征。在 2D 图像区域学习码本系数后通过加权稀疏聚合提升到 3D 高斯，Top-$K$ 过滤提取每个高斯的主导多视角系数（取各视角中一致度最高的 $K$ 个系数聚合），实现 $400\times$ 训练加速和 $3\times$ 训练内存效率提升。

\paragraph{LangFlash：前馈推理} 前述方法均需逐场景优化。LangFlash（2026.05）~\cite{2605_23287} 将语言特征场的构建从逐场景优化推向前馈推理：基于冻结的 NoPoseNet 从稀疏无位姿图像单次前向预测几何和语义，引入语义分组模块用 Siamese ViT 解码器产生多尺度密集语义图，再通过语言特征聚合模块将 SAM+CLIP 的语言特征分配到各语义组形成语义高斯表示，支持低延迟 3D 重建和语言一致性场景理解。

\paragraph{Lee et al.：3D 体分割范式} 前述方法均在 2D 像素上监督然后蒸馏到 3D。Lee et al.（2024.08）~\cite{2408_07416} 重新定义开放词汇分割为 3D 体分割问题：直接在 3D 光线采样点上用尺度无关的语言嵌入监督训练，避免了 2D 到 3D 蒸馏中的信息损失，再将学到的语言场迁移到 3DGS，首次在不牺牲训练时间或精度的前提下实现实时渲染速度，并提出 3D 查询和评估协议。

\paragraph{OpenSplat3D：SAM 实例+对比学习} OpenSplat3D（2025.06）~\cite{2506_07697} 结合 SAM 实例掩码和对比损失训练实例特征，利用视觉-语言模型的语言嵌入实现文本驱动的实例识别，在 LERF-mask、LERF-OVS 和 ScanNet++ 验证集上展示有效性。

\paragraph{OpenInsGaussian：上下文感知特征聚合} OpenSplat3D 的逐掩码特征缺乏上下文。OpenInsGaussian（2025.10）~\cite{2510_18253} 针对两个核心问题：上下文感知特征提取为每个掩码补充丰富语义上下文（解决预处理中个体掩码语义不足），注意力驱动特征聚合选择性融合多视角特征以缓解对齐误差和不完整性。

\paragraph{COS3D：实例-语言协作场} OpenInsGaussian 的实例和语言特征各自独立。COS3D（2025.10）~\cite{2510_20238} 提出协作场（collaborative field），包含实例场和语言场两个子场：训练时通过实例到语言的特征映射和高效两阶段策略捕获两个场的内在关系；推理时通过自适应语言到实例的提示精炼弥合两个场的特性差异。

\paragraph{ReferSplat：指代 3DGS 新任务} 前述方法做开放词汇实例分割，ReferSplat（2025.08）~\cite{2508_08252} 则定义了 Referring 3DGS（R3DGS，指代 3D 高斯分割）新任务——根据自然语言描述分割 3D 高斯场景中的目标对象，构建了 Ref-LERF 数据集，通过空间感知范式显式建模 3D 高斯与自然语言表达式的关系，在 R3DGS 和 3D 开放词汇分割基准上均达到 SOTA。

\textbf{动机。}开放词汇分割返回所有匹配类别的实例，而用户往往只需要定位特定描述的那个对象（如"被杯子挡住的水果"）；ReferSplat 定义了指代分割新任务，通过空间感知范式显式建模 3D 位置与自然语言表达式的关系。

\paragraph{CaRF：相机条件化对齐} ReferSplat~未显式建模跨视角差异。CaRF（2025.11）~\cite{2511_03992} 引入 Camera-conditioned Alignment Modulation（CAM）将相机几何注入高斯-文本交互，以及 Gaussian-level Cross-view Logit Alignment（GCLA）显式对齐同一高斯在不同视角下的指代表达响应，将跨视角差异转化为可优化目标，在~Ref-LERF~上~mIoU~提升~16.8\%。

\paragraph{TrackRef3D：自动化 track-then-label} CaRF 仍需手动标注。TrackRef3D（2026.05）~\cite{2605_26576} 提出完全自动化的 track-then-label 范式：通过轨迹感知语义共识模块（TSCM）聚合跨视角预测建立规范语义身份，配合混合训练策略联合优化粗类别语义和细粒度指代线索，实现无需手动标注的开放世界指代分割。

\paragraph{Segment then Splat：先分割后重建} 前述方法均遵循"先重建后分割"范式。Segment then Splat（2025.03）~\cite{2503_22204} 颠覆了这一流程——先将高斯分组为不同对象集再重建，使场景天然分割为独立对象，同时支持静态和动态场景。

\textbf{动机。}"先重建后分割"范式中，一个高斯可能同时属于多个对象，导致分割边界模糊；Segment then Splat 将分割提前到重建阶段，每个对象独立重建，从根本上消除高斯-对象错位问题。这一设计消除了几何和语义歧义以及动态场景中的高斯-对象错位问题。

\paragraph{PanoGS：金字塔 tri-plane 全景分割} Segment then Splat 基于简单分组。PanoGS（2025.03）~\cite{2503_18107} 采用金字塔 tri-plane 建模潜在连续参数特征空间，用 3D 特征解码器回归多视角融合的 2D 特征云，提出语言引导的图割将高斯聚合为超基元，再通过 SAM 引导的图聚类实现 3D 一致实例分割。

\textbf{动机。}Segment then Splat 基于简单分组，粒度控制粗糙；PanoGS 通过金字塔 tri-plane 建模连续特征空间，配合语言引导的图割实现更精细的实例分割。

\subsubsection{交互式与点击式分割}

本路线介绍支持用户通过点击、涂画或文本提示选择 3D 对象的交互式分割方法，是实现场景编辑的关键交互方式。这类方法的核心目标是在交互响应速度和分割精度之间取得平衡。

\paragraph{Click-Gaussian：两级粒度交互分割}

在 Click-Gaussian 出现之前，交互式 3DGS 分割方法（如 Gaussian Grouping、SAGA 等）为每个高斯赋予单一粒度的身份编码——这种全局粒度设计无法同时支持"选择整张桌子"和"选择桌上的杯子"两种层级的交互需求。Click-Gaussian（2024.07）~\cite{2407_11793} 的核心创新在于学习两级粒度（粗粒度和细粒度）的可区分特征场，使同一场景能响应用户在不同语义层级上的点击选择。

\textbf{动机。}前述交互式分割方法为每个高斯赋予单一粒度的身份编码，无法同时支持"选择整张桌子"和"选择桌上的杯子"两种层级的交互需求；Click-Gaussian 通过粒度先验设计，使细粒度特征天然继承粗粒度的区分信息。

\textbf{特征设计：粒度先验。}直觉上，粗粒度和细粒度之间存在天然的层级依赖——如果两个对象在粗粒度上已经不同（例如桌子和椅子），那么它们的任何细粒度部分自然也不同。Click-Gaussian 利用这一\emph{粒度先验}（granularity prior）设计特征空间：为每个高斯赋予 $D{=}24$ 维特征向量 $f_i$，将其拆分为粗粒度部分 $f_i^c \in \mathbb{R}^{12}$ 和细粒度补充部分 $\hat{f}_i^c \in \mathbb{R}^{12}$，细粒度特征定义为 $f_i^f = [f_i^c;\, \hat{f}_i^c]$（级联而非独立）。这样，细粒度特征天然继承了粗粒度的区分信息，避免了两个层级独立学习时的冲突。两级特征通过同一个前向过程在单次光栅化中渲染出来。

\textbf{训练策略：对比学习 + 全局特征引导。}训练分为三个阶段。首先，利用 SAM 的自动掩码生成模块为所有训练视角产生 2D 分段，按分段面积排序——面积最大的掩码作为粗粒度标签、最小的作为细粒度标签。然后，在两级掩码上分别施加余弦相似度对比损失：同掩码内的像素特征最大化相似度（$L_{\text{pos}}^{\text{cont}}$），不同掩码的像素特征约束低于余弦边距（细粒度 $\delta^f{=}0.75$，粗粒度 $\delta^c{=}0.5$）。关键设计是在细粒度负对比中，对粗粒度分量施加 stop-gradient（$F^f = \text{sg}(F^c) \oplus \hat{F}^c$），确保细粒度学习不干扰已稳定的粗粒度特征。

然而，跨视角的 SAM 掩码存在天然的不一致性——同一对象在不同视角可能被 SAM 分成不同数量的段落，直接用这些噪声标签逐视角训练会导致特征场不稳定。为此，Click-Gaussian 提出 Global Feature-guided Learning（GFL）机制：从第 2000 次迭代开始，每隔一定步数，对所有训练视角渲染 2D 特征图并对每个掩码做平均池化，得到该掩码的全局特征表示；然后用 HDBSCAN 聚类算法（粗粒度 eps $= 10^{-2}$，细粒度 eps $= 10^{-3}$）将这些特征聚为 $C^l$ 个全局候选簇。训练时，每个高斯的特征被引导向其最相似的全局簇靠近（余弦相似度阈值 $\theta^g{=}0.9$）、远离其他簇。由于全局簇是从所有视角的渲染特征聚合而来，它们代表了场景中最具代表性的特征模式，有效平滑了 SAM 掩码的跨视角噪声。

整个训练过程约 3000 次迭代，在 NVIDIA RTX A5000 上耗时约 13 分钟。训练完成后，用户点击场景中的任意位置即可在约 10ms 内检索到对应的粗/细粒度全局簇并提取相关高斯——比 SAGA 快约 15 倍，比 Feature3DGS 快约 130 倍。在 LERF-Mask 和 SPIn-NeRF 数据集上，Click-Gaussian 在 mIoU 指标上均超越所有基线，尤其在细粒度分割上优势显著。

\paragraph{GaussianCut：图割式免训练交互分割} 与 Click-Gaussian 的两级特征场设计不同，GaussianCut（2024.11）~\cite{2411_07555} 走了一条无需特征训练的路线：将场景建模为图结构，利用图割算法最小化能量函数来划分前景和背景高斯，支持点击、涂画和文本等多种交互方式。

\textbf{动机。}Click-Gaussian 等方法需要训练特征场，训练完成后虽然可反复使用但缺乏灵活性（粒度固定）；GaussianCut 不训练任何特征，每次交互都基于当前需求重新运行图割优化，代价是每次交互需重新计算。

\paragraph{ArtisanGS：交互式工具套件} ArtisanGS（2026.02）~\cite{2602_10173} 从工程实用性出发，提供围绕高斯选择和分割的交互式工具套件：快速 AI 驱动的 2D 选择到 3DGS 选择的传播（允许用户干预修正错误），辅以灵活的手动选择和分割工具，可实现任何非结构化 3DGS 场景的任意二元分割。与 Click-Gaussian 的自动两级粒度不同，ArtisanGS 的粒度完全由用户在交互过程中动态决定。

\paragraph{B$^3$-Seg：贝叶斯最优视角采样} B$^3$-Seg（2026.02）~\cite{2602_17134} 从贝叶斯视角重新定义 3DGS 分割：将分割建模为序列 Beta-Bernoulli 贝叶斯更新，通过解析式 Expected Information Gain（EIG，期望信息增益）主动选择下一个最优视角，保证自适应单调性和次模性，产生 $(1-1/e)$ 近似的最优视角采样策略。

\textbf{动机。}现有分割方法假设视角已给定，但并非所有视角都对分割同等有用——某些视角包含更多辨别性信息；B$^3$-Seg 主动选择信息增益最大的视角进行标注，以最少的视角获得最准确的分割。与 Click-Gaussian 需要预先训练特征场不同，B$^3$-Seg 完全不需要相机参数和训练过程，但以多视角采集为代价换取了无需训练的优势。

\paragraph{Split\&Splat：独立重建式实例分割} Split\&Splat（2026.02）~\cite{2602_03809} 从另一个角度出发：先利用深度信息跨视角传播实例掩码生成视角一致的 2D 掩码，为每个对象独立重建再合并回场景，天然支持全景分割、对象检索和 3D 编辑等下游任务，在 ScanNetv2 上取得 SOTA。

\textbf{动机。}全局重建后分割会导致跨对象的混合高斯，分割边界模糊；Split\&Splat 为每个对象独立重建，使每个高斯天然归属于单一对象，从根本上消除边界歧义。

\paragraph{PointGauss：点云引导多对象分割} PointGauss（2025.08）~\cite{2508_00259} 提出点云引导的多对象分割框架，通过基于点云分割的高斯基元解码器在 1 分钟内生成 3D 实例掩码。

\textbf{动机。}纯 3DGS 分割方法仅依赖 2D 视觉特征，缺乏三维空间结构的直接引导；PointGauss 利用预训练点云分割模型的 3D 空间理解能力引导高斯基元的分组，结合 2D 视觉和 3D 几何两种信息源。并贡献了 DesktopObjects-360 多对象分割数据集（27000+ 2D 掩码、全局一致标注、360\textdegree{} 覆盖）。

\paragraph{SAGOnline：零样本在线分割} SAGOnline（2025.08，Segment-Any-Gaussian Online）~\cite{2508_08219} 实现零样本在线分割：集成视频基础模型（SAM 2）生成时序一致的 2D 掩码，通过 Rasterization-aware Geometric Consensus 机制将 2D 预测确定性地映射为 3D 高斯标签，达到 27ms/帧的速度和 92.7\%/95.2\% mIoU。

\textbf{动机。}离线分割方法需要所有视角同时可用；SAGOnline 利用 SAM 2 的时序一致性实现在线零样本分割，无需任何训练即可将新帧的 2D 分割确定性映射为 3D 高斯标签。

\paragraph{多视角视频分割数据集} MUVOD（2025.07，Multi-View Object Detection）~\cite{2507_07519} 是首个多视角视频对象分割数据集和 3D 分割基准，填补了 NeRF/3DGS 场景理解领域缺乏统一评估标准的空白。

\textbf{动机。}3DGS 分割方法的评估缺乏统一的多视角视频基准，不同方法在不同数据集上比较，难以公平评估；MUVOD 提供了标准化的多视角视频标注和评估协议。

\paragraph{Seg-Wild：互联网无约束照片分割} Seg-Wild（2025.07）~\cite{2507_07395} 处理互联网无约束照片集合的分割挑战——这些图像存在不一致的光照和瞬态遮挡。

\textbf{动机。}室内/受控环境下的分割方法假设拍摄条件一致，但互联网照片存在光照剧变、瞬态物体和不同外观模式；Seg-Wild 引入 Spiky 3D Gaussian Cutter（SGC）平滑异常高斯，并设计了评估野外场景分割质量的基准。

\paragraph{CrashSplat：车辆损伤 3D 分割} CrashSplat（2025.09）~\cite{2509_23947} 专注于车辆损伤的 3D 分割，提出基于 Z-buffer 和深度-不透明度正态分布模型的无学习单视角分割方法。

\textbf{动机。}多视角一致性方法假设损伤在多个视角可见，但实际车辆损伤（如划痕、小凹痕）往往仅在特定角度可见；CrashSplat 的无学习方法仅需单视角即可完成损伤分割，适用于仅在一个视角可见的损伤类型（如划痕、小凹痕），在多视角一致性方法不适用的场景中表现突出。

\paragraph{SAMa：材质选择} SAMa（2024.11，Segment Anything Material）~\cite{2411_19322} 将 SAM2 的视频先验扩展到材质领域：构建材质中心视频数据集训练模型，通过深度投影将 2D 预测提升到 3D 点云，实现无需逐资产优化的材质选择（秒级完成），可替换 text-to-3D 输出中的漫反射纹理为 PBR（Physically Based Rendering，基于物理的渲染）材质。

\textbf{动机。}SAM2 在物体分割上表现优异，但无法区分同一物体的不同材质（如木质桌面和金属桌腿）；SAMa 通过在材质中心视频数据集上微调，使 SAM2 的视频先验适用于材质级分割。

\subsubsection{全景分割}

本路线介绍同时完成语义分割和实例分割的 3DGS 全景分割方法。全景分割（对"事物"类别的像素级分类和对"物品"类别的实例级分割）是最全面的场景分割任务。3DGS 全景分割的挑战在于：3DGS 缺乏参数化结构（不像 NeRF 的 MLP 天然具有平滑性），因此对噪声 2D 掩码更加敏感。

\paragraph{3DGS 全景分割基础} PLLS（2024.10）~\cite{2410_17505} 首次将全景分割引入 3DGS 框架。

\textbf{动机。}在 NeRF 时代，MLP 的参数化结构天然提供平滑映射可隐式平滑 2D 掩码噪声；但 3DGS 每个高斯独立编码语义属性，噪声被忠实地记忆而非平滑，导致渲染掩码不一致。PLGS 的核心动机是解决 3DGS 不满足平滑性假设这一根本问题。在 NeRF 时代，Panoptic Lifting 和 Contrastive Lift 利用 MLP 的参数化结构天然提供特征到掩码的平滑映射，即使 2D 掩码存在跨视角噪声，MLP 也能隐式地将其平滑掉。然而 3DGS 没有这样的参数化结构。每个高斯独立编码语义属性，不存在全局平滑约束，因此即使视角仅有微小变化，2D 分割模型（如 Mask2Former）产生的噪声也会被高斯忠实地记忆而非平滑，导致渲染掩码出现明显的不一致性和伪影。直接沿用 NeRF 时代的噪声处理策略（测试时增强、分割一致性损失等）也无法获得满意结果。

为此，PLGS 从结构化和降噪两个层面系统性地弥补 3DGS 的平滑性缺陷。在结构化层面，PLGS 基于 Scaffold-GS 构建全景感知的高斯模型：通过体素化操作 $\mathbf{V} = \{[\mathbf{P}/\epsilon]\} \cdot \epsilon$（其中 $\epsilon$ 为体素尺寸）将点云离散化为锚点集，每个锚点携带可学习特征 $\hat{f}_v$，再由轻量 MLP 解码器（如语义解码器 $\mathbf{F}_{\text{sem}}$）为周围的 $k$ 个高斯预测语义概率 $\hat{y}$ 和不透明度等属性，从而引入锚点级别的局部平滑。在降噪层面，语义场采用三阶段策略：(1) 一致性增强掩码——对于参考视角的每个像素 $\hat{\mathbf{u}}$，通过深度引导的跨视角变形（利用相机内参 $\mathbf{K}_{\text{cam}}$ 和外参 $\mathbf{T}_w^c$）收集邻域 $N_w$ 个视角的对应像素集合 $\Psi(\hat{\mathbf{u}})$，执行投票过滤 $q(\hat{\mathbf{u}}) = \prod_{\mathbf{u} \in \Psi(\hat{\mathbf{u}})} \mathbbm{1}\{s_{\hat{\mathbf{u}}} = s_{\mathbf{u}}\}$，仅保留所有邻域视角语义标签一致的像素；(2) 可靠语义锚点——将过滤后的语义标签投影到 3D 并通过聚类生成全局一致的语义点云替代 SfM 稀疏点作为初始化，同时将语义点云反向投影回各视角生成 2D 一致语义掩码用于监督训练中新生长的高斯；(3) 自训练伪标签——将模型渲染掩码中语义标签一致的区段与原始噪声掩码的区段合并作为伪标签，在训练后期进一步增强鲁棒性。实例场方面，将 2D 实例掩码投影到 3D 空间并与定向包围盒（oriented bounding box）匹配，结合线性分配损失和一致实例掩码监督生成跨视角一致的实例掩码，二者互补（消融实验表明结合使用比单独使用分别提升 3.7 和 1.5 PQ$^{\text{SQ}}$）。在 ScanNet 上，PLGS 以远超 NeRF 方法的训练和渲染速度达到 SOTA，相比 Gaussian-Grouping 在 mIoU 上提升 8.1、PQ$^{\text{SQ}}$ 上提升 11.1。

\paragraph{端到端全景重建} PLGS 虽然首次将全景分割引入 3DGS，但其流水线仍是多阶段的——先训练 3DGS、再训练语义场和实例场。

\textbf{动机。}多阶段流水线中语义场和实例场的训练无法受益于彼此的监督信号，且训练时间线性叠加；PanopticSplatting 将语义和实例预测统一到端到端框架中，使两者相互促进。PanopticSplatting（2025.03）~\cite{2503_18073} 进一步提出端到端全景重建系统：为每个高斯附加可学习的语义特征 $f_\text{sem}$ 和实例特征 $f_\text{ins}$，这些特征与视角无关（因为语义和实例在不同视角下应保持一致），通过 $\alpha$-混合光栅化直接渲染 2D 语义图和实例图。核心架构引入 query-guided Gaussian Segmentation 和局部交叉注意力——learnable queries 通过交叉注意力与高斯特征交互，端到端预测实例掩码和语义标签，无需后处理关联。局部交叉注意力限制在视锥体内（只关注当前视角附近的高斯），有效减少训练内存使模型可处理包含更多高斯和对象的大场景。同时设计标签混合（label blending）策略——将多视角预测的标签以加权方式混合为一致的全局标签——促进一致的 3D 分割并减少噪声浮游体，以及标签扭曲（label warping）利用深度信息将相邻视角的标签投影到当前视角作为额外监督——增强多视角一致性和分割精度。在 ScanNet-V2 和 ScanNet++ 上超越 NeRF-based 和 Gaussian-based 方法。

\paragraph{OnlinePG：在线全景建图} 上述全景分割方法均为离线处理——需要所有视角的图像一次性输入。OnlinePG（2026.03，Online Panoptic Gaussian）~\cite{2603_18510} 实现在线全景建图：采用滑动窗口的局部到全局范式，增量处理 RGB-D 流。

\textbf{动机。}前述全景分割方法均为离线处理——需要所有视角的图像一次性输入；OnlinePG 通过滑动窗口增量处理 RGB-D 流，支持实时在线场景理解。在局部阶段，对当前帧的 2D 分割结果通过段聚类算法融合为一致的 3D 实例——协同利用几何邻近性（3D 空间距离）和语义相似性（特征距离）将不一致的 2D 片段合并为完整实例。在全局阶段，构建显式空间属性网格存储场景语义信息，通过双向二部图高斯实例匹配将新检测到的局部实例与全局地图中已有实例进行最优匹配更新，保持实时效率。

\paragraph{MaskField：掩码特征场分解} MaskField（2024.07）~\cite{2407_01220} 从新的视角解决开放词汇分割——不直接蒸馏高维 CLIP 特征（LangSplat 等方法的瓶颈所在），而是分解为掩码特征场和查询。

\textbf{动机。}LangSplat 直接蒸馏 512 维 CLIP 特征，内存开销大且容易过拟合；MaskField 将高维特征蒸馏分解为低维掩码嵌入加查询的方式，利用 SAM 的对象形状先验避免复杂正则化，特别适合 3DGS 等高斯基元数量巨大的场景。掩码特征场为每个 3D 点预测一个低维掩码嵌入，自然引入 SAM 分割的对象形状作为边界先验，避免复杂正则化；查询时通过计算掩码嵌入与文本/点提示的相似度定位目标对象。这种分解绕过了训练中对密集高维 CLIP 特征的直接处理，特别适合 3DGS 等显式场景表示中高斯基元数量巨大的场景。

\paragraph{Fast-SegSim：2DGS 实时开放词汇分割} Fast-SegSim（2026.04）~\cite{2604_10951} 基于 2DGS（而非 3DGS）实现实时开放词汇分割重建：引入 Precise Tile Intersection 减少光栅化冗余——精确计算 2D 高斯 surfel 与像素块的交集面积，避免不必要的特征累积——以及 Top-K Hard Selection 策略利用 2D 高斯的几何稀疏性，仅保留对每个像素贡献最大的 $K$ 个 surfel 的语义特征，实现超过 40 FPS 的渲染速度。

\textbf{动机。}3DGS 的体积渲染在语义特征渲染时效率低下，因为体积高斯在深度方向有大量重叠；2DGS 的 surfel 天然贴合表面，几何稀疏性使得仅保留 top-K 贡献者即可实现高保真语义渲染。生成的语义标签还可用于微调下游感知任务（在目标导航中成功翻倍导航成功率）。

\paragraph{PanSt3R：联合预测几何与全景分割} PanSt3R（2025.06，Panoptic Stereo Transformer）~\cite{2506_21348} 则走得更远——联合预测 3D 几何和多视角全景分割，完全消除逐场景测试时优化需求。

\textbf{动机。}所有前述方法均需要逐场景优化或训练；PanSt3R 基于 MUSt3R 的 Transformer 架构，在单次前向传播中同时输出 3D 点和全景分割结果，实现无需优化的全景重建。基于 MUSt3R 的 Transformer 架构增强语义感知和多视角全景分割能力，为每个像素输出 3D 点、类别和实例 ID，通过 Mask2Former 风格的 mask 预测和全局合并步骤生成一致的 3D 全景分割。

\paragraph{Unified-Lift：端到端对象感知提升} Unified-Lift（2025.03）~\cite{2503_14029} 设计了端到端的对象感知提升方法：为每个高斯附加可学习特征，引入可学习的对象级码本显式建模场景中的各个对象，通过关联学习模块和噪声标签过滤模块实现鲁棒的码本学习，在 LERF-Masked、Replica 和 Messy Rooms 三个基准上均超越现有方法。

\paragraph{Trace3D：高斯实例追踪} Trace3D（2025.08）~\cite{2508_03227} 引入 Gaussian Instance Tracing（GIT，高斯实例追踪），为高斯增加跨视角实例权重矩阵，利用高斯在 3D 中的固有一致性识别和纠正 2D 分割不一致，并提出 GIT 引导的自适应密度控制来分裂和修剪歧义高斯。

\textbf{动机。}2D 分割模型在不同视角产生不一致的掩码，直接反投影会导致同一 3D 高斯获得矛盾的标签；Trace3D 利用高斯在 3D 空间中的固有一致性作为仲裁器，识别并纠正 2D 分割的跨视角不一致。

\paragraph{Binary-Gaussian：二进制编码全景分割} Binary-Gaussian（2025.12）~\cite{2512_00944} 提出粗到细的二进制编码方案。

\textbf{动机。}传统方法为每个高斯存储高维语义特征向量，内存开销随类别数线性增长；Binary-Gaussian 将类别特征压缩为单个整数，通过二进制到十进制映射大幅降低内存，同时渐进训练策略分解全景分割为独立子任务减少类别冲突。将每个高斯的类别特征通过二进制到十进制映射压缩为单个整数，大幅降低内存开销，并通过渐进训练策略分解全景分割为一系列独立子任务减少类别冲突，同时微调不透明度解决光度渲染与语义分割的不兼容。

\paragraph{CoRe-GS：机器人粗到细分割} CoRe-GS（2025.09，Coarse-to-fine Region Gaussian Segmentation）~\cite{2509_04859} 面向机器人应用提出粗到细的高斯分割。

\textbf{动机。}全场景分割对机器人应用来说计算浪费——机器人通常只关心操作目标附近的区域；CoRe-GS 仅优化与感兴趣区域（POI）相关的高斯，大幅减少训练时间，并通过基于颜色的过滤机制去除分割产生的浮游体伪影。先生成分割就绪的 GS 表示，再仅优化与感兴趣区域（POI）相关的高斯，大幅减少训练时间，引入基于颜色的过滤机制去除分割引起的浮游体。

\paragraph{LISA-3D：语言-图像分割到 3D} LISA-3D（2025.12）~\cite{2512_01008} 将语言-图像分割提升到 3D：用几何感知 LoRA（Low-Rank Adaptation，低秩适配）层改装 LISA 模型，通过可微重投影损失强制跨视角一致性（仅需 11.6M 可训练参数），支持零样本部署。

\textbf{动机。}LISA 是强大的 2D 语言-图像分割模型，但仅在 2D 空间操作；LISA-3D 通过几何感知 LoRA 层和 3D 重投影损失将 LISA 的分割能力扩展到三维，以极少的可训练参数实现跨视角一致分割。

\paragraph{G2P：高斯属性迁移到点云} G2P（2026.01，Gaussian-to-Point）~\cite{2601_03510} 从相反方向出发——将 3DGS 的外观感知属性（不透明度解决几何歧义、尺度实现精确边界定位）迁移到点云以增强分割判别性，无需 2D 或语言监督。

\textbf{动机。}大多数方法从 2D/语言蒸馏到 3DGS；G2P 反向操作，利用 3DGS 固有的外观感知属性（不透明度、尺度）增强原始点云的分割判别性，无需额外监督信号。

\paragraph{Easy3D：高效 3D 交互分割} Easy3D（2025.04）~\cite{2504_11024} 提出高效的 3D 交互分割方法，将体素稀疏编码器与轻量 Transformer 解码器结合实现隐式点击融合，在包括 3DGS 重建场景在内的多种几何分布上泛化良好。

\textbf{动机。}交互式分割方法通常针对特定几何分布训练，跨域泛化能力有限；Easy3D 通过体素稀疏编码器提取几何无关的特征表示，在包括 3DGS 在内的多种几何分布上保持良好泛化。

\paragraph{GaussianBeV：BEV 感知框架} GaussianBeV（2024.07）~\cite{2407_14108} 将 3DGS 引入 BEV（Bird's Eye View，鸟瞰图）感知框架：用 3D 高斯精细表示场景并 splatting 到 BEV 特征图，在线集成到单阶段模型中，在 nuScenes BEV 语义分割上取得新的 SOTA。

\textbf{动机。}传统 BEV 感知方法从多视角图像直接推断 BEV 特征，缺乏显式的三维几何表示；GaussianBeV 用 3D 高斯作为中间表示提供精细的三维几何先验，通过 splatting 将其投影到 BEV 特征图增强感知精度。

\paragraph{小麦穗实例分割：农业应用} 小麦穗实例分割（2026.03）~\cite{2603_14309} 展示了 3DGS 在农业领域的应用。

\textbf{动机。}农业场景中密集分布的小麦穗在 2D 图像中严重遮挡，传统 2D 实例分割方法难以准确计数；该方法利用 3DGS 提供多视角一致性，通过 3D-to-2D 多视角投影和 Grounded SAM（一种将检测与分割结合的 SAM 变体）实现零样本 2D 分割，再融合为 3D 实例方案训练有监督分割网络。通过 3D-to-2D 多视角投影和 Grounded SAM 实现零样本 2D 分割，再融合为 3D 实例方案训练有监督分割网络。

\paragraph{农业车辆分割：合成数据训练} 农业车辆分割（2025.06）~\cite{2506_05009} 利用 3DGS 和 GOF 生成农业车辆的 3D 资产放入模拟环境，PTv3 模型仅用合成数据训练即达到 91.35\% mIoU。

\paragraph{Robust Prior-Guided Segmentation：先验引导标签重分配} Robust Prior-Guided Segmentation（2026.05）~\cite{2605_16065} 利用 SAM-HQ 生成精确 2D 掩码，引入先验引导的标签重分配方法通过学习先验强制多视角一致性。

\textbf{动机。}标准 SAM 生成的掩码在不同视角间可能不一致，直接用于 3D 分割会产生噪声累积；该方法利用 SAM-HQ（SAM 的高质量变体）生成更精确的 2D 掩码，并通过先验引导的标签重分配强制跨视角一致性。

\begin{takeaway}
3D 分割与语义理解通过将2D基础模型的语义能力蒸馏到3DGS中实现3D场景理解。技术路线从SAM掩码反投影到学习特征场（Feature 3DGS、LangSplat），再到农业和自动驾驶等垂直领域的应用。关键挑战是跨视角语义一致性和实例级精细分割。
\end{takeaway}

\paragraph{3DGS 分类去歧：属性辅助分类} 3DGS 分类去歧（2025.03）~\cite{2503_08352} 发现 GS 点云的尺度、旋转和不透明度属性有助于区分线状与平面物体以及透明/反射物体。

\textbf{动机。}传统点云分类方法仅利用三维坐标，无法区分形状相似但材质不同的物体（如透明玻璃杯和不透明陶瓷杯）；3DGS 的额外属性（不透明度、尺度）提供了区分这类物体的判别性线索。构建了首个真实世界 GS 点云分类数据集（20 类、每类 200 对象）。

%=====================================================================

\subsection{近期方法脉络：开放词汇分割、4D高斯分割、Open-vocabulary segmentation of 3D Gaussians等}

这一组近期工作延续 3DGS、4D Gaussian Splatting、Language-embedded 3DGS，主要把改动集中在应用、增量改进、扩展、加速等环节。下面只保留相对既有路线的变化，便于和前文方法直接对照。

\paragraph{GeoSAM-3D}
GeoSAM-3D（2026.05）~\cite{2606_00447} 延续 3DGS 在“开放词汇分割”上的基本设定，主要改变是：GeoSAM-3D：冻结图像/视频基础模型，基于单目3DGS重建和高斯质心图的测地热核传播提示。3DGS 的几何-渲染统一表示保持为底层地图。
其设计动机是：提出轻量开放词汇3D分割设置：用户在一帧点击或命名物体，系统在Gaussian scene graph上按测地热核传播提示，减少相邻但不连通物体之间的泄漏。
具体而言，方法先从单目视频重建3DGS场景，再以高斯质心构图，把2D基础模型特征和提示映射到图节点，通过heat-kernel geodesic距离传播mask。
适用边界是：依赖单目3DGS重建质量和2D基础模型；开源/仓库状态描述性质较强，系统成熟度仍需验证。

\paragraph{Intrinsic 4D Gaussian Segmentation from Scene Cues}
Intrinsic 4D Gaussian Segmentation from Scene Cues（2026.06）~\cite{2606_18623} 延续 4D Gaussian Splatting 在“4D高斯分割”上的基本设定，主要改变是：Intrinsic 4D Gaussian Segmentation：利用场景内部线索进行4D高斯分割。3DGS 的几何-渲染统一表示保持为底层地图。
其设计动机是：提出面向动态4DGS的内生分割思路，把运动、外观和结构线索用于高斯primitive分组。
具体而言，方法避免对每帧每视角调用SAM等2D基础模型，而是在重建出的4D Gaussian representation中进行一致分割。
适用边界是：分割质量依赖4DGS重建质量；相似运动或遮挡对象仍难分离。

\paragraph{Open-Vocabulary and Referring Segmentation for 3D Gaussians Using 2D Detectors}
Open-Vocabulary and Referring Segmentation for 3D Gaussians Using 2D Detectors（2026.06）~\cite{2606_30638} 延续 Language-embedded 3DGS 在“Open-vocabulary segmentation of 3D Gaussians”上的基本设定，主要改变是：GaussDet：用 2D detectors 替代直接 CLIP 蒸馏和固定实例数分组。3DGS 的几何-渲染统一表示保持为底层地图。
其设计动机是：解决 language-driven 3DGS 方法中实例分组需要预设数量或 bottom-up grouping 噪声大、CLIP 语义受限的问题。
具体而言，从 2D detectors 获取更强的目标级语义和区域线索，再投射或融合到 3D Gaussian 表示，实现开放词表查询和指代表达分割。
适用边界是：依赖 2D 检测器的开放词表能力和多视图一致性；遮挡和小物体可能仍有噪声。

\paragraph{DLGStream}
DLGStream（2026.06）~\cite{2606_28840} 延续 Language-embedded 3DGS 在“Language-embedded dynamic Gaussian video”上的基本设定，主要改变是：面向流式 FVV 的低帧尺寸高 FPS 语言 Gaussian 表示。3DGS 的几何-渲染统一表示保持为底层地图。
其设计动机是：针对现有 language Gaussian 表示难以满足 FVV 低帧尺寸和高帧率要求的问题，提出可流式动态语言嵌入表示。
具体而言，在 3DGS/4DGS 自由视点视频表示中蒸馏或压缩 CLIP 等语言特征，使动态场景能够进行开放词表查询，同时满足 streaming 约束。
适用边界是：语言特征压缩可能牺牲细粒度语义；动态场景中的时序一致查询和边界精度仍需关注。

\paragraph{Online Segment 3D Gaussians via Launching Virtual Drones}
Online Segment 3D Gaussians via Launching Virtual Drones（2026.07）~\cite{2607_01628} 延续 3DGS segmentation 在“Online interactive segmentation of Gaussian scenes”上的基本设定，主要改变是：Virtual drones 生成/采样多视角线索，免 per-scene setup 的快速 segmentation。3DGS 的几何-渲染统一表示保持为底层地图。
其设计动机是：解决现有方法在 raw 3DGS scene 上开始交互前需要几十秒到数分钟 setup 的瓶颈。
具体而言，通过 launching virtual drones 在场景中主动获取或渲染有用视角线索，用于在线定位和分割目标 Gaussians，而不进行昂贵的预处理蒸馏。
适用边界是：虚拟视角策略可能受遮挡、场景尺度和用户提示类型影响；极细粒度对象边界仍需验证。

\paragraph{Consistent Scene Understanding in 3D Gaussian Splatting via Multi-Cue Mask Refinement}
Consistent Scene Understanding in 3D Gaussian Splatting via Multi-Cue Mask Refinement（2026.07）~\cite{2607_01708} 延续 Feature-field 3DGS 在“Multi-cue mask refinement for Gaussian feature fields”上的基本设定，主要改变是：多线索 mask refinement，引导 3DGS feature field 优化。3DGS 的几何-渲染统一表示保持为底层地图。
其设计动机是：解决 2D foundation segmentation 直接用于 3DGS 时跨视图碎片化和预测不一致的问题。
具体而言，流程包括 Multi-Cue Extraction、Multi-Cue Mask Refinement 和用一致 2D instance masks 优化 3D Gaussian feature fields，使对象级交互和高保真表示更可靠。
适用边界是：多线索质量依赖 2D 模型和几何估计；复杂遮挡和高度相似实例仍可能混淆。
\section{开放词汇理解与特征蒸馏}
\label{sec:10_3}

开放词汇场景理解和特征蒸馏是 3DGS 场景理解的核心能力。前者旨在突破预定义类别的限制，使系统能根据任意自然语言描述在三维空间中定位和识别物体；后者关注如何将二维基础模型的语义和几何特征有效迁移到三维高斯空间。两者相辅相成：开放词汇能力赋予场景理解灵活性，特征蒸馏则为这种灵活性提供了技术基础。本节从开放词汇识别与查询、特征场蒸馏等角度展开讨论。

\subsection{开放词汇识别与特征场建模}
\label{ss:10_3}

本节讨论2条技术路线：（1）开放词汇识别与查询；（2）特征场蒸馏。

%=====================================================================
\subsubsection{开放词汇识别与查询}

本路线围绕开放词汇三维场景理解的技术演进展开，从早期的 CLIP 特征蒸馏、语言嵌入三维高斯方法，到面向开放词汇的分割、检测和推理系统，按技术路线分类讨论。

%---------------------------------------------------------------------
\paragraph{CLIP 特征蒸馏的基础范式}

将 CLIP 等视觉-语言模型的二维语义特征提升到三维空间，是实现开放词汇理解的基本策略。Feature 3DGS（2023.12）~\cite{2312_03203} 奠定了这一方向的基础。

\textbf{动机。}SAM、CLIP 等二维基础模型能提取强大的语义特征，但仅存在于二维像素空间中，缺乏跨视角一致性；Feature 3DGS 通过将这些二维特征蒸馏到三维高斯基元中，使每个高斯拥有与视角无关的语义属性。本方法修改了 3DGS 的高斯基元属性——为每个高斯基元附加一个任意维度的语义特征向量 $f \in \mathbb{R}^N$，取代了仅渲染颜色的传统设定——而 3DGS 的几何表示（位置、协方差、不透明度）保持不变。标准 3DGS 中，每个高斯基元的可优化属性仅包括位置 $\mu \in \mathbb{R}^3$、旋转四元数 $q \in \mathbb{R}^4$、缩放因子 $s \in \mathbb{R}^3$、不透明度 $\alpha \in \mathbb{R}$ 以及球谐系数（最高 3 阶，用于建模视角相关颜色 $c \in \mathbb{R}^3$）；Feature 3DGS 在此基础上增加语义特征 $f \in \mathbb{R}^N$，其中 $N$ 可以是任意维度，从而将属性集扩展为 $\Theta_i = \{\mu_i, q_i, s_i, \alpha_i, c_i, f_i\}$。

特征渲染的核心思想是复用 3DGS 原有的 $\alpha$-混合 机制：正如 RGB 图像通过逐高斯加权求和得到像素颜色，语义特征图也用完全相同的权重渲染。直觉上，如果一组高斯基元在 RGB 渲染中共同构成了某个物体的外观，那么用同一组权重渲染语义特征，就能保证语义信息与外观信息在空间上对齐。具体而言，像素 $p$ 处的渲染颜色和语义特征分别由以下公式计算：

\begin{equation}
\hat{C}(p) = \sum_{i=1}^{N} c_i \, \alpha_i \, T_i, \qquad
\hat{F}(p) = \sum_{i=1}^{N} f_i \, \alpha_i \, T_i,
\end{equation}

其中 $T_i = \prod_{j=1}^{i-1}(1 - \alpha_j)$ 是透射率，$N$ 是沿射线按深度排序后与该像素重叠的高斯基元集合。两个公式共享完全相同的 $\alpha$-混合 权重 $\alpha_i T_i$，区别仅在于前者混合 RGB 颜色 $c_i$，后者混合语义特征 $f_i$——这意味着语义特征与 RGB 图像具有完全一致的像素级空间对齐关系。

这一设计源于一个核心动机：SAM、CLIP-LSeg 等二维基础模型能够从单张图像中提取强大的语义特征（例如 LSeg 的特征维度为 $M=512$，SAM 的特征维度为 $M=256$），但这些特征仅存在于二维像素空间中，缺乏跨视角的一致性——同一张椅子在不同视角下可能得到不一致的分割结果。Feature 3DGS 通过将这些二维特征蒸馏到三维高斯基元的语义特征向量中，使每个高斯拥有与视角无关的语义属性，从而在新视角合成时也能产生语义一致的特征图。相比 NeRF 方案中特征分支与辐射场分支共享网络层导致互相干扰的问题（2022）~\cite{kobayashi2022decomposing}，Feature 3DGS 的显式表示天然地将外观与语义解耦，允许等权重的联合优化（损失函数中 RGB 损失权重 $\lambda=1.0$，特征损失权重 $\beta=0.2$）而不互相影响。

然而，高维特征渲染面临一个效率瓶颈：直接渲染 512 维特征图极其耗时。Feature 3DGS 引入了一个轻量级的加速模块——在光栅化输出后接一个 $1 \times 1$ 卷积解码器，将低维特征 $f \in \mathbb{R}^N$（$N \ll M$）上采样到目标维度 $M$，从而在保持下游任务性能的前提下大幅提升训练效率。在实现层面，RGB 图像与特征图共享同一个基于 16$\times$16 分块的 tile-based 光栅化流程，两者被渲染到相同分辨率但在不同通道维度上，确保语义特征的像素级精度与 RGB 图像完全一致。Feature 3DGS 支持三种交互模式：基于 SAM 的点/框提示分割（promptable segmentation）、无提示零样本分割（promptless zero-shot segmentation）、以及基于 CLIP-LSeg 的文本驱动语义分割，实现了二维基础模型的全部功能向三维空间的迁移。该方法揭示了将高维特征直接附在高斯基元上时面临的分辨率不匹配和通道一致性等关键挑战，并提出了针对性的架构和训练改进。

Occam's LGS（2024.12，Occam's Language Gaussian Splatting）~\cite{2412_01807} 对此范式进行了极简化重新审视。

\textbf{动机。}LangSplat 引入场景专属自编码器压缩高维特征，但增加了训练复杂度；Occam's LGS 指出基于概率建模的加权多视角聚合即可达到最优效果，无需任何压缩，实现两个数量级的训练加速。该工作指出，现有的语言高斯方法过度依赖复杂的多视角特征聚合流程，而基于概率建模的加权多视角聚合策略即可在无需任何压缩的前提下达到最优效果，实现两个数量级的训练加速。

\paragraph{FastLGS：语义特征网格映射} FastLGS（2024.06，Fast Language Gaussian Splatting）~\cite{2406_01916} 提出语义特征网格映射方案，将多视角 CLIP 特征（基于 SAM 掩码提取）存储在特征网格中，并通过低维映射训练语义场，实现高分辨率下的实时开放词汇查询，速度比 LERF 快 98 倍。

\textbf{动机。}LangSplat 为每个高斯直接存储高维语言特征，在高分辨率场景中内存开销巨大且查询缓慢；FastLGS 用特征网格存储映射关系，通过低维映射间接查询，大幅降低内存和计算开销。

\paragraph{Dr.\,Splat：直接语言嵌入注册} Dr.\,Splat（2025.02，Direct Registration Splatting）~\cite{2502_16652} 提出直接语言嵌入注册（Direct Language Embedding Registration）策略，绕过渲染过程而直接将 CLIP 嵌入与三维高斯关联。

\textbf{动机。}LangSplat 等方法通过渲染-比较循环优化语言特征，$\alpha$-混合导致不同对象的特征互相干扰；Dr.\,Splat 绕过渲染直接注册，为每个像素射线上的主导高斯分配语言特征，避免了渲染导致的语义模糊。该方法为每个像素射线上的主导高斯分配语言特征，并引入在通用大规模图像数据上训练的乘积量化（Product Quantization）来紧凑表示嵌入，无需逐场景优化，在三维语义分割、目标定位和选择任务上显著优于依赖渲染的方法。

%---------------------------------------------------------------------
\paragraph{语言嵌入三维高斯表示}

LangSplat（2023.12）~\cite{2312_16081} 将开放词汇三维场景理解从 NeRF 体积渲染范式迁移到 3DGS 可微分光栅化框架，首次实现了支持任意文本查询的三维语言场构建。在 LangSplat 之前，LERF 通过在 NeRF 密集体中存储 CLIP 特征来建立语言场，但体渲染的高计算开销严重制约了查询效率；而 Feature 3DGS 虽已将语义特征蒸馏引入 3DGS，却主要面向分割和编辑任务，未提供端到端的开放词汇查询能力。LangSplat 的核心动机正是填补这一空白：利用 3DGS 显式点云表示天然的高效渲染优势，为三维空间中的文本驱动定位与分割构建统一的语言场。

在实现层面，LangSplat 的技术路线包含三个紧密耦合的阶段。第一阶段，利用 SAM（Segment Anything Model）为每帧训练图像生成过完备的二维掩码集合，再通过 CLIP 视觉编码器为每个掩码区域提取与文本语义对齐的高维嵌入向量。由于 CLIP 直接在像素级操作会产生大量冗余且语义模糊的特征，SAM 的语义级掩码提供了天然的分组粒度，显著提高了特征质量。第二阶段，将二维语言特征提升到三维：为每个三维高斯基元 $\mathbf{g}_i$ 附加一个与 CLIP 特征维度相同的语言属性向量 $\mathbf{l}_i \in \mathbb{R}^D$（$D$ 通常为 512），通过 3DGS 的 $\alpha$-混合 光栅化机制将三维语言特征投影回二维，与第一阶段获得的二维 CLIP 特征计算损失并反向传播优化：

\begin{equation}
  \mathcal{L}_{\text{lang}} = \sum_{\mathbf{p}} \left\| \hat{\mathbf{F}}_{\text{CLIP}}(\mathbf{p}) - \mathbf{F}_{\text{CLIP}}(\mathbf{p}) \right\|_2^2,
\end{equation}

其中 $\hat{\mathbf{F}}_{\text{CLIP}}(\mathbf{p}) = \sum_{i} \alpha_i \prod_{j=1}^{i-1}(1-\alpha_j) \mathbf{l}_i$ 是沿像素射线 $\mathbf{p}$ 通过 $\alpha$-混合 渲染的语言特征。第三阶段，为缓解 512 维 CLIP 特征直接存储在数百万高斯上带来的巨大显存压力和过拟合风险，LangSplat 引入场景专属自编码器（Scene-Aware AutoEncoder），将高维语言特征压缩为低维潜表示 $\mathbf{z}_i \in \mathbb{R}^d$（$d \ll D$，实践中取 $d=48$），训练时同时在潜空间和原始空间监督重建质量，推理时仅需解码查询点的特征，从而兼顾内存效率与查询精度。

经过上述流程，LangSplat 在 LERF-Multi 和 Mip-NeRF 360 数据集上以显著优势超越了 LERF，首次证明了基于 3DGS 的语言场在效率和精度上均可优于 NeRF 方案。然而其局限性也同样明确：场景专属自编码器需要逐场景训练，SAM + CLIP 的多阶段流水线计算开销较大，且高斯基元的语言属性完全依赖二维监督信号蒸馏而缺乏三维几何约束，这些瓶颈催生了后续一系列改进工作。

LangSplatV2（2025.07）~\cite{2507_07136} 针对 LangSplat 的速度瓶颈进行优化。

\textbf{动机。}LangSplat 的场景专属自编码器是推理速度的主要障碍，每次查询都需解码高维特征；LangSplatV2 用稀疏编码替代自编码器，假设每个高斯是全局词典中的稀疏编码，完全消除解码器需求，实现 476.2 FPS 的高维特征渲染。该工作发现重量级解码器是推理速度的主要障碍，转而假设每个高斯是全局词典中的稀疏编码，学习三维稀疏系数场从而完全消除解码器需求。通过 CUDA 优化的稀疏系数 splatting，实现 476.2 FPS 的高维特征渲染和 384.6 FPS 的三维开放词汇查询，分别比 LangSplat 快 42 倍和 47 倍。

\paragraph{Gen-LangSplat：通用自编码器替代} Gen-LangSplat（2025.10，Generalized LangSplat）~\cite{2510_22930} 解决了 LangSplat 需要逐场景训练自编码器的瓶颈，用在大规模 ScanNet 数据集上预训练的通用自编码器替代场景专属自编码器，消除了逐场景优化步骤，在保持查询性能的同时大幅提升可扩展性。

\textbf{动机。}LangSplat 的场景专属自编码器需要为每个新场景从零训练，限制了可扩展性；Gen-LangSplat 用在大规模数据上预训练的通用自编码器，推理时直接应用于新场景无需额外训练。

\paragraph{SLGaussian：稀疏视角前馈语义场} SLGaussian（2024.12，Sparse-view Language Gaussian）~\cite{2412_08331} 面向稀疏视角场景，提出前馈式三维语义场构建方法。

\textbf{动机。}LangSplat 需要密集多视角输入才能构建高质量语言场；SLGaussian 面向仅两个视角的极端稀疏条件，通过视频跟踪确保 SAM 分割的多视角一致性，实现 30 秒内的场景推断。通过视频跟踪确保 SAM 分割的多视角一致性，并用低维索引表示高维 CLIP 特征，在两个视角的稀疏条件下实现 30 秒内的场景推断和 0.011 秒的开放词汇查询。

\paragraph{LangFlash：无位姿前馈语言特征场} LangFlash（2026.05）~\cite{2605_23287} 进一步将前馈范式推向无位姿稀疏图像输入。该框架在单次前向传播中预测高斯基元的几何和语义属性，并提出稀疏语义编码方案——全局语义词典与逐基元局部权重相结合——降低表示复杂度，建立了无位姿、语言对齐的三维重建新范式。

\paragraph{SCOUP：稀疏码本提升加速} 三维语言高斯的稀疏编码加速方向上，SCOUP（Sparse COde UPlifting）（2026.05）~\cite{2605_13600} 通过将语言表示学习与三维高斯优化解耦，在二维图像区域学习稀疏码本表示，再通过加权稀疏聚合将系数提升至三维高斯。该方法实现高达 400 倍的训练加速，训练期间内存效率提高 3 倍，同时在多个基准上匹配或超越现有方法的查询精度。

%---------------------------------------------------------------------
\paragraph{多视角语义一致性与特征注册}

二维视觉-语言模型固有的多视角不一致性，是影响三维语义场质量的核心障碍。LangSplat 通过 SAM 掩码蒸馏 CLIP 特征并学习低维语言场来部分规避此问题，但其依赖的 SAM 掩码本身在多视角间缺乏显式一致性约束，后续工作从多个维度对此进行了改进。

\paragraph{Language Embedded 3D Gaussians：密集像素级特征} 早期方法 Language Embedded 3D Gaussians（2023.11）~\cite{2311_18482} 与 LangSplat 同期提出，走了一条不同的路线。

\textbf{动机。}LangSplat 采用 SAM 掩码区域裁剪后提取 CLIP 特征，掩码粒度影响特征质量；该方法直接从 CLIP 和 LSeg 的密集像素级特征出发，通过量化压缩和不确定性学习的平滑嵌入来处理多视角不一致性。LangSplat 采用 SAM 掩码区域裁剪后提取 CLIP 特征，而该方法直接从 CLIP 和 LSeg 的密集像素级特征出发，通过专用量化方案将高维特征压缩到低维码本空间以缓解内存需求，并引入基于不确定性学习的平滑嵌入过程来对抗多视角特征不一致性和点状表示的高频偏置。

\paragraph{OpenGaussian：三维点级开放词汇理解} OpenGaussian（2024.06）~\cite{2406_02058} 则在表示粒度上超越了 LangSplat 的区域级查询能力，首次实现基于 3DGS 的三维点级开放词汇理解。

\textbf{动机。}LangSplat 将 CLIP 特征蒸馏到每个高斯的低维向量中，查询精度受限于区域级的粗糙匹配；OpenGaussian 学习 6 维实例特征并通过由粗到精的码本离散化，实现精确的三维点级语义查询。其核心差异在于：LangSplat 将 CLIP 特征蒸馏到每个高斯的低维向量中，而 OpenGaussian 学习 6 维实例特征 $f \in \mathbb{R}^6$（无需额外的维度压缩或量化步骤），通过无跨帧关联的 SAM 掩码训练具有三维一致性的实例特征，并提出由粗到精的两阶段码本进行特征离散化，最后通过实例级三维-二维特征关联将三维点链接到 CLIP 特征。

\paragraph{CLIP-GS：室内语义一致性} CLIP-GS（2024.04，Contrastive Language-Image Pre-training for Gaussian Splatting）~\cite{2404_14249} 进一步聚焦于室内场景的语义一致性问题，其关键洞察是 LangSplat 的高维语言特征未考虑物体内部语义的自然统一性。

\textbf{动机。}LangSplat 的高维语言特征在物体内部存在不必要的方差——同一张椅子的不同部位可能获得略微不同的语言特征，导致查询不稳定；CLIP-GS 利用物体内语义的自然统一性学习紧凑表示，并通过三维一致性正则化约束语义一致性。为此，CLIP-GS 提出语义属性紧凑化（SAC）和三维一致性正则化（3DCR）：SAC 利用物体内语义的自然统一性学习紧凑表示，3DCR 同时在二维和三维域施加语义一致性约束。在 ScanNet 和 Replica 上分别取得 21.20\% 和 13.05\% 的 mIoU 提升，渲染速度超过 100 FPS，在效率与精度上均优于 LangSplat。

\paragraph{VALA：可见性感知语言聚合} VALA（Visibility-Aware Language Aggregation）（2025.09）~\cite{2509_05515} 从光栅化机理出发，关注两个被 LangSplat 忽视的问题：背景高斯与主导前景高斯获得相同特征、以及多视角语言嵌入的噪声。

\textbf{动机。}LangSplat 的 $\alpha$-混合将所有贡献高斯平等对待，背景高斯虽贡献小但获得与前景高斯相同的语言特征训练信号，引入语义噪声；VALA 通过可见性感知门控仅保留对渲染像素有实质贡献的高斯，并利用余弦空间的流式加权几何中值融合多视角噪声特征。LangSplat 的 $\alpha$-混合 特征渲染将所有贡献高斯平等对待，而 VALA 通过可见性感知门控仅保留对渲染像素有实质贡献的高斯，并提出余弦空间中的流式加权几何中值来融合多视角噪声特征，从根本上改善了多视角特征聚合的质量。

\paragraph{LaGa：视角依赖语义建模} LaGa（Language Gaussians）（2025.05）~\cite{2505_24746} 在语义建模维度上更进一步，首次系统研究三维场景中的视角依赖语义（view-dependent semantics）现象。

\textbf{动机。}LangSplat 假设语义是视角无关的，但实际场景中同一物体从不同视角可能呈现不同语义（如近距离看到的是"木质纹理"而远距离看到的是"桌面"）；LaGa 通过将场景分解为物体并基于多视角语义加权聚合，首次系统建模这种视角依赖语义现象。该方法通过将场景分解为物体、聚类语义描述子并基于多视角语义进行加权聚合，在 LERF-OVS 数据集上取得 +18.7\% mIoU 的显著提升。

\paragraph{GAGS：分割粒度一致性} GAGS（2024.12，Granularity-Aware Gaussian Splatting）~\cite{2412_13654} 从分割粒度角度解决多视角不一致性，指出 LangSplat 所依赖的 SAM 掩码粒度在不同视角下的不稳定性是特征不一致的重要来源。

\textbf{动机。}SAM 掩码的粒度因视角距离而异——近处物体被过度细分为小掩码，远处物体被合并为大掩码，导致同一物体在不同视角的特征不一致；GAGS 将提示点密度与相机距离关联以稳定分割粒度。GAGS 提出将 SAM 的提示点密度与相机距离关联以改善多视角分割一致性，并学习粒度因子以无监督方式选择多视角一致的特征用于蒸馏，从而在特征提取源头解决一致性问题。

\paragraph{econSG：SAM-CLIP 互精炼} econSG（2025.04，Efficient Consistent Semantic Gaussian Splatting）~\cite{2504_06003} 沿着类似的思路更进了一步。

\textbf{动机。}LangSplat 在二维视角上分别提取 SAM 和 CLIP 特征再独立蒸馏到三维，两个独立的蒸馏过程各自引入不一致性；econSG 通过置信区域引导的正则化让 SAM 和 CLIP 相互精炼，在融合后的三维特征上直接降维，避免了二维视角上分别操作导致的多视角不一致。不同于 LangSplat 在二维视角上分别提取 SAM 和 CLIP 特征再蒸馏到三维，econSG 通过置信区域引导的正则化（CRR）相互精炼 SAM 和 CLIP 以获得精确的语义特征，并在融合后的三维特征上直接降维，避免在二维视角上分别操作导致的多视角不一致，实现了训练效率与语义质量的双重最优。

\paragraph{LangSurf：语言嵌入表面场} LangSurf（2024.12，Language-Embedded Surface）~\cite{2412_17635} 则从几何与语义的耦合角度提出了根本性的改进。

\textbf{动机。}LangSplat 的语言特征自由漂浮在三维空间中，与物理表面无显式关联，语义查询可能定位到空间中不对应任何表面的"幽灵"位置；LangSurf 通过几何监督将语言高斯平铺在物体表面上，确保语义特征附着于真实几何。LangSplat 的语言特征自由漂浮在三维空间中，与物理表面无显式关联；LangSurf 提出语言嵌入表面场（Language-Embedded Surface Field），通过几何监督和对比损失将语言高斯平铺在物体表面上，确保语义特征附着于真实几何。同时引入层次化上下文感知模块，不像 LangSplat 仅依赖 SAM 掩码提取局部特征，而是先提取图像级全局特征再分层掩码池化，在二维和三维开放词汇分割上大幅超越 LangSplat。

\paragraph{GLS：几何感知统一重建与分割} GLS（Geometry-aware 3D Language Gaussian Splatting）（2024.11）~\cite{2411_18066} 在 LangSurf 的几何-语义联合框架基础上进一步统一了三维表面重建与开放词汇分割。

\textbf{动机。}LangSurf 将语义特征附着于表面但未利用表面几何来精炼语义；GLS 利用表面法线先验分析高斯法线误差来精炼物体结构，同时用 DEVA 掩码（一种带时序一致性的分割器）替代普通 SAM 掩码以提升视角一致性。GLS 利用预训练的表面法线估计模型（DSINE）提供法线先验，通过分析高斯法线误差来精炼物体结构，同时用 CLIP 特征引导实例特征并用 DEVA 掩码（一种带时序一致性的分割器）替代 LangSplat 的普通 SAM 掩码以保持视角一致性，在 MuSHRoom、ScanNet++ 和 LERF-OVS 数据集上均达到最优。

\paragraph{3D Vision-Language Gaussian Splatting：跨模态光栅化器} 3D Vision-Language Gaussian Splatting（2024.10）~\cite{2410_07577} 最后从光栅化层面提出创新。

\textbf{动机。}LangSplat 将颜色和语言特征视为相互独立的模态，使用相同的 $\alpha$-混合公式渲染两者；但语义光栅化与颜色光栅化有本质差异——半透明对颜色有意义但对语义无意义（一个半透明的红色物体仍然是"红色"，但半透明区域的语义归属应该更明确）。LangSplat 将颜色和语言特征视为相互独立的模态，使用相同的 $\alpha$-混合 公式渲染两者；而该方法观察到语义光栅化与颜色光栅化具有本质差异（例如半透明对颜色有意义但对语义无意义），据此提出跨模态光栅化器，使用模态融合和平滑语义指示器增强语义光栅化，并通过相机视角混合技术改善已有视角与合成视角之间的语义一致性。

%---------------------------------------------------------------------
\paragraph{FAST-Splat：直接语义编码} FAST-Splat（2024.11）~\cite{2411_13753} 从表示范式上与 LangSplat 形成了鲜明对比：LangSplat 先用 SAM+CLIP 提取高维语言特征再通过自编码器压缩至低维空间，而 FAST-Splat 采用自底向上的方法直接为每个高斯分配语义编码，结合哈希表实现开放词汇语义相似度查询。

\textbf{动机。}LangSplat 的 SAM+CLIP 多阶段流水线和场景专属自编码器计算开销大且训练慢；FAST-Splat 跳过特征蒸馏和压缩，直接编码语义到哈希表，训练速度快 6--8 倍，渲染速度快 18--51 倍。这种直接编码策略跳过了特征蒸馏和压缩的繁琐流水线，训练速度比竞争方法快 6--8 倍，渲染速度快 18--51 倍，GPU 内存减少约 6 倍，同时保持或超越现有方法的分割性能。

\paragraph{LUDVIG：免训练图扩散} LUDVIG（2024.10，Language-Unsupervised Distillation via Graph）~\cite{2410_14462} 走得更远，提出完全免训练的二维特征提升方法。

\textbf{动机。}LangSplat 需要逐场景优化语言场参数，训练耗时；LUDVIG 通过图扩散机制利用 DINOv2 的成对相似性和三维几何直接精炼三维特征，仅使用 DINOv2 特征即获得与 SAM 相当的分割结果，从根本上消除了训练阶段。LangSplat 需要逐场景优化语言场参数，而 LUDVIG 通过图扩散机制精炼三维特征，利用 DINOv2 的成对相似性和三维几何改善粗分割掩码，仅使用 DINOv2 特征即获得与 SAM 相当的分割结果，从根本上消除了训练阶段。

\paragraph{Gradient-Weighted Feature Back-Projection：梯度反投影} Gradient-Weighted Feature Back-Projection（2024.11）~\cite{2411_15193} 提出另一种免训练方法。

\textbf{动机。}LangSplat 需要逐场景优化语言场参数；该方法发现一个简洁的数学性质——渲染颜色对高斯颜色的梯度等于该高斯在 $\alpha$-混合 中的权重，据此可用推理时梯度直接反投影 2D 特征到 3D 高斯，无需任何训练。其核心公式揭示了与 LangSplat 的本质差异：对渲染颜色 $C(x,y)$ 关于第 $k$ 个高斯的颜色 $c_k$ 求导，梯度 $\partial C / \partial c_k = \alpha_k T_k$ 正好等于该高斯在渲染中的权重。据此，该方法利用推理时的梯度反传计算每个高斯的影响力权重，将二维特征反投影到预训练的三维高斯中，无需任何训练即可在二维和三维分割上取得高质量结果。

\paragraph{SLAG：多 GPU 并行嵌入} SLAG（2025.05，Scalable Language-Embedded Gaussian）~\cite{2505_08124} 从系统架构层面突破单 GPU 限制，提出多 GPU 并行的语言增强高斯框架。

\textbf{动机。}LangSplat 的逐场景损失函数优化在大规模场景中计算瓶颈严重；SLAG 通过归一化加权平均直接计算嵌入，消除损失函数需求，在多 GPU 配置上实现数量级加速。不同于 LangSplat 的逐场景损失函数优化，SLAG 通过归一化加权平均直接从三维高斯场景参数计算嵌入，消除损失函数计算需求，在 16 GPU 配置上实现比 OpenGaussian 快 18 倍的嵌入计算，为大规模部署提供了工程级解决方案。

\paragraph{Training-Free Hierarchical Scene Understanding：超点图} Training-Free Hierarchical Scene Understanding（2025.04）~\cite{2504_13153} 则在表示结构上进行创新：LangSplat 将语言特征附加到每个高斯上，而该方法构建超点图（superpoint graph）直接从高斯基元出发划分场景。

\textbf{动机。}LangSplat 为每个高斯独立存储语言特征，在大场景中内存开销巨大；该方法通过超点图将场景分割为少量语义均匀的超点，在超点级别而非高斯级别处理语义，使语义场重建速度比现有方法快 30 倍以上。通过对比高斯分区将场景分割为语义均匀的超点，再用层次化合并策略按 SAM 掩码层级聚合为多级超点图。设计高效的重投影策略将二维语义特征提升到超点上，避免昂贵的多视角迭代训练，语义场重建速度比现有方法快 30 倍以上。

%---------------------------------------------------------------------
\paragraph{Hi-LSplat：层次化语言高斯}

三维场景中物体固有的层次结构（如物体-部件关系）要求语义理解支持多粒度查询。LangSplat 的语言场是扁平的——所有高斯在同一语义层级上编码语言特征，无法区分"椅子"与"椅腿"的层次关系。Hi-LSplat（2025.06，Hierarchical Language Splatting）~\cite{2506_06822} 针对这一不足提出层次化语言高斯表示，通过构建三维层次语义树和分层实例聚类解决二维语义特征导致的视角不一致问题，并引入实例级和部件级对比损失捕捉全面的层次语义。

\textbf{动机。}LangSplat 的语言场是扁平的，无法区分"椅子"与"椅腿"的层次关系；Hi-LSplat 通过三维层次语义树显式建模物体-部件的层次关系，使查询可在不同语义粒度上进行。该工作还构建了两个层次化语义数据集来评估模型区分不同语义层次的能力。

\paragraph{TreeGaussian：显式树状拓扑} TreeGaussian（2026.04）~\cite{2604_03309} 进一步将层次结构显式建模为树状拓扑。不同于 Hi-LSplat 通过聚类隐式发现层次关系，TreeGaussian 通过树引导的级联对比学习框架显式建模层次语义关系，构建多级对象树实现结构化学习，并提出两阶段级联对比策略从全局到局部逐步精炼特征表示，同时引入一致性分割检测（CSD）机制和基于图的去噪模块。

\paragraph{FMLGS：部件级语义偏差} FMLGS（2025.04，Fine-grained Multi-Level Gaussian Segmentation）~\cite{2504_08581} 侧重于部件级细粒度查询能力。

\textbf{动机。}LangSplat 的扁平语言场无法区分"椅子"和"椅腿"等层次语义，且语言查询本身存在歧义；FMLGS 通过语义偏差策略在父级物体的语义空间中插值定位子级部件，解决语言歧义问题。该方法基于 SAM2 构建一致的物体和部件级语义，并提出语义偏差策略通过插值细粒度目标的语义特征解决语言歧义，速度比 LERF 快 98 倍。与 TreeGaussian 的树状拓扑不同，FMLGS 通过语义偏差插值在父级物体的语义空间中定位子级部件。

\paragraph{SuperGSeg：结构化超高斯} SuperGSeg（2024.12，Super-Gaussian Segmentation）~\cite{2412_10231} 提出结构化超高斯（Super-Gaussians）概念，采用三阶段训练范式。

\textbf{动机。}为每个高斯独立编码语言特征需要数百万个向量，内存开销极大；SuperGSeg 通过可学习聚类将锚点分组为少量语义均匀的超高斯，仅需为每个超高斯学习一个语言向量，大幅降低内存开销。先用 Scaffold-GS 重建几何，再用可学习聚类网络将锚点分组为语义均匀的超高斯（$S \ll N'$），最后为每个超高斯学习语言特征 $\phi^l$，仅需 $S$ 个语言向量而非数百万高斯的独立编码。这一设计在中等 GPU 内存开销下实现层次化场景理解。

\paragraph{GOI：语义空间超平面} GOI（2024.05，Gaussian Open-vocabulary Instance）~\cite{2405_17596} 引入可优化的语义空间超平面，通过二维参照表达式分割模型微调特征空间中的分割边界，并用可训练特征聚类码本（TFCC，Trainable Feature Clustering Codebook）压缩高维语言嵌入，实现更精确的开放词汇查询定位。

\textbf{动机。}LangSplat 的语言特征未在特征空间中建立明确的分割边界，查询时依赖余弦相似度的经验阈值；GOI 用可优化的超平面显式定义分割边界，使开放词汇查询的定位更加精确。

%---------------------------------------------------------------------
\paragraph{ObjectGS：对象感知锚点框架}

区分三维空间中的独立物体实例是场景理解的关键能力。LangSplat 和早期的特征蒸馏方法主要在点级别或像素级别编码语义，缺乏显式的对象级身份建模，难以直接区分同一类别下的不同实例。ObjectGS（2025.07）~\cite{2507_15454} 提出对象感知框架，将每个物体建模为生成神经高斯的局部锚点并共享对象 ID。

\textbf{动机。}LangSplat 和早期特征蒸馏方法主要在点级别编码语义，缺乏显式的对象级身份建模，难以区分同一类别下的不同实例；ObjectGS 通过为每个锚点分配 one-hot 对象 ID，使同一物体的所有高斯共享一致的实例身份。与 LangSplat 对每个高斯独立编码语言特征不同，ObjectGS 为每个锚点分配 one-hot 对象 ID 编码，配合分类损失强制语义约束，使同一物体的所有高斯共享一致的实例身份。该方法还将三维重建与实例分割视为联合优化任务——语义信息帮助消除几何歧义，精确重建反过来提升分割质量——在开放词汇和全景分割任务上取得最优表现，同时无缝支持网格提取和场景编辑。

\paragraph{OP2GS：双不透明度表示} ObjectGS 通过共享 ID 建立实例一致性，但高斯的半透明特性和 $\alpha$-混合机制使得二维前景掩码无法一一映射到唯一的前景高斯集合，不准确的深度估计和重叠区域会产生``标签污染''。OP2GS（2026.05，Object-Point-2-Gaussian Splatting）~\cite{2605_20044} 针对此问题提出双不透明度（dual-opacity）表示：原始不透明度 $\alpha$ 负责视觉重建，专用实例不透明度 $\alpha^*$ 负责对象掩码渲染。

\textbf{动机。}高斯的半透明特性和 $\alpha$-混合机制使得二维前景掩码无法一一映射到唯一的前景高斯集合，不准确的深度估计和重叠区域会产生"标签污染"；OP2GS 从物理上解耦视觉存在与实例占据，用标量占据值替代高维特征编码身份。与此前使用高维特征向量编码身份的方法（如 OpenGaussian、Feature 3DGS）相比，OP2GS 用标量占据值 $\alpha^*$ 与整数标签 $l$ 的组合替代高维特征，从物理上解耦视觉存在与实例占据，利用三维高斯不透明度固有的视角一致性作为稳定的掩码渲染信号，消除误标高斯对渲染的影响。

\paragraph{ExtrinSplat：外建范式解耦几何与语义} ObjectGS 和 OP2GS 都将语义特征嵌入三维高斯并逐场景优化，存储开销随高斯数量线性增长。ExtrinSplat（2025.09，Extrinsic Splatting）~\cite{2509_22225} 则另辟蹊径，提出外建范式（extrinsic paradigm）将几何与语义完全解耦。

\textbf{动机。}ObjectGS 和 OP2GS 将语义特征嵌入高斯基元，存储开销随高斯数量线性增长；ExtrinSplat 将语义存储在高斯基元之外的对象组中，存储开销减少数个数量级，且无需训练即可完成场景适配。该方法不嵌入任何高维特征，而是将高斯聚类为多粒度重叠的三维对象组，并通过识别和排除边界歧义的中性点净化对象边界，再由视觉-语言模型解读生成轻量级文本假设。用户查询直接与独立对象组的特征匹配，无需为每个高斯点存储唯一语义向量，存储开销减少数个数量级，场景适配时间从 LangSplat 的数小时缩短至数分钟，且无需训练。

\paragraph{OpenGaFF（续）：几何条件化特征场} 前述方法侧重于实例身份建模，但对高斯几何与语义属性之间的耦合关系关注不足——高斯的位移或分裂可能导致语义特征的空间漂移。OpenGaFF（2026.05）~\cite{2605_06088} 引入高斯特征场将语义建模为高斯几何和外观的连续函数，使语义属性与空间位置自然对齐而非独立存储，并设计结构化码本和码本引导的注意力机制，通过查询嵌入与码本条目的相似匹配实现鲁棒的开放词汇推理。

\paragraph{GALA：自监督实例特征场} 在实例特征的学习策略上，GALA（2025.08，Group-Align-Language Attention）~\cite{2508_14278} 与上述方法形成鲜明对比：它不依赖二维掩码的跨视角对齐，而是通过自监督对比学习蒸馏场景特定的三维实例特征场。

\textbf{动机。}ObjectGS 和 OP2GS 等方法依赖二维掩码的跨视角对齐，但 SAM 掩码本身在不同视角间不一致；GALA 通过自监督对比学习直接在三维空间中学习实例特征，完全绕过二维掩码对齐的难题。GALA 基于 Scaffold-GS 架构，采用两阶段训练——第一阶段自监督重建场景几何和实例特征场，第二阶段渲染至二维并通过双可学习码本的交叉注意力模块将实例特征映射为泛化语言特征。线性注意力设计使仅用二维训练即可实现二维和三维无缝分割，码本与注意力权重熵损失强制一对一映射，确保实例内特征一致性并支持开放词汇查询。

\paragraph{CCL-LGS：显式监督对比码本} GALA 依靠自监督学习实例特征，而 CCL-LGS（2025.05，Contrastive Codebook Learning for Language Gaussian Splatting）~\cite{2505_20469} 则构建了显式监督的对比码本学习框架来系统解决多视角语义冲突问题。

\textbf{动机。}GALA 依靠自监督学习实例特征，缺乏显式的语义监督信号；CCL-LGS 构建显式监督的对比码本，通过强制类内紧凑和类间可分性系统解决多视角语义冲突。该方法先用零样本跟踪器对齐 SAM 生成的二维掩码获得跨视角一致的实例分组，再用 CLIP 提取鲁棒的多视角语义编码，最后通过对比码本学习强制类内紧凑和类间可分性。低维语言特征向量被集成到高斯基元中，通过可微分渲染和语义分割损失进行优化，使语义渗透到三维空间的同时保持与二维先验的一致性。

%---------------------------------------------------------------------
\paragraph{EmbodiedSplat：具身在线前馈重建}

逐场景优化的局限性推动了对泛化能力的研究。LangSplat 及后续方法均需对新场景从零开始数十分钟至数小时的训练，无法满足具身智能体在线感知和实时决策的需求。EmbodiedSplat（2026.03）~\cite{2603_04254} 首先从具身感知的需求出发，提出满足在线、实时、可泛化、大规模场景和开放词汇五项条件的在线前馈三维高斯语义重建框架。

\textbf{动机。}LangSplat 及后续方法均需对新场景从零开始数十分钟至数小时的训练，无法满足具身智能体在线感知和实时决策的需求；EmbodiedSplat 采取"二维提升至三维"的前馈路线，在重建过程中直接将像素级 CLIP 特征反投影至三维空间。与 LangSplat 采用的``三维渲染至二维再比较''（render-and-compare）范式不同，EmbodiedSplat 采取相反的``二维提升至三维''（2D-to-3D）路线：在在线重建过程中，像素级 CLIP 特征沿局部高斯三元组直接反投影至三维空间，形成高斯四元组。为控制在线累积的内存开销，该方法引入在线稀疏系数场和 CLIP 全局码本，将 CLIP 嵌入绑定到三维高斯的同时最小化内存消耗，并通过三维 U-Net 聚合高斯点云生成三维几何感知的 CLIP 特征。其轻量化变体进一步将推理速度提升至接近实时。

\paragraph{SemGS：双高斯稀疏视角重建} EmbodiedSplat 面向具身在线场景，而 SemGS（2026.03）~\cite{2603_02548} 聚焦于稀疏视角输入下的泛化语义场重建。SemGS 提出双高斯表示——每个输入像素关联两个互补高斯：颜色高斯用于辐射建模，语义高斯用于语义推理。两者共享三维位置和不透明度以保证几何一致性，但各自维护协方差矩阵和分支特定属性（颜色用球谐函数编码，语义用类别分布表示）。架构上采用双分支特征提取器，颜色分支使用 6 块 Swin Transformer，语义分支使用 3 块 Transformer 并额外增加 2 个残差块用于语义精炼，两分支共享浅层 CNN 骨干以利用颜色外观的结构和纹理线索增强语义推理。相机姿态被注入 Transformer 注意力机制中建模视角间几何关系，区域平滑损失进一步约束语义预测的局部一致性。

\paragraph{OVGaussian：跨场景泛化分割（数据驱动）} EmbodiedSplat 和 SemGS 通过前馈网络实现泛化，但它们依赖二维语义先验的直接注入，缺乏对三维空间关系的显式建模。OVGaussian（2025.01）~\cite{2501_00326} 从数据规模的角度切入，首次构建大规模三维高斯语义数据集 SegGaussian（包含 ScanNet++ 的 280 个场景和 Replica 的 8 个场景），提出泛化语义光栅化（GSR，Generalizable Semantic Rasterization）通过三维神经网络学习预测高斯语义属性——在大型高斯场景集合上训练后，网络可直接预测新场景中每个高斯点的语义向量。

\textbf{动机。}EmbodiedSplat 和 SemGS 通过前馈网络实现泛化但缺乏对三维空间关系的显式建模；OVGaussian 从数据驱动角度切入，通过大规模语义数据集训练跨场景泛化的语义光栅化网络。同时设计跨模态一致性学习（CCL）框架，利用二维图像和三维高斯的开放词汇标注训练，确保三维语义表示与二维视觉和文本嵌入的对齐，实现跨场景、跨域和跨视角的泛化分割。

\paragraph{Scene-Agnostic Object-Centric：对象中心全局码本} 上述方法在像素级或点级进行语义推理，但缺乏显式的对象级表示。Scene-Agnostic Object-Centric Representation Learning（2026.04）~\cite{2604_09045} 提出数据集级别的对象中心监督方案。

\textbf{动机。}所有前述方法在像素级或点级进行语义推理，缺乏显式的对象级表示；该方法通过解耦槽注意力学习全局、场景无关的对象码本，提供跨场景一致的对象监督信号，从根本上解决二维基础模型掩码跨视角不一致的问题。该方法基于 GOLD 框架，利用 DINO 特征通过解耦槽注意力（Disentangled Slot Attention）学习全局、场景无关的对象码本——背景编码器处理 DINO 特征生成背景槽，同时码本维护独立的对象条目。学习到的全局码本提供跨视角和跨场景一致的、以身份锚定的对象监督信号，无需逐场景微调即可进行对象识别和分割，从根本上解决了二维基础模型掩码跨视角不一致的问题。

\paragraph{ProFuse：渲染免监督直接注册} 前三类方法在特征表示和训练范式上寻求泛化，但仍依赖可微分渲染进行语义监督。ProFuse（2026.01，Proposal-Fused Gaussian Segmentation）~\cite{2601_04754} 则提出完全绕过渲染监督的替代方案。

\textbf{动机。}所有基于可微分渲染的方法存在渲染-比较过程中的语义失真——$\alpha$-混合会将不同对象的特征混合在一起；ProFuse 通过直接注册（而非渲染-比较）将跨视角掩码分组为三维上下文提案，避免渲染造成的语义模糊。该方法通过密集对应引导的预注册阶段初始化高斯几何并链接跨视角分割掩码形成三维上下文提案（Context Proposals），每个提案记录推断为同一场景内容的跨视角掩码分组。随后通过上下文引导的注册阶段，利用基于透射率和不透明度的可见性加权为每个提案计算全局语言特征并分配给对应高斯。这种``直接注册''策略避免了渲染-比较的语义失真问题，每场景约 5 分钟即可完成，兼具语义保真度和效率优势。

\paragraph{LESV：稀疏体素光栅化特征注册} 尽管 ProFuse 绕过了渲染监督，其底层仍使用非结构化的三维高斯表示，在高斯重叠区域存在空间模糊问题。LESV（2026.04，Language-Embedded Sparse Voxel rasterization）~\cite{2604_01388} 换用稀疏体素光栅化（SVRaster）从根本上解决这一结构性问题。

\textbf{动机。}3DGS 的非结构化椭球在高斯重叠区域存在空间模糊问题，$\alpha$-混合导致语义渗透伪影；LESV 用稀疏体素替代非结构化椭球，通过严格深度排序消除弹出伪影，实现确定性的特征注册。与 3DGS 的非结构化椭球不同，SVRaster 将场景表示为高分辨率层次网格中的稀疏体素集合，通过视角依赖的 Morton 码排序实现严格深度排序，从根本上消除非结构化 3DGS 中常见的弹出伪影。在此基础上，LESV 通过单目深度和法线先验正则化建立稳定的几何基础，利用三线性插值体密度和光线-体素交距计算不透明度，实现确定性的、置信度感知的特征注册，有效抑制三维高斯中常见的语义渗透伪影。

%---------------------------------------------------------------------
\paragraph{开放词汇分割与检测系统}

面向实际应用的开放词汇系统需要处理实例分割、对象检测和场景查询等复合任务。与 Feature 3DGS（2023）~\cite{feature3dgs} 和 LangSplat（2023）~\cite{langsplat} 专注于像素级语义分割不同，本节的方法将开放词汇能力扩展到实例级识别、三维检测和交互式场景推理。

\paragraph{OpenSplat3D：实例对比学习} OpenSplat3D（2025.06）~\cite{2506_07697} 将 SAM 实例掩码与对比损失公式结合引导实例特征，利用视觉-语言模型的语言嵌入支持灵活的文本驱动实例识别。

\textbf{动机。}LangSplat 依赖 CLIP 特征直接蒸馏，但 CLIP 特征的像素级操作产生大量冗余且语义模糊的表示；OpenSplat3D 引入低维（8维）实例嵌入加方差正则化，在 $\alpha$-混合渲染中保持特征一致性，同时大幅降低特征维度。与 LangSplat 依赖 CLIP 特征直接蒸馏不同，OpenSplat3D 为每个高斯基元引入 8 维实例嵌入 $f_n^i \in \mathbb{R}^8$，通过方差正则化损失 $\mathcal{L}_{\text{var}} = \frac{1}{N}\sum_n \text{Var}(\{f_n^i\})$ 确保 $\alpha$-混合渲染中特征的一致性，损失权重设为 $\lambda_{\text{var}} = 0.5$，实例对比损失权重 $\lambda_{\text{inst2d}} = 0.1$，负样本边界 $m=1.0$。在 LERF-mask、LERF-OVS 和 ScanNet++ 数据集上展示了有效性，单场景优化耗时约 20--45 分钟（RTX 3090）。

\paragraph{COS3D：协作式实例-语言场} COS3D（2025.10，Collaborative Object-Scene 3D）~\cite{2510_20238} 转而构建实例场和语言场的协作式提示-分割框架（collaborative field），通过实例到语言的特征映射和高效两阶段训练策略捕捉内在关系，推理时通过自适应语言到实例的提示精炼实现高质量分割。

\textbf{动机。}LangSplat 只有单一语言场，缺乏对实例身份的显式建模；COS3D 构建实例场和语言场两个子场，通过显式特征桥接使两者相互增强，超越单一语言场的表达力。该方法在实例-语言两个场之间建立显式的特征桥接，超越了 LangSplat 单一语言场的表达力。

\paragraph{OpenInsGaussian：上下文感知聚合} OpenInsGaussian（2025.10）~\cite{2510_18253} 在此基础上进一步聚焦实例高斯分割中的两个核心问题——单掩码上下文不足和多视角特征融合不一致。

\textbf{动机。}OpenSplat3D 的逐掩码特征缺乏上下文信息（单个掩码内的特征无法区分外观相似的不同实例），且多视角特征简单平均忽略了视角质量差异；OpenInsGaussian 通过 ConvNeXt-Base CLIP 提供更丰富的上下文特征，并通过注意力驱动聚合自适应加权多视角信号。分别提出上下文感知特征提取（使用 ConvNeXt-Base CLIP 而非 ViT-B/16，以更好地处理大分辨率输入）和注意力驱动特征聚合模块。其 6 维实例特征 $f \in \mathbb{R}^6$ 通过对比学习训练，注意力聚合根据多视角语义相似度自适应加权而非简单平均，在 ScanNet 的 19/15/10 类子集上取得优于 OpenGaussian 和 LangSplat 的 mIoU。

\paragraph{VoteSplat：Hough 投票机制} VoteSplat（2025.06）~\cite{2506_22799} 走了一条截然不同的技术路线：将经典 Hough 投票机制与 3DGS 结合，利用 SAM 提取实例并生成二维投票图，将空间偏移向量 $\delta_i \in \mathbb{R}^3$ 嵌入高斯基元构建三维空间投票，通过深度畸变正则化 $\mathcal{L}_{\text{depth}}$ 约束偏移向量沿深度方向的聚集。

\textbf{动机。}LangSplat 依赖 512 维 CLIP 特征，训练成本高昂；VoteSplat 将语义匹配转化为实例 ID 投影，大幅降低特征维度和训练成本，同时利用 Hough 投票的空间聚集性实现更鲁棒的实例定位。与 LangSplat 依赖 512 维 CLIP 特征不同，VoteSplat 将语义匹配转化为实例 ID 投影，大幅降低与高维 CLIP 特征相关的训练成本，在 3D-OVS 和 LERF-OVS 上优于 LangSplat 和 OpenGaussian。

\paragraph{Beyond Averages：嵌入包检索} Beyond Averages（2025.09）~\cite{2509_12938} 则从根本上绕开了语义的可微分渲染瓶颈。

\textbf{动机。}Feature 3DGS 和 LangSplat 通过 $\alpha$-混合渲染传播语义特征，但混合过程会将不同对象的特征模糊在一起；Beyond Averages 绕过渲染，直接在对象级高斯组上聚合多视角 CLIP 特征，避免了渲染导致的语义模糊。这正是 Feature 3DGS 和 LangSplat 的核心技术路径。利用 Gaussian Grouping 预分解的对象级高斯，为每个对象组聚合多视角 CLIP 特征创建全面的``嵌入包''（bag of embeddings），推理时以 $k=5$ 的 top-relevancy 均值计算匹配分数，实现准确的对象级开放词汇检索。该方法避免了 $\alpha$-混合渲染中的语义模糊，在 3D-OVS 和 LERF 数据集的 3D 对象提取任务上表现突出。

\paragraph{NVSMask3D：新视角硬视觉提示} NVSMask3D（2025.04，Novel View Synthesis Mask 3D）~\cite{2504_14638} 同样避免了端到端的二维-三维特征蒸馏。

\textbf{动机。}端到端特征蒸馏方法需要逐场景优化且受限于 2D-3D 对齐误差；NVSMask3D 利用 3DGS 的高效新视角合成能力生成多样化视角图像，直接用 2D 硬视觉提示技术增强 CLIP 的零样本识别能力，无需任何 2D-3D 优化或微调。利用 3DGS 的高效新视角合成能力，通过相机插值围绕目标物体质心生成多样视角（基于 Weiszfeld 算法计算几何中位数），用分割高斯作为对象描述性硬视觉提示增强 CLIP 的零样本识别能力。核心创新在于将 2D 硬视觉提示技术迁移到 3D 场景中，通过加权特征平衡（WFB）融合选定和插值视角的特征。无需二维-三维优化或微调，在 Replica 和 ScanNet++ 上显著超越 OpenMask3D 基线。

\paragraph{GaussianVLM：原生三维视觉-语言模型} GaussianVLM（2025.07）~\cite{2507_00886} 提出首个基于高斯的原生三维视觉-语言模型（VLM，Vision-Language Model），采用 SceneSplat 的语言感知骨干为每个高斯预测语言特征，实现早期模态对齐——这比 LangSplat 的后期特征投影更深入。其双稀疏化器设计（任务引导稀疏化器 + 位置引导稀疏化器）将密集高斯表示蒸馏为紧凑的任务相关标记，实现无检测器的区域兴趣（ROI）提取，在室内场景理解任务上展示了端到端语言-三维推理能力。

\paragraph{ChatSplat：多轮对话式理解} ChatSplat（2024.12）~\cite{2412_00734} 进一步扩展交互维度，构建三层交互系统——对象级、视角级和场景级——支持基于聊天的三维空间交互。

\textbf{动机。}LangSplat 仅支持简单的文本查询（"找到椅子"），无法进行多轮对话式推理（"那个红色杯子旁边的蓝色物体是什么"）；ChatSplat 引入 LLM 语言嵌入和可学习归一化，将交互从单次查询扩展到多轮对话式场景理解。与 LangSplat 的 CLIP 嵌入分布不同，LLM 语言嵌入（如 LLaVA）具有更复杂多样的分布特征，ChatSplat 引入可学习归一化技术（learnable normalization）解决此挑战，将交互从简单的文本查询扩展到多轮对话式三维场景理解。

\paragraph{ReasonGrounder：认知推理定位} ReasonGrounder（2025.03）~\cite{2503_23297} 将开放词汇能力推向更高层的认知推理。

\textbf{动机。}LangSplat 等方法仅支持直接语义查询（"找到水果"），无法处理需要推理的隐式描述（"被杯子挡住的水果"）；ReasonGrounder 结合大视觉-语言模型的推理能力和层次化三维特征高斯场，实现认知推理级的三维目标定位。结合大视觉-语言模型和层次化三维特征高斯场，根据物理尺度进行自适应分组，支持隐式语言描述（如``被杯子挡住的水果''）下的三维目标定位与推理，即使目标被遮挡也能准确识别。该方法通过尺度感知的层次化分组解决了 LangSplat 在处理细粒度语义关系时的粒度不匹配问题。

%---------------------------------------------------------------------
\paragraph{特定领域的开放词汇应用}

开放词汇场景理解正在向特定领域渗透。这些领域面临的共同挑战是 LangSplat 和 Feature 3DGS 的通用设计难以应对的特殊条件：动态变形（手术）、极端尺度差异（果园水果）、退化成像（水下）、密集遮挡（以自我为中心视角）等。

\paragraph{SurgTPGS：手术场景语义} SurgTPGS（2025.06）~\cite{2506_23309} 将文本提示式高斯引入手术场景，结合 SAM 和视觉-语言模型的三维语义特征学习策略以及语义感知变形跟踪，实现手术工具和解剖结构的实时文本提示三维查询。与通用 LangSplat 使用静态高斯不同，SurgTPGS 采用灵活变形建模方案（FDM），通过傅里叶和多项式基函数 $\Phi(\tau; \omega, \mu, \sigma)$ 参数化形变，将语义特征与变形高斯绑定，使文本提示能在动态手术过程中持续追踪解剖结构。

\paragraph{CountingFruit：果园水果计数} CountingFruit（2025.06）~\cite{2506_01109} 面向果园场景，构建语言引导的三维水果计数框架 FruitLangGS，在高斯 $G_i = \{\mu_i, \Sigma_i, c_i, o_i, f_i\}$ 中嵌入压缩的 CLIP 对齐语义向量 $f_i$，通过双阈值余弦相似度机制检索目标提示相关的高斯——仅当余弦相似度同时超过上界阈值（确定正样本）且不落入模糊区间时才计入匹配——在九个数据集上实现高达 99.7\% 的实例计数召回率。与 LangSplat 面向全场景的语义分割不同，FruitLangGS 的自适应 splatting 策略针对果园场景中树冠不规则、果实密集分布的几何特点优化高斯基元分配。

\paragraph{SWAGSplatting：水下退化成像} SWAGSplatting（2025.09）~\cite{2509_00800} 面向水下退化成像条件，将 CLIP 语义先验集成到 3DGS 中增强鲁棒性，提出语义一致性损失对齐几何重建与高级语义，以及自适应高斯基元重分配策略缓解低能见度区域的表示不足。水下场景的色彩偏差和散射退化直接威胁 LangSplat 所依赖的 CLIP 特征质量，SWAGSplatting 通过显式建模语义-几何对齐来补偿。

\paragraph{EgoSplat：以自我为中心视角} EgoSplat（2025.03）~\cite{2503_11345} 面向以自我为中心（egocentric）场景，利用 SAM2 的分割和跟踪能力选择性聚合互补的跨视角特征，并提出实例感知的时空瞬态预测模块减少动态干扰物。与 LangSplat 假设静态场景不同，以自我为中心视频中频繁出现的手部遮挡和运动物体对语言场构建构成严重干扰，EgoSplat 的瞬态预测模块显式分离动态与静态成分。

\paragraph{PlanarGS：室内平面先验} PlanarGS（2025.10）~\cite{2510_23930} 面向室内场景，利用视觉-语言分割模型生成语言提示的平面先验（LP3），通过交叉视角融合和几何先验检查精炼区域提案，为高斯优化引入平面一致性约束和几何先验约束。LangSplat 缺乏对室内场景中平面结构（墙壁、地板、桌面）的几何先验，PlanarGS 将平面检测与语言特征融合统一在 3DGS 框架中。

\paragraph{Taking Language Embedded 3DGS into the Wild：互联网照片集} Taking Language Embedded 3DGS into the Wild（2025.07）~\cite{2507_19830} 将语言嵌入三维高斯扩展到无约束互联网照片集。

\textbf{动机。}LangSplat 假设受控拍摄条件（固定光照、无瞬态物体），但互联网照片集存在光照剧变、瞬态物体（行人、车辆）和不同外观模式；该方法通过多外观 CLIP 特征和瞬态/外观不确定性图，用不确定性感知的自编码器将这些挑战统一处理。提出多外观 CLIP 特征和瞬态/外观不确定性图引导优化，引入瞬态不确定性感知自编码器和后集成策略。与 LangSplat 假设受控拍摄条件不同，互联网照片集存在光照剧变、瞬态物体（行人、车辆）和不同外观模式，该方法通过不确定性感知的自编码器将这些挑战统一处理。

\paragraph{ReLaGS：层次化语言场景图} ReLaGS（2026.03，Relational Language Gaussian Splatting）~\cite{2603_17605} 构建层次化语言蒸馏高斯场景及其三维语义场景图，通过高斯剪枝精炼场景几何，鲁棒的多视角语言对齐策略聚合噪声二维特征。

\textbf{动机。}LangSplat 仅能进行语义查询，无法建模物体间的关系（如"桌子上的杯子在椅子旁边"）；ReLaGS 在层次化语言高斯上显式构建三维场景图，节点对应层次化实体，边编码语义关系，实现关系级场景理解。其核心是定义 $L$ 层抽象层次的层次化表示 $\mathcal{S}^{(1)}, \ldots, \mathcal{S}^{(L)}$，并在其上显式构建三维场景图——节点对应层次化实体，边编码语义关系（由大语言模型提升或图神经网络预测）。在开放词汇分割、场景图生成和关系引导检索上验证有效性，超越了 LangSplat 仅能进行语义查询的能力。

\paragraph{CAGS：局部图上下文传播} CAGS（2025.04）~\cite{2504_11893} 通过构建局部图在高斯间传播上下文特征，减少不一致粒度导致的噪声，并采用掩码中心对比学习平滑 SAM 衍生的跨视角特征。与 Feature 3DGS 的逐像素特征蒸馏不同，CAGS 的图传播机制利用高斯的空间邻域关系实现特征平滑，从根本上解决了跨视角 SAM 掩码不一致的问题。

%---------------------------------------------------------------------
\paragraph{三维场景问答与推理}

将语言理解从定位和分割拓展到更高层次的推理和问答是前沿方向。不同于 LangSplat 的开放词汇查询仅输出语义相关性图，本节的方法需要模型具备组合推理、时序理解和自由形式回答的能力。

\paragraph{SplatTalk：零样本三维视觉问答} SplatTalk（2025.03）~\cite{2503_06271} 利用泛化三维高斯框架（基于 FreeSplat 扩展）生成可直接输入预训练 LLM 的三维标记，实现仅依赖带位姿图像的有效零样本三维视觉问答。其关键创新在于特征压缩流水线：先用 LLaVA-OV（SigLIP 视觉编码器 + Qwen2 语言模型）提取 $27 \times 27 \times 3584$ 维高维特征，再通过自编码器将每个 3584 维特征向量压缩至 256 维（压缩比约 14 倍），以 256 维低维特征作为伪真值训练 3DGS 语言表示。推理时直接将三维高斯的语言特征送入 LLM 生成自由形式回答，无需点云或深度图输入，在 ScanQA、SQA3D 和 MSR3D 上超越需要 3D 结构信息的专门训练模型。

\paragraph{4-LEGS：四维语言-视觉定位} 4-LEGS（2024.10）~\cite{2410_10719} 将语言-视觉特征提升到基于 Dynamic 3DGS 的四维表示，实现从文本提示在视频中时空定位事件的交互式接口。4-LEGS 使用 ViCLIP-L-14 提取时空语言特征，通过 MLP 自编码器将 768 维特征压缩为 128 维潜特征嵌入每个高斯。训练时冻结预训练的动态 3DGS 参数，每个时间步优化 2000 次迭代（约 5 分钟/步），推理速度 10 FPS。其自注意力机制对每个高斯与其 20 个最近邻计算注意力权重，增强局部语义一致性。该方法将 LangSplat 的静态语义查询扩展到时序维度，实现从``人物''到``正在玩杂耍的人''的精细时空定位。

\paragraph{Semantic Consistent Language Gaussian Splatting：点级查询} Semantic Consistent Language Gaussian Splatting（2025.03）~\cite{2503_21767} 提出点级查询框架，通过对分割掩码进行跟踪建立语义一致的监督信号，并引入基于锚定的查询方法先检索蒸馏后的真值再查询个体高斯，在 LERF、3D-OVS 和 Replica 上分别取得 +4.14、+20.42 和 +1.7 的 mIoU 提升。与 LangSplat 的逐像素独立查询不同，其锚定查询机制先在蒸馏后的真值上匹配再索引到个体高斯，有效缓解了开放词汇查询中的语义漂移问题。

\paragraph{GVR：零样本视觉定位} GVR（Grounding via View Retrieval）（2025.09）~\cite{2509_15871} 提出零样本视觉定位框架，将三维视觉定位转化为二维检索任务，利用对象级视角检索从多个视角收集定位线索，避免三维标注成本和逐场景训练。与 LangSplat 需要逐场景优化语言场不同，GVR 的检索式范式无需任何三维端训练即可实现视觉定位。

\paragraph{PanopticQuery：四维动态消歧推理} PanopticQuery（2026.04）~\cite{2604_05638} 将查询时推理扩展到四维动态场景，基于 4D Gaussian Splatting 构建多视角语义共识机制，跨视角和时间帧聚合二维语义预测。该方法构建了 Panoptic-L4D 数据集，其核心设计理念为``可消解歧义''：场景中故意包含多个视觉上相似的对象，要求模型通过属性、动作、空间关系或交互四个维度进行消歧推理（如区分``睡觉的猫''和``走路的猫''），将 LangSplat 的静态语义分割能力推进到动态场景的细粒度关系推理。

\subsubsection{特征场蒸馏}

本路线介绍将二维基础模型的特征通过可微分渲染或其他映射机制融入三维高斯基元以实现超越外观的语义理解的方法。特征场蒸馏（feature field distillation）的目标是将二维基础模型（如 DINO、CLIP、SAM）提取的丰富语义和几何特征有效迁移到三维高斯空间，赋予其超越外观的语义理解能力。与前两节聚焦于开放词汇查询和语言场构建不同，本节更关注特征表示本身的效率、紧凑性和多模态融合。

%---------------------------------------------------------------------
\paragraph{语义特征蒸馏——从二维到三维的特征迁移}

语义特征蒸馏的核心是将 DINO、CLIP 等模型提取的二维语义特征高效地嵌入三维高斯表示。

Feature 3DGS（2023.12）~\cite{2312_03203} 奠定了这一方向的基础，通过可微分渲染将任意维度的语义特征蒸馏到三维高斯上，支持新视角语义分割和语言引导编辑。该方法揭示了高维特征与 RGB 图像之间的分辨率差异和通道不一致等根本性挑战。（详见本节前文对该方法的完整分析。）

\paragraph{FHGS：各向同性-各向异性协调} FHGS（特征均匀化高斯泼溅）（2025.05）~\cite{2505_19154}，全称~Feature-Homogenized Gaussian Splatting，从物理模型角度出发，指出~3DGS~光栅化管线的各向异性渲染机制与语义特征所需的各向同性表示之间存在根本矛盾：颜色渲染沿射线方向对不同高斯赋予不同权重（各向异性），而语义特征需要在相邻视角间保持一致（各向同性）。该框架包含三个核心组件。

（1）\textbf{通用特征融合架构}：通过空间哈希表加速三维点与二维特征的对应关系检索，从任意预训练分割模型生成二维特征真值，并为每个高斯基元初始化冻结的语义特征向量 $f_i$。

（2）\textbf{非可微特征融合机制（NDFD）}：语义特征 $f_i$ 在前向过程中不参与可微分渲染，而是直接利用贡献权重 $w_i$ 计算特征损失。梯度仅通过 $w_i$ 回传优化高斯的几何参数（位置、旋转、缩放、不透明度），隐式引导高斯基元向特征一致的区域移动。这样各向异性的颜色渲染专注于光照和阴影建模，而语义特征的各向同性通过权重驱动的分布优化实现，两者互不干扰。

（3）\textbf{物理启发的双驱动机制}：受物理场论中电场的直觉类比启发，将像素射线视为电场线，二维特征图中的真值特征视为场源，通过"场"的概念统一几何一致性和语义一致性的优化目标。

\paragraph{CUS-GS：紧凑统一结构化表示} CUS-GS（2025.11）~\cite{2511_17904} 提出紧凑统一结构化表示，首次将结构化 3DGS 的空间紧凑性与语义导向方法的多模态理解能力整合到统一框架中。该方法的设计动机是：现有方法要么追求结构效率但缺乏语义理解（如 Scaffold-GS、Octree-GS），要么引入语义能力但牺牲空间一致性和结构正则性（如 Feature 3DGS、LangSplat）。CUS-GS 的核心架构包含两个层次：体素化锚支架（voxel anchor scaffold）将三维空间划分为体素网格，每个体素 $v$ 维护一个可学习潜在特征向量 $z_v \in \mathbb{R}^{d_f}$，联合编码局部外观、几何和语义上下文，并通过共享轻量 MLP 解码出该体素内 $N$ 个高斯基元的全部属性（不透明度、协方差、颜色）；多模态特征记忆库（multimodal memory bank）从 CLIP、DINOv2、SEEM 等六个基础模型提取密集特征图，通过余弦相似度阈值 $\tau$ 进行冗余压缩，构建紧凑的多模态场景描述符集合。在语义检索阶段，CUS-GS 采用分层查询适应机制：先从体素级查询 $q_v$ 聚合粗粒度语义上下文，再通过高斯级偏移查询 $q_{v,n}$ 生成细粒度语义特征，确保空间连贯性和语义自适应性。此外，该方法引入特征感知的生长与剪枝策略，动态调整高斯密度以优化语义区域的表达。最终仅用 6M 参数实现与最优方法相当的竞争力——比最近的对手小一个数量级（35M）。

\paragraph{CF3：自顶向下紧凑特征场} CF3（Compact and Fast 3D Feature Fields）（2025.08）~\cite{2508_05254} 提出自顶向下的紧凑特征场构建流程，其核心观察是：为颜色细节优化而训练的大量高斯基元中，许多对表达特征场是冗余的。CF3 的流水线分为三个阶段。第一阶段是特征提升（feature lifting）：利用预训练的 3DGS 场景，对每个高斯基元计算其二维投影区域内多视角特征的加权平均——直觉上，如果一个高斯在多个视角中都可见，其三维特征应该是这些视角二维特征的加权组合。具体公式为 $f_i \approx \frac{\sum_{m=1}^{M} \sum_{p \in P_{i,m}} w_{i,m,p} \, F_{m,p}}{\sum_{m=1}^{M} \sum_{p \in P_{i,m}} w_{i,m,p}}$，其中 $w_{i,m,p}$ 是渲染权重。该方案无需联合训练 RGB 与特征，可快速获得与联合优化方法相当的特征质量。同时，CF3 通过计算提升特征的方差来识别几何不准确的异常高斯，过滤掉方差超过前 0.01\% 的基元。第二阶段是特征压缩：与 LangSplat 在特征提升前训练自编码器不同，CF3 先提升特征再训练逐高斯自编码器——自编码器架构为五层 MLP $[128, 64, 32, 16, 3]$，将 $D$ 维提升特征压缩到仅 3 维潜在空间，等价于将压缩特征视为三通道 RGB 颜色，可直接复用现有 3DGS 光栅化器，训练损失包含 MSE 重建损失、余弦相似度损失和轻量结构保持正则化项。第三阶段是自适应稀疏化：通过迭代剪枝（基于全局贡献度 $C(g_i)$ 排序）和合并（基于马氏距离 $d_M$ 识别重叠高斯对，用矩匹配方法计算合并后属性）将高斯基元数量大幅缩减，最终仅用 Feature-3DGS 5\% 的高斯构建竞争性的三维特征场。

\paragraph{Quantile Rendering：分位数渲染加速} Quantile Rendering（2025.12）~\cite{2512_20927} 提出分位数渲染策略，仅稀疏采样沿每条射线有主导影响的高斯而非密集采样所有交叉高斯，集成到泛化三维神经网络 GS-Net 中实现 512 维特征图的约 43.7 倍渲染加速。

\paragraph{4DLangVGGT：四维语言定位前馈} 4DLangVGGT（2025.12）~\cite{2512_05060} 提出首个基于 Transformer 的四维语言定位前馈框架，包含捕获动态场景时空几何表示的 StreamVGGT 和将几何感知特征投射到语言对齐语义空间的语义桥接解码器（SBD），可在多动态场景上联合训练并直接部署。

\paragraph{Feature4X：二维到四维特征蒸馏} Feature4X（2025.03）~\cite{2503_20776} 提出通用框架将任意二维视觉基础模型的功能扩展到四维领域，通过动态优化策略将多个模型能力统一到单一表示中，首次将视频基础模型（如 SAM2、InternVideo2）的特征蒸馏到显式四维特征场。

%---------------------------------------------------------------------
\paragraph{几何特征蒸馏与配准}

除了语义特征，几何特征（如深度、法线、表面属性）的蒸馏对于三维重建和配准任务同样重要。与 Feature 3DGS 和 LangSplat 专注于语言-视觉特征不同，本节的方法关注几何结构信息的提取和学习。

\paragraph{SceneSplat：原生三维高斯理解} SceneSplat（2025.03）~\cite{2503_18052} 提出首个原生在三维高斯上运行的大规模室内场景理解方法，配合自监督学习方案从未标注场景中解锁丰富的三维特征学习。在视觉-语言预训练中，SceneSplat 直接将高斯基元与 VLM 的图像嵌入空间对齐（而非降维到文本嵌入），保留更丰富的隐层语义信息；在自监督预训练中，采用多目标框架整合重建和自蒸馏对齐。并构建 SceneSplat-7K 数据集（7916 个场景）作为 3DGS 场景理解的首个大规模基准。该方法的重要意义在于表明 3DGS 参数本身即可作为有效的表示学习输入，无需额外的点云或网格转换。

\paragraph{GaussianCross：跨模态自监督} GaussianCross（2025.08）~\cite{2508_02172} 从跨模态输入出发，提出跨模态自监督三维表示学习架构，将尺度不一致的三维点云转换为统一的立方体归一化高斯表示，通过三属性自适应蒸馏 splatting 模块构建三维特征场，协同捕获外观、几何和语义线索。其立方体归一化将点云 $(x, y, z)$ 映射到 $[-1, 1]^3$，消除不同传感器尺度差异；三属性蒸馏 splatting 分别处理颜色、法线和语义三种属性。在 ScanNet200 上全量微调精度提高 9.3\% mIoU 和 6.1\% AP$_{50}$，证明了统一高斯表示对跨模态学习的有效性。

\paragraph{Learning Convex Decomposition via Feature Fields：几何驱动凸分解} Learning Convex Decomposition via Feature Fields（2026.03）~\cite{2603_09285} 通过学习特征场实现首个前馈式开放世界凸分解模型，学习连续特征场后通过自监督纯几何目标聚类产生凸分解，可泛化到网格、CAD 模型乃至高斯表示。与 LangSplat 的语义特征场不同，该方法的特征场完全由几何信号驱动，无需任何语义标注即可实现有意义的形状分解。

%---------------------------------------------------------------------
\paragraph{三维表示学习与预训练}

3DGS 作为场景表示正被探索用于三维视觉的预训练范式，利用其高效渲染和显式结构加速表示学习。与 Feature 3DGS 将 2D 特征蒸馏到高斯不同，本节的方法将 3DGS 本身作为表示学习的引擎或目标。

\paragraph{GS-PT：3DGS 集成点云自监督} GS-PT（2024.09）~\cite{2409_04963} 首次将 3DGS 集成到点云自监督学习中。基于 ShapeNet 的 55 类 50,000+ CAD 模型，为每个模型构建三元组 $T_i = (P_i, I_i, D_i)$：均匀采样 $N_p = 2{,}048$ 个点的点云、从四个正交视角渲染的 RGB 图像和深度图。3DGS 分支利用渲染图像提供跨模态对比学习的监督信号，点云分支使用层次化 Transformer 编码器处理输入。在 ModelNet40 上达到 93.8\% 分类精度（仅次于 Point-M2AE 的 94.0\%），在 ScanObjectNN 的三个变体上分别提升 11.9\%、8.84\% 和 8.89\%，在 5-way 20-shot 少样本学习上超越 Point-MAE +0.9\%，在三维目标分类、真实世界分类、少样本学习和分割等下游任务上超越现有自监督方法。

\paragraph{GS$^3$：渲染式点云预训练} GS$^3$（2024.11）~\cite{2411_18667} 提出基于 3DGS 的高效渲染式自监督点云预训练框架，通过比较渲染 RGB 图像和真实图像预训练点云编码器。其核心思路是用点云编码器预测的高斯参数（位置、协方差、颜色、不透明度）渲染 RGB 图像，以与真实图像的差异作为自监督信号。与 GS-PT 需要额外的 3DGS 分支不同，GS$^3$ 的点云编码器直接输出高斯参数用于渲染，实现约 9 倍预训练加速和不到 0.25 倍的内存开销，同时保持有竞争力的下游任务性能。

\paragraph{Learning Representations from 3D Gaussian Splats：高斯属性分类} GS-PT 和 GS$^3$ 都利用 3DGS 作为数据增强或自监督信号源，而 Learning Representations from 3D Gaussian Splats（2026.05）~\cite{2605_29549} 则直接将高斯属性（位置、缩放、旋转、不透明度、颜色 SH 系数）作为分类模型的输入特征。该研究对 PointNet、PointNet++、PointNeXt、DGCNN、GAT、SplineCNN 等多种架构在 ModelNet40、ShapeNet、ShapeSplat（30,843 样本、55 类）和 MACGS（4,288 样本、30 类）数据集上进行系统基准测试，揭示了架构族之间的系统性差异：PointNet 系列的自连接（self-connectivity）设计在高斯数据上表现稳健，SplineCNN 的局部卷积核机制最适应高斯的非均匀分布，而 DGCNN 的动态 k-NN 图在属性异质性下难以构建有意义的连接。高斯的额外属性（缩放、旋转、不透明度）相比纯三维坐标可提升分类精度，但贡献随架构而异。

\paragraph{3D Gaussian Point Encoders：显式逐点嵌入} 3D Gaussian Point Encoders（2025.11）~\cite{2511_04797} 引入基于三维高斯混合的显式逐点嵌入，通过自然梯度和 PointNet 蒸馏的优化技术学习高斯基，运行速度比同等精度 PointNet 快 2.7 倍，内存减少 46\%，FLOP 减少 88\%。

\paragraph{XSPLAIN：可解释性框架} XSPLAIN（2026.02）~\cite{2602_10239} 追求可解释性——提出首个专为 3DGS 分类设计的前置可解释性框架。仅使用几何属性 $g_i = (x_i, s_i, q_i, o_i) \in (\mathbb{R}^3, \mathbb{R}^3, \mathbb{R}^4, \mathbb{R}^1)$（排除颜色以提高鲁棒性），通过体素聚合 PointNet 骨干（网格尺寸 $G=7$，潜特征维度 $C=256$）和可逆正交变换 $U$ 解耦特征通道实现可解释性，并精确补偿分类器权重 $W'_{\text{cls}} = W_{\text{cls}} U^T$ 以保证变换不改变预测。原型纯度（prototype purity）从 13.0\% 提升至 22.5\%，用户研究（N=51）显示参与者 48.4\% 的时间选择该方法为最佳解释（$p < 0.0001$），删除测试中移除 top-1 体素导致精度下降 2.36\%（MACGS）和 6.82\%（Toys 数据集）。

\paragraph{Joint Semantic and Rendering Enhancements：Chebyshev 描述子} Joint Semantic and Rendering Enhancements（2026.01）~\cite{2601_02339} 提出三维语义高斯建模的联合增强框架，引入各向异性三维高斯 Chebyshev 描述子利用 Laplace-Beltrami 算子捕捉细粒度三维形状细节。该描述子将高斯的各向异性协方差 $\Sigma_i = R_i S_i S_i^T R_i^T$ 转化为频域表示，通过 Chebyshev 多项式在谱图上传播局部几何信息。同时，语义感知的高斯剪枝和 SH 调整模块根据局部梯度和语义线索自适应复制或移除高斯，并通过跨场景知识迁移模块维护模式基（pattern basis）持续更新学到的形状模式。

\paragraph{Geometry Meets Vision：几何接地实验指导} Geometry Meets Vision（2025.10）~\cite{2510_03104} 为上述所有蒸馏方法的特征选择提供了实验指导。系统研究几何接地（geometry-grounding）在蒸馏场中的作用，发现视觉-几何特征包含更精细的结构细节，但在语义对象定位上与纯视觉特征无显著差异，且位姿估计精度反而降低，揭示纯视觉特征在更广泛下游任务中的更高通用性。这一发现对 LangSplat 和 Feature 3DGS 等方法具有直接指导意义：若目标仅为语义分割或对象定位，引入几何信号可能是不必要的。

%---------------------------------------------------------------------
\paragraph{多模态特征融合与统一表示}

将来自不同基础模型的异构特征统一到三维高斯表示中，是实现全面场景理解的关键。Feature 3DGS 和 LangSplat 分别展示了 DINO 和 CLIP 特征蒸馏的范式，但它们各自只处理单一模态。本节的方法旨在融合多种特征源或设计更灵活的特征表示机制。

\paragraph{CLIP-GS（多模态）：序列化高斯标记} CLIP-GS（2024.12）~\cite{2412_19142} 提出基于 3DGS 的多模态表示学习框架，引入 GS Tokenizer 生成序列化高斯标记。具体而言，先用 FPS（最远点采样）和 kNN 将高斯组织为补丁，每个高斯点的属性（位置 $P \in \mathbb{R}^3$、SH 颜色 $C \in \mathbb{R}^{3(sh+1)^2}$、不透明度 $o \in \mathbb{R}^1$、缩放 $S \in \mathbb{R}^3$、旋转 $R \in \mathbb{R}^4$）被编码为标记序列，通过从点云模型初始化权重的 Transformer 层处理，利用与 CLIP 视觉-文本嵌入的对比损失（包括创新的图像投票损失 image voting loss）学习统一多模态表示。在多模态检索、零样本和少样本分类上超越纯点云模型，表明 3DGS 的丰富属性为多模态对齐提供了比纯坐标更有效的表示空间。

\paragraph{SpectralGaussians：跨光谱渲染} SpectralGaussians（2024.08）~\cite{2408_06975} 则将特征融合扩展到光谱维度，提出基于 3DGS 的跨光谱渲染框架。基于 Gaussian Shader 架构估计 BRDF 参数和光照，利用 Gaussian Grouping 按语义信息分组高斯，从配准的多视角光谱图和分割图生成语义有意义的光谱高斯。改进的物理渲染方法逐光谱估计反射率和光照，首次在 3DGS 框架中实现了材料识别、光谱分析和语义分割的端到端统一。

\paragraph{Feature Splatting：特征向量替代球谐} Feature Splatting（2024.05）~\cite{2405_15518} 转向高斯基元本身的表示能力。将高斯的颜色信息从球谐函数替换为逐高斯特征向量 $f_G \sim \mathcal{N}(0, 1)$（FeatSplat），通过两层 MLP（隐藏层 64 神经元，SiLU 激活）解码特征向量为 RGB 像素值，拼接相机位置 $x_{\text{cam}} \in \mathbb{R}^3$ 和像素嵌入 $e_p \in [-1,1]^2$ 条件化解码过程。32 维特征向量的表达力显著超越球谐函数（FeatsuSplat-32 在 Mip-360、T\&T 和 Deep Blending 上全面超越标准 3DGS），显著改善低重叠视角的新视角合成质量。同时，将 MLP 输出层扩展为 67 维（3 维 RGB + 64 维语义类别）即可支持逐像素语义标签生成，在 ScanNet++ 的 64 类封闭词汇分割任务上验证了特征向量的多功能性。

\paragraph{FHGS（续）：物理启发双驱动} FHGS（2025.05）~\cite{2505_19154} 的非可微特征融合机制和电势场启发的双驱动优化策略，为多模态语义特征的统一嵌入提供了物理启发的解决方案，解决了 Feature 3DGS 可微分渲染中多源特征梯度冲突的问题。

\paragraph{VolSegGS：动态体积分割} VolSegGS（2025.07）~\cite{2507_12667} 利用可变形三维高斯表示动态体积场景，通过高斯的视角无关颜色进行粗级分割并用亲和力场网络精炼，将分割结果嵌入高斯使其变形可连续跟踪分割区域。与 LangSplat 的静态语义场不同，VolSegGS 的分割结果随高斯变形在时间维度上连续传播。

\paragraph{LaGS：四维全景占据跟踪} LaGS（2026.02）~\cite{2602_23172} 提出用于四维全景占据跟踪的潜在高斯表示，受 COTR 架构启发，先融合观测到作为稀疏点中心潜在表示的三维高斯中，再将聚合特征 splat 到三维体素网格解码。其潜在高斯编码器在 3D 特征体量上细化特征，随后全景掩码解码器以查询为中心的方式解码语义和实例掩码，通过跟踪注意力机制维持跨帧身份。在 Occ3D nuScenes 和 Waymo 上取得四维全景占据跟踪的最优性能。

\paragraph{Explicit Memory through Online 3DGS：显式三维记忆} Explicit Memory through Online 3DGS（2025.10）~\cite{2510_23521} 利用在线三维高斯显式记忆增强类别无关视频分割，开发 FastSAM-Splat 和 SAM2-Splat 两种融合技术，用显式三维记忆改善基础模型的预测一致性和准确性。与 LangSplat 将语言嵌入高斯不同，该方法将分割结果本身作为三维记忆存储在在线构建的高斯场中，通过三维几何一致性纠正二维分割的时序不一致。

\begin{takeaway}
开放词汇理解与特征蒸馏将2D视觉语言模型的知识迁移到3DGS中，使3D场景理解不再受限于预定义类别。从 Feature 3DGS 的逐高斯特征嵌入到 LangSplat 的语言场，技术路线日趋成熟。VolSegGS 将语义分割扩展到动态场景，LaGS 实现四维全景占据跟踪，Explicit Memory 用三维记忆纠正二维分割的时序不一致。
\end{takeaway}

%=====================================================================
\section{目标检测、场景图与场景补全}
\label{sec:10_2}

3D 目标检测和 6D 位姿估计是从场景理解到具身智能的关键桥梁。场景图和语义场景补全则为更高层次的空间推理提供了结构化表示。本节从目标检测与位姿估计、场景图与语义建图、占用网格与语义场景补全三个角度展开讨论。

\subsection{检测与位姿估计}
\label{ss:10_4}

本节讨论2条技术路线：（1）3D目标检测；（2）6D位姿估计。

\subsubsection{3D目标检测}

本路线介绍利用 3DGS 高斯基元直接编码的 3D 空间信息进行目标检测的方法。3DGS 用于目标检测的核心优势在于：高斯基元直接编码了 3D 空间信息，无需从 2D 特征推断深度。然而，核心挑战也正源于此：3DGS 的训练目标仅是光度一致性，高斯的空间分布并非为检测而优化。

\paragraph{边界引导与聚焦采样的高斯检测} 3DGS-DET（2024.10）~\cite{2410_01647} 首次将 3DGS 引入室内 3D 目标检测，修改了 3DGS 的训练监督与后处理采样——引入 2D Boundary Guidance 增强空间分布、用 Box-Focused Sampling 替代随机下采样——解决了标准 3DGS 在检测任务中面临的两类结构性缺陷。

标准 3DGS 的训练目标仅为光度一致性（即渲染图像与真实图像的 $L_1$ + D-SSIM 损失），这导致高斯 blob 在 3D 空间中的分布呈现两个不利特征。第一，前景与背景的边界模糊：由于仅有像素级颜色监督，高斯 blob 会向物体边界两侧扩散，3D 空间中物体与背景的几何分界变得不清晰，对后续检测中的前景判别造成干扰。第二，背景 blob 冗余：2D 图像中大量像素属于背景区域，SfM 初始化会产生稠密的背景高斯，这些噪声 blob 在送入检测器前需要被随机下采样，而随机采样无差别地丢弃前景与背景 blob，导致检测所需的物体信息丢失。

\textbf{2D Boundary Guidance} 针对 blob 分布模糊问题，将 2D 边界先验注入 3DGS 训练过程。具体地，对每张输入图像，先用 Grounded SAM 生成类别级分割掩码，再通过 Suzuki-Abe 算法提取掩码边界并叠加类别颜色信息，生成边界标注图像 $I_b$。训练时将标准渲染损失的监督目标从原始图像 $I$ 替换为边界标注图像 $I_b$：

\begin{equation}
\label{eq:10_bg_loss}
\mathcal{L}_{\text{render}}^{\text{BG}} = (1-\lambda)\, L_1(\hat{I}, I_b) + \lambda\, L_{\text{D-SSIM}}(\hat{I}, I_b)
\end{equation}

直觉上，3DGS 为最小化该损失必须精确渲染边界线的颜色和位置，这迫使高斯 blob 在边界处凝聚并对齐，从而在 3D 空间中形成清晰的前景-背景分界。多视角一致性则由损失本身自然保证——同一 3D 边界从不同视角渲染时必须呈现一致的颜色和位置。

\textbf{Box-Focused Sampling} 针对背景 blob 冗余问题，用概率采样替代随机下采样。首先，通过 Grounding DINO 检测 2D 目标框，将每个框 $b_j$ 连同其最高类别概率 $p_j = \max_c p^c_j$ 投影到 3D 视锥体 $F_j$。对于落在视锥体内的每个高斯 blob $g_i$，其属于物体的概率为：

\begin{equation}
\label{eq:10_prob_agg}
p_{\text{agr}}(g_i) = \max_{v \in V}\; p(g_i \mid g_i \in F_j^v)
\end{equation}

其中跨视角取最大值用于聚合来自不同视锥体的先验，未落入任何视锥体的 blob 被赋予小概率 $p_{\text{bg}}=0.01$。最终依概率 $p_{\text{agr}}$ 对全部高斯进行采样：物体 blob 因概率高而更易被保留，背景 blob 因概率低而大量被丢弃，从而在固定采样预算内最大化物体信息密度。

在 ScanNet 上，3DGS-DET 比 NeRF-based 方法 NeRF-Det++ 在 mAP@0.25 上提升 +6.0，mAP@0.5 上提升 +7.8；在 ARKITScenes 上 mAP@0.25 提升 +14.9。

\paragraph{表面表示与高斯-体素融合} 与 3DGS-DET 通过 2D 先验改造 3DGS 训练管线不同，Gaussian-Det（2024.10）~\cite{2410_01404} 从几何先验角度出发：物体被连续表面包裹——皮肤包裹身体、皮革覆盖沙发、钣金包裹汽车——这一先验可用于推导对象性（objectness）。Gaussian-Det 首先将高斯基元重构为表面表示 $G^{\text{surf}} = \{(x_p, n_p, A_p)\}_{p=1}^P$，其中法向量 $n_p$ 通过求解协方差矩阵的特征值问题 $\Sigma_p n_p = \lambda n_p$ 获得（对应最小特征值的特征向量），横截面积 $A_p$ 取椭球最大截面。在此基础上，Closure Inferring Module（CIM）分两阶段精化检测提案。

\textbf{第一阶段——部分表面推断}：由于高斯光栅化仅受 2D 信号监督，局部表面特征推断本质上是欠定的，导致非结构化离群值（2023）~\cite{2023_sugar}。CIM 引入变分项 $F^V$ 建立概率特征残差 $F^{\text{cand}} \leftarrow F^{\text{cand}} + \lambda F^V$，其中 $F^V$ 从标准正态先验 $p_\theta(F^V) = \mathcal{N}(F^V; 0, I)$ 中采样，通过优化 ELBO 拟合后验 $q_\phi(F^V|F^{\text{cand}})$ 来显式建模这一不确定性。

\textbf{第二阶段——整体表面闭合聚合}：利用封闭曲面对常向量场的净通量为零这一几何定理（$\Phi = \oint_S T \cdot n \, dS = 0$），CIM 对每个提案 $p_k$ 内的部分表面数值计算通量 $\Phi_k = \sum_{g_i \in G^{\text{cand}}_k} T \cdot n_i A_i$。通量越接近零表明表面闭合度越高，提案质量越好，从而为对象性预测提供信息先验。

GVSynergy-Det（2025.12）~\cite{2512_23176} 采用与 Gaussian-Det 截然不同的双表示协同策略，洞察到连续高斯表示擅长精细表面细节而离散体素提供结构化空间上下文。其管线包含四个模块：2D-to-3D 特征提升构建体素表示 $V$，像素对齐 3DGS 模块预测深度并生成包含 64 维潜在特征的高斯基元，跨表示增强模块融合两种表示，最后由检测头输出 3D 包围盒。核心创新在于跨表示增强机制：先将高斯基元空间分箱为与检测体素网格对齐的规则体素 $V_g$ 并通过高斯特征编码器 $P_g$ 编码为 256 维特征，再与占据掩码 $O$ 相乘得到 $\hat{V}_g^o = \hat{V}_g \odot O$ 以标识有效几何区域；自适应权重网络 $W$ 基于拼接特征 $[V, \hat{V}_g^o]$ 动态计算体素权重 $\alpha_v$ 和高斯权重 $\alpha_g$，最终增强特征为 $V_e = F([\alpha_v V, \alpha_g \hat{V}_g^o])$。这种占据感知设计确保仅在几何真实存在处增强体素特征，在 ScanNetV2 和 ARKitScenes 上无需深度或密集 3D 监督即达到 SOTA。

\paragraph{雷达与自动驾驶场景} RadarGaussianDet3D（2025.09）~\cite{2509_16119} 面向 4D 汽车雷达的 3D 检测，与传统室内方法（如 3DGS-DET 针对多视角 RGB 输入）不同，它处理的是雷达点云这种天然稀疏且噪声更大的输入模态。现有方法依赖柱状编码器，每个点仅贡献单个 BEV 网格导致稀疏特征图。RadarGaussianDet3D 设计了 Point Gaussian Encoder（PGE）来弥补这一缺陷，PGE 包含四个步骤：（1）\textbf{局部特征聚合（LFA）}：对每个雷达点的球形邻域 $\mathcal{N}_i$ 应用轻量 PointNet 聚合位置偏移加特征，丰富单点表征；（2）\textbf{全局特征聚合（GFA）}：通过自注意力机制捕获全局上下文 $f_g = \text{FFN}(\text{Attn}(f_l)) + f_l$；（3）\textbf{高斯属性预测}：拼接原始、局部和全局特征后线性投影得到尺度 $s'$、旋转四元数 $q$ 和视角无关特征 $f^g$，其中不透明度固定为 1、位置偏移固定为 0 以避免早期训练中的自由度过大问题，尺度经 sigmoid 激活并限制上界 $s_\text{max}$；（4）\textbf{BEV 高斯 splatting}：将预测的高斯基元 $G = \{(\mu_i, s_i, q_i, o_i, f^g_i)\}_{i=1}^N$ 通过平行投影 splat 到 BEV 平面，用特征 $f^g$ 替代颜色渲染，生成密集 BEV 特征图 $F \in \mathbb{R}^{C \times H \times W}$。相比传统柱状编码器每个点仅影响单个网格，PGE 的每个高斯可自适应地影响多个相邻网格。同时提出 Box Gaussian Loss（BGL）：将真实和预测包围盒 $b = [x, y, z, l, w, h, \theta]$ 转化为 3D 高斯分布 $\mathcal{N}(\mu, \Sigma)$，其中 $\mu = [x,y,z]$，$\Sigma = R_b S_b S_b^T R_b^T$，缩放矩阵 $S_b = \frac{1}{a}\text{diag}(\frac{l}{2}, \frac{w}{2}, \frac{h}{2})$，$a > 0$ 为缩放超参数，通过反向 KL 散度 $D_{\text{KL}}(\hat{\mathcal{N}}, \mathcal{N})$ 统一优化包围盒的全部属性。RaGS（2025.07）~\cite{2507_19856} 是首个利用 3DGS 融合 4D 雷达和单目线索的方法：采用级联流水线——视锥定位初始化（FLI）将前景像素反投影为粗高斯中心，迭代多模态聚合（IMA）显式利用图像语义和隐式整合 4D 雷达速度几何精化高斯，多级高斯融合（MGF）将高斯场渲染为分层 BEV 特征进行检测，4D 雷达的速度维度为高斯空间分布提供运动锚定。

\paragraph{MVSDet：高斯正则化辅助检测} 与 3DGS-DET 修改训练管线不同，MVSDet（2024.10）~\cite{2410_21566} 仅将高斯光栅化作为辅助正则化信号：利用平面扫描进行几何感知检测，引入概率采样和软加权机制决定像素特征在 3D 体上的放置位置，并用像素对齐高斯光栅化正则化深度预测，以极小计算开销改善检测性能。

\paragraph{GFreeDet：高斯重建替代 CAD 模板} GFreeDet（2024.12）~\cite{2412_01552} 在 BOP Challenge 2024 中获得 model-free 2D 检测最佳方法奖和最佳快速方法奖，利用高斯光栅化从参考视频直接重建物体替代 CAD 模板——不进行检测，而是用高斯重建替代 CAD 模板以支撑下游检测器，在 BOP-H3 基准上达到与 CAD-based 方法可比的性能。

\paragraph{RCGDet3D：增强雷达特征编码} RCGDet3D（2026.05）~\cite{2605_21112} 重新思考 4D 雷达-相机融合检测，聚焦于增强雷达特征编码而非复杂融合策略——与 RadarGaussianDet3D 用高斯替代柱状编码器的思路互补，RCGDet3D 保持传统编码架构但强化特征表示，在保持推理速度的同时提升检测精度。

\paragraph{NeRF-based 预训练视角：理论分析} NeRF-based 预训练视角（2026.03）~\cite{2603_28090} 探讨了是否需要视角变换来将 NeRF 预训练应用于 3D 感知模型，为基于 NeRF/3DGS 的检测预训练范式提供了理论分析。

\subsubsection{6D位姿估计}

本路线介绍利用 3DGS 可微渲染构建位姿优化管线以实现 6D 位姿估计和跟踪的方法。

6D 位姿估计（3D 平移 + 3D 旋转）是机器人抓取和增强现实的核心任务。3DGS 在此任务中的独特价值在于：高斯基元的可微渲染可用于构建位姿可微的渲染管线，通过比较渲染图像和观测图像迭代优化位姿。这一思路与经典的 Bundle Adjustment 一脉相承，但 3DGS 的显式表示使渲染速度比 NeRF 快数个数量级。

\paragraph{高斯光栅化驱动的在线位姿估计} 6DOPE-GS（2024.12）~\cite{2412_01543} 是首个利用高斯光栅化进行在线 6D 位姿估计和跟踪的方法。其核心思路是同时优化 6D 对象位姿和 3D 对象重建——利用 2DGS 的快速可微渲染能力，通过渲染损失同时优化两者。与基于 NeRF 的 BundleSDF（平均跟踪频率仅 $0.4$ Hz）相比，6DOPE-GS 利用高斯光栅化替代体渲染，从根本上降低了计算开销。具体地，给定当前位姿估计 $\mathbf{T}_t$ 和高斯模型 $\mathcal{G}$，渲染图像 $\hat{I}(\mathbf{T}_t, \mathcal{G})$ 与观测图像 $I_t$ 之间的光度差异驱动位姿更新：

\begin{equation}
\label{eq:10_pose_opt}
\mathbf{T}_{t+1} = \mathbf{T}_t - \eta \frac{\partial \mathcal{L}_{\text{photo}}(\hat{I}(\mathbf{T}_t, \mathcal{G}), I_t)}{\partial \mathbf{T}_t}
\end{equation}

直觉上，式~\ref{eq:10_pose_opt}~的含义是：如果渲染图像与观测图像在某处不一致，梯度会告诉位姿应该向哪个方向调整才能减少差异。与传统 ICP 依赖点云配准、PnP 依赖 2D-3D 对应关系不同，这一基于可微渲染的范式直接在像素空间建立损失，端到端地联合优化位姿和场景表示——无需显式提取特征对应关系，避免了特征匹配失败时的级联误差。

6DOPE-GS 的完整管线分为三个阶段。

\textbf{粗位姿初始化}阶段利用 SAM2 进行对象分割，LoFTR 建立帧间点对应关系，通过 RANSAC 式最小二乘求解粗位姿并构建关键帧记忆池。

\textbf{联合优化}阶段是方法的核心：采用增量式 2D 高斯光栅化，将每个高斯基元投影到关键帧位姿下进行 $\alpha$-混合 渲染，渲染损失（包含颜色一致性 $\mathcal{L}_c$、深度一致性 $\mathcal{L}_d$、深度失真 $\mathcal{L}_{dd}$ 和法线对齐 $\mathcal{L}_n$ 四项）同时反传至高斯参数和以 $\mathrm{SE}(3)$ 参数化的位姿变换 $\mathbf{T}_k$，实现位姿与对象模型的联合精化。这一增量机制的关键优势在于：2DGS 将每个高斯建模为定向平面盘（surfel），相比 3DGS 的体积表示更适合精确的表面几何重建，渲染深度和法线估计更可靠。

为确保联合优化的稳定性和效率，6DOPE-GS 提出两项关键机制。

\textbf{动态关键帧选择}在以二十面体顶点和面心构建的锚点集上聚类关键帧位姿，从每个锚点簇中选择对象掩码最大的关键帧参与优化——这既确保了空间视角的全覆盖（避免视角退化导致重建空洞），又限制了参与优化的关键帧数量以控制光栅化的线性计算开销。优化过程中，该方法还基于重建损失的中位数绝对偏差（MAD）实时移除异常关键帧，防止错误位姿估计污染高斯训练。

\textbf{不透明度百分位数修剪}替代了原始 3DGS 中需要大量工程调参的自适应密度控制：每 100 步移除不透明度处于底部 5\% 的高斯基元，直至第 95 百分位不透明度超过阈值；配合基于位置梯度的分裂和克隆操作，在训练稳定性与模型紧凑性之间取得平衡。优化收敛后冻结高斯属性，仅精化位姿 500 步。

\textbf{在线位姿图优化}阶段将精化后的关键帧位姿注入全局位姿图，对新到达帧基于视锥可见性选择重叠帧构建局部位姿图，通过成对几何一致性最小化像素级重投影误差，实现逐帧连续跟踪。整个系统以双线程并行运行：在线跟踪线程约 $4.1$ Hz，高斯对象场线程平均每关键帧 $0.23$ 秒。在 HO3D 和 YCBInEOAT 数据集上，6DOPE-GS 在匹配 SOTA 性能的同时提供 $5\times$ 加速。

\paragraph{位姿可微渲染管线} 与 6DOPE-GS 同时优化位姿和高斯模型不同，GS2Pose（2024.11）~\cite{2411_03807} 采用两阶段解耦框架。粗估计阶段设计带极化注意力机制的 U-Net（Pose-Net），在 3DGS 模型从多视角生成的 RGB-NOCS 图像对监督下训练，从 RGB 图像预测 NOCS 图像并经改进 ICP 对齐计算粗位姿 $T_{cm}^{\text{coarse}}$。精化阶段（GS-Refiner）借鉴 Bundle Adjustment 思想：利用李代数 $\mathfrak{se}(3)$ 参数化位姿变换，将 $\mathrm{SE}(3)$ 上的非线性优化转化为无约束优化——具体地，将位姿 $\mathbf{T} = [\mathbf{R} | \mathbf{t}]$ 映射为李代数向量 $\xi \in \mathbb{R}^6$（3 个旋转分量 + 3 个平移分量），在李代数空间中通过 3DGS 可微渲染管线计算光度损失 $\mathcal{L}_{\text{photo}}(\hat{I}(\xi), I_{\text{in}})$ 对 $\xi$ 的梯度并迭代更新。这一李代数扩展确保优化始终在 $\mathrm{SE}(3)$ 流形上进行，避免旋转矩阵正交性约束被破坏。GS-Refiner 还选择性更新 3DGS 模型的球谐颜色参数以适应测试时的光照变化，在 T-LESS、LineMod-Occlusion 和 LineMod 上分别提升 1.4\%、2.8\% 和 2.5\% 精度。GS2POSE（2025.10）~\cite{2510_16777}（同名不同团队）采用类似的粗-精两阶段思路，但改进了颜色更新策略：通过更鲁棒的颜色参数自适应机制增强对遮挡和光照变化的鲁棒性，在遮挡场景下的位姿精度优于 GS2Pose。

\paragraph{单图像与闭合形式位姿估计} 6DGS（2024.07）~\cite{2407_15484} 从完全不同的角度解决位姿估计——不是优化对象位姿，而是从单个 RGB 图像和 3DGS 场景模型估计\textbf{相机位姿}。核心创新是定义 Ellicell（Ellipsoid Cell）：将每个 3D 高斯椭球体的表面等面积划分为若干单元，在每个单元中心沿从椭球中心到单元中心的方向生成射线，构建射线集 $V = \{\langle v_o, v_d, v_c \rangle\}_{i=1}^N$，其中 $v_o \in \mathbb{R}^3$ 为椭球中心（射线原点），$v_d \in \mathbb{R}^3$ 为从中心到单元中心的方向，$v_c \in \mathbb{R}^3$ 为通过 3DGS 渲染函数 $\mathcal{R}$ 合成的射线颜色。Ellicell 的划分是确定性的（基于等面积球面采样），比 Monte Carlo 随机采样以更少射线达到更高精度。射线-像素绑定通过注意力图 $A$ 实现：计算目标图像像素特征 $\phi(I_t)$ 和射线特征 $\phi(V)$ 之间的注意力权重，注意力值越高表示该射线-像素对越可能对应有效交点。最终通过加权最小二乘求最优射线束的交集直接计算相机中心，旋转则由射线方向与像素方向的对应关系确定。这一方法以闭合形式求解 6DoF 位姿，无需初始化和迭代优化，在真实场景上旋转精度提升 12\%、平移精度提升 22\%，同时达到近实时约 16 FPS（比 iNeRF 快约 94$\times$）。SplatPose（2025.03）~\cite{2503_05174} 提出 Dual-Attention Ray Scoring Network（DARS-Net），通过几何域注意力机制解耦位置对齐和角度对齐，显式建模方向依赖性以缓解旋转歧义，配合粗到细优化管线通过密集 2D 特征对齐精化位姿估计。

\paragraph{SGPose：稀疏视角位姿估计} 与 6DOPE-GS 依赖密集关键帧和在线高斯重建不同，以下方法在更约束的条件下完成位姿估计。SGPose（2024.09）~\cite{2409_02581} 解决稀疏视角（仅约 10 个视角）下的位姿估计问题：不依赖 SfM 管线，而是从随机立方体初始化开始，将高斯基元建模为 2D 椭圆盘（而非 3D 椭球），在其局部切平面上参数化协方差矩阵。通过几何感知深度渲染——利用射线的齐次坐标变换推导精确的像素-高斯深度对应关系——在极稀疏视角下提供可靠的深度监督信号。同时，通过在线合成视角扭曲生成额外训练视角，SGPose 在 Occlusion LM-O 等基准上在稀疏视角约束下仍超越现有方法。

\paragraph{UnPose：零样本扩散位姿估计} UnPose（2025.08）~\cite{2508_15972} 提出零样本、无 CAD 模型的 6D 位姿估计框架，与 6DOPE-GS 的核心差异在于：6DOPE-GS 从关键帧逐步构建高斯对象场，而 UnPose 从单视角 RGB-D 帧出发，利用 Wonder3D 扩散模型合成 6 个新视角图像并建立初始 3DGS 对象场 $Q_o$。关键技术是对扩散模型进行不确定性估计：通过末层拉普拉斯近似（LLLA）对预训练扩散模型的噪声预测进行贝叶斯推断，在无需额外训练的情况下获得像素级认知不确定性 $\text{Var}(x_0)$。该不确定性贯穿整个管线——初始化阶段调制扩散视角的置信度以引导位姿图优化，3DGS 映射阶段将其作为损失权重 $\mathcal{L}_t = \sum_p \text{Var}(p) \cdot \mathbb{1}(\text{vis}(p) > 0.99) \cdot (L_1(D(p)) + 0.5 L_1(C(p)))$，确保高不确定性区域（如扩散幻觉产生的伪影）不污染高斯场更新。后端优化通过 Mast3R 匹配帧间对应关系构建位姿图并闭环，实现增量精化。

\paragraph{GSGTrack：单目 RGB 位姿跟踪} GSGTrack（2024.12）~\cite{2412_02267} 基于 RGB 视频进行 6DoF 位姿跟踪（无需深度），与 6DOPE-GS 的差异在于：6DOPE-GS 用 2DGS 并依赖 RGB-D 输入，而 GSGTrack 仅用单目 RGB 并引入广义立体匹配网络 DUSt3R 提供粗几何先验。GSGTrack 采用在线 3DGS 表示与图基几何优化联合框架。核心创新包括：（1）\textbf{距离变换轮廓损失}：替代传统 IoU 轮廓损失，通过计算每个像素到轮廓边缘的欧氏距离加权损失 $L_s = \frac{1}{|\Omega|} \sum_p (D_S(p) \cdot (1-\hat{S}(p)) + D_{\hat{S}}(p) \cdot (1-S(p)))$，在重叠极小时仍提供有效梯度；（2）\textbf{在线几何结构图}：有向图中每个节点存储帧的 6DoF 位姿和 3D 表示，每条边存储 DUSt3R 的像素对齐点图和置信度图，避免重复计算且抑制跟踪漂移；（3）\textbf{图像对剪枝策略}：基于位姿一致性（PnP-RANSAC 估计与跟踪位姿比较）、几何相似性和深度置信度三重准则过滤低质量匹配对，确保鲁棒优化。

\paragraph{FloPE：机器人授粉位姿估计} 上述方法聚焦于通用检测和位姿估计，而以下工作将 3DGS 的位姿相关能力拓展到特定应用场景。FloPE（2025.03）~\cite{2503_11692} 将 6D 位姿估计应用于机器人授粉——利用 3DGS 生成带精确位姿标注的逼真合成数据集，通过知识蒸馏将高容量教师模型的知识迁移到轻量学生模型，在计算受限的机器人平台上实现 0.6cm 平均位姿误差和 78.75\% 授粉成功率。与通用位姿方法不同，FloPE 的核心贡献在于将 3DGS 的逼真渲染能力转化为数据增强工具，弥合了实验室环境与田间部署之间的域差距。

\paragraph{光场位姿跟踪：多方向光线表示} 光场位姿跟踪（2025.12）~\cite{2512_13007} 利用光场图像的多视角特性提取语义和几何特征，转化为视角依赖高斯 splats 作为统一对象表示。与标准 3DGS 位姿方法依赖 RGB 或 RGB-D 输入不同，光场输入天然编码了多方向的光线信息，使高斯 splats 能够捕获反射、折射等复杂视觉行为，在包含此类挑战性外观的物体上与基于精确 CAD 模型的跟踪器竞争。

\paragraph{LineGS：3D 线段位姿表示} LineGS（2024.12）~\cite{2412_00477} 将 3DGS 的应用范围从面扩展到线——利用边缘处高密度 Gaussian 点分布的天然优势（边缘处高斯密度高且排列有序），精化和优化传统几何方法生成的初始 3D 线段。与面基方法关注区域覆盖不同，LineGS 关注一维几何基元的精确重建，在几何精度和模型紧凑性上显著提升，为线特征驱动的位姿估计提供了新的 3D 表示路径。

场景图将 3D 场景组织为结构化的语义表示：节点代表对象实例，边编码对象之间的空间和语义关系。这种结构化表示为机器人任务规划和交互式导航提供了推理基础。

\paragraph{自适应聚类与场景图生成}

GaussianGraph（2025.03）~\cite{2503_04034} 是将~3DGS~与场景图深度融合的代表性工作，其核心洞察是：现有~3DGS~语义方法存在两个关键缺陷——(1) 将高维~CLIP~特征压缩后附加到高斯基元上会损失细粒度语义信息，导致分割精度下降；即使~OpenGaussian~等替代方案尝试通过~K-Means~将未压缩特征匹配到高斯簇，但需要手动指定簇数且随机初始化往往导致次优的~2D-3D~关联；(2) 现有方法仅关注~CLIP~特征学习，缺乏对象级空间关系，因此无法处理涉及空间推理和歧义类别的复杂查询。

为解决第一个问题，GaussianGraph~提出了"控制-跟随"（Control-Follow）聚类策略，其直觉是：与其对所有~$M$~个高斯点进行聚类（计算开销大且受特征噪声影响），不如先选取语义特征稳定的稀疏"控制点"作为聚类种子，然后让其余高斯点"跟随"分配到最近的簇。具体而言，在控制阶段（Control Stage），系统根据实例特征在连续迭代帧间的变化量筛选稳定点（要求~$\frac{1}{k_2 - k_1}\sum_{\text{iter}=k_1+1}^{k_2}\|\nabla I_n^{\text{iter}}\| < \epsilon$），然后从中下采样得到~$m$（$m \ll M$）个控制点。以控制点的实例特征和~3D~坐标为联合特征，采用~BIRCH~算法获得初始聚类中心~$E_c = \{e_1, e_2, ..., e_{k_c}\}$。在跟随阶段（Follow Stage），所有高斯点根据特征~$G^f = [I_n, XYZ]$~被分配到最近的初始簇；若最小距离超过阈值，则形成新簇。最后，对包含多个邻近对象的单一簇进行基于实例特征的细化拆分，确保每个簇对应一个语义一致的~3D~对象。这一策略的关键优势在于：聚类数量由场景尺度和特征分布自动确定，无需人工指定。

在场景图构建阶段，GaussianGraph~首先利用~2D~基础模型（SAM~生成全分割掩码，LLaVA~生成对象类别和密集描述，Grounding DINO~定位前景对象）提取对象属性；然后通过~VLM~提取对象间的功能性和空间性关系三元组~$<\text{subject}|\text{predicate}|\text{object}>$。为过滤~2D~推理引入的不可靠关系，系统引入了四个~3D~校正子模块：（1）接触验证——对于"in"或"on"等需要直接接触的空间关系，将两个对象的~3D~点集投影到地平面并计算凸包交集，仅保留~$\text{Convex}(P_i^{\text{proj}}) \cap \text{Convex}(P_j^{\text{proj}}) \neq \emptyset$~的关系；（2）方向验证——通过计算对象中心间向量与坐标轴的点积来验证"in front of"或"behind"等方向关系；（3）距离消歧——对于同类别多对象的定位查询（如"最靠近黑板的椅子"），利用~3D~高斯簇间的欧氏距离选择目标；（4）邻近距离阈值——对于"next to"等邻接关系，若~3D~中心距离超过场景尺度的十分之一则丢弃该关系。实验表明，GaussianGraph~在~LERF、Replica~和~ScanNet~三个数据集上的语义分割和~3D~对象定位（Sr3D+/Nr3D）任务中均超越此前最优方法。

\paragraph{动态场景图的增量更新}

GaussianGraph~的场景图是静态的——一旦构建完成便无法适应环境变化。然而在真实的人机协作场景中，对象的移动、消失和新出现是常态。DynamicGSG（2025.02）~\cite{2502_15309} 正是为了解决这一瓶颈，将~3DGS~场景图从"一次性构建"推进到"持续自适应"的范式。

系统的核心创新在于将高斯建图与实例分组联合优化。具体而言，每个高斯基元除了标准参数（颜色、位置、半径、不透明度）外，还附加了低维实例特征~$e \in \mathbb{R}^3$~和实例标识符~$idx$。在渲染时，实例特征通过与颜色相同的~$\alpha$-混合方式投射到~2D~平面，生成像素级实例特征图~$\hat{E}_t$。为防止优化过程中高斯点跨越对象边界（导致分割伪影和边界模糊），DynamicGSG~设计了特征一致性损失来监督实例分组：
$$\mathcal{L}_f = \lambda_1|\hat{E}_t - E_t| + \lambda_2(1 - \text{SSIM}(\hat{E}_t, E_t))$$
其中~$E_t$~为从~200~色实例调色板映射的~2D~真值实例特征图。最终的联合优化损失整合了颜色、深度和特征三路监督：$\mathcal{L} = w_c\mathcal{L}_c + w_d\mathcal{L}_d + w_f\mathcal{L}_f$，在优化高斯建图质量的同时确保实例级分组的准确性。

在动态更新方面，系统首先通过~VIO~获取初始位姿估计，再利用渲染的彩色图和深度图与观测之间的~L1~损失进行迭代位姿精化。然后检测环境变化：对于~50\%~以上高斯点落入当前视锥的对象，渲染其~RGB~掩码并与观测图像的对应区域计算~SSIM；若~$S(j) < \theta_{\text{change}}$，则判定该对象已被移动或替换，将其归入消失对象集合并移除关联的过期高斯点。新出现的对象则通过~YOLO-World~检测和~SAM~分割进行识别、初始化并注册到场景图中。关键设计是：所有更新均在实例级局部执行，仅修改变化区域的~Gaussian~场景图和层次化图结构（包含资产对象、普通对象和独立对象三层），而非重建整个场景。实验在~Replica、ScanNet++~数据集和真实实验室环境中验证了系统的动态更新能力、高保真重建质量和语言引导的对象检索效果。

\paragraph{光学-几何-语义统一建图}

如果说~DynamicGSG~聚焦于场景图的时间维度（动态更新），OmniMap（2025.09）~\cite{2509_07500} 则在空间模态维度上实现了突破——将光学（高保真渲染）、几何（精细网格重建）和语义（开放词汇理解）三种能力统一到一个在线建图框架中。

OmniMap~采用紧耦合的~3DGS-Voxel~混合表示架构。直觉是：3DGS~擅长精细的外观建模和新视角渲染，但高斯点的空间分布不均匀且缺乏结构化约束；体素网格提供规则的空间结构和概率建模能力，但分辨率受限于体素大小。两者互补：体素通过~TSDF-Fusion~提供规则化的空间覆盖，3DGS~通过可微渲染提供高保真外观。具体而言，系统并行维护一个~3DGS~场~$\mathcal{G}_t$（高斯基元集合）和一个开放词汇体素图~$\mathcal{M}_t$（概率实例体素和嵌入码本），体素图存储实例级语义理解，3DGS~场提供精细渲染和几何。

实现层面，OmniMap~的三项核心创新如下。（1）自适应相机建模：在图像多视角稠密建图中，OmniMap~对相机内参和畸变参数进行可微建模，以补偿运动模糊和曝光变化对重建质量的影响，避免了因相机参数不准导致的重建退化。（2）基于体素的高斯初始化与法线监督：不同于暴力回投影每一帧关键帧（导致观察频繁区域高斯密度过高），OmniMap~利用~TSDF-Fusion~的体素计数来识别新增体素~$V_t^+ = \{v^j | c_{t-1}^j = 0 \text{ and } c_t^j > 0\}$，仅从新增体素中心初始化高斯点，实现规则且高效的增量扩展。同时引入法线方向的几何信号监督，约束高斯基元的表面朝向一致，提升几何重建精度。（3）概率融合的实例级理解：OmniMap~首先通过~YOLO-World~检测开放集实例，再用~TAP~模型生成像素级分割掩码和描述性标题；随后利用~LLM（SBERT）编码标题为语义嵌入。在~3D~层面，采用概率建图框架——每个体素存储概率实例~ID~元组，通过链式法则将问题分解为实例关联任务~$P(I_t | \mathcal{M}_{t-1}, Q_t)$~和活图演化任务~$P(\mathcal{M}_t | \mathcal{M}_{t-1}, Q_t, I_t)$，以增强前端分割噪声下的建图鲁棒性。实验表明，OmniMap~在~Replica~和~ScanNet~数据集上的渲染保真度（PSNR/SSIM/LPIPS）、几何精度（F-score）和开放词汇语义分割（mIoU）上均超越此前方法，并支持多域场景问答、交互式编辑、机械臂操控和移动导航等下游应用。

\paragraph{大规模三维空间智能数据管线}

上述方法展示了~3DGS~场景理解的算法进步，但其训练和评估受限于~3D~空间数据的稀缺性——现有数据集要么依赖昂贵的专业~3D~扫描设备，要么局限于~50~个封闭类别。Holi-Spatial（2026.03）~\cite{2603_07660} 不直接提出新的场景理解算法，而是构建了首个全自动化大规模~3D~空间智能数据管线和数据集，为空间智能模型的训练提供了可扩展的数据基础。

其数据管线包含三个阶段。（1）几何优化：从原始视频流出发，首先利用~SfM~解算相机参数，再结合空间基础模型生成稠密点云，最后通过~3DGS~逐场景优化（集成多视角深度一致性正则化和表面重建约束）消除大尺度飘浮物，获得与物理表面对齐的干净~3DGS~重建。（2）图像级感知：均匀采样关键帧后，利用~Gemini-3-Pro~生成场景描述，并通过动态类标签记忆~$M_t$~确保跨帧语义一致性（每帧提示~VLM~优先复用已有标签，形式化为~$M_t = M_{t-1} \cup \phi(I_t)$）。SAM3~执行开放词汇实例分割后，利用优化后的~3DGS~渲染深度图将~2D~掩码反投影为~3D~有向包围盒（OBB），并引入几何感知滤波策略（掩码腐蚀 + 网格引导的深度滤波）抑制深度边缘飘浮物。（3）场景级细化：采用由粗到精策略——先通过空间聚类合并同类别且~3D~重叠度超过~$\theta_{\text{merge}} = 0.2$~的冗余检测，保留最高置信度的观测视角；再进行全局重力对齐修正不一致的旋转偏差；最后由~VLM~生成实例级详细标题。

在此基础上构建的~Holi-Spatial-4M~数据集包含：12K~个优化后的~3DGS~场景（源自~ScanNet、ScanNet++~和~DL3DV-10K），1.3M~个~2D~实例掩码，320K~个~3D~边界框，320K~条实例标题，1.2M~个~3D~定位对，以及~1.25M~个空间问答对。问答对分为以摄像机为中心的任务（如旋转和运动方向）和以对象为中心的任务（如对象间距离、全局/局部方向、尺寸测量），确保几何推理、关系推理和语义推理的均衡覆盖。利用该数据集微调视觉语言模型，在空间推理和~3D~定位任务上取得了显著提升。

\paragraph{语义特征驱动的跨域全局对齐}

前面讨论的方法都假设输入数据具有足够的视角重叠。但在实际场景中，来自无约束网络图片的局部重建往往几乎没有重叠——例如游客拍摄米兰大教堂时，绝大多数照片集中在正门，少量覆盖后墙，视角偏差导致传统~SfM~管线生成多个断裂的部分重建或错误合并。

Scene Grounding in the Wild（2026.03）~\cite{2603_26584} 的核心洞察是：尽管真实照片与低质量~3D~模型在外观上差异巨大，但底层语义是共享的——同一栋建筑的语义特征无论拍摄条件如何都保持一致。基于此，该工作提出了一条利用语义特征进行跨域对齐的技术路线：首先利用~Google Earth Studio~渲染密集的伪合成视角作为参考模型，通过~COLMAP~构建全局一致的~3DGS~表示，并在每个高斯点上蒸馏~DINOv2~语义特征向量（使每个高斯同时拥有颜色和特征向量）。然后将真实世界图像集视为一个"元图像"（meta-image）$\bar{I}$，将对其的全局~6DoF+尺度变换~$T$~的对齐问题建模为逆优化问题——冻结~3DGS~模型参数，引入~7~参数向量（6~个刚体变换参数~$SE(3)$~和~1~个尺度参数），以语义~L1~损失~$L_{\text{sem}}$~为目标函数进行迭代优化。

关键挑战在于异常值（outliers）：参考模型中的飘浮物导致渲染视角被遮挡，真实图片中可能包含瞬态遮挡，这些异常值会严重干扰优化收敛。为此，系统采用最小修剪平方（Least Trimmed Squares, LTS）鲁棒优化策略——每次迭代忽略~$L_{\text{sem}}$~值超过前一次迭代中位数损失的图像，逐步剔除异常值的干扰。实验在自建的基准数据集上进行评估（包含来自~WikiScenes~的局部重建和~Google Earth Studio~的参考模型，并提供精确的对齐真值），在平移误差、旋转误差和元图像变换准确率（MTA）上均超越~COLMAP、gDLS+++~和~SP+LG~等基线方法，以及~DUSt3R、MASt3R、VGGT~等前馈~3D~模型。

\paragraph{语义控制的大规模户外场景重建}

上述方法主要面向室内场景和"圆形"拍摄模式，而大规模户外场景提出了不同的挑战：摄像机沿路径前进而非环绕对象旋转，每处位置仅有少量帧，且需要移除的不再是小物体而是大型场景元素（如建筑物、道路、水面）。SCGS（2024.09）~\cite{2409_15959}（Semantics-Controlled GS）正是为此场景设计的语义驱动重建方法。

SCGS~的核心思想是：与其像传统方法那样在高斯基元上附加可学习特征向量再训练分类器，不如直接为每个高斯基元分配一个确定性的语义类别~ID，将分割问题重构为高斯参数空间中的直接分类任务。具体实现上，SCGS~修改了高斯光栅化过程——在标准~RGB~渲染管线之外并行运行一条语义渲染管线，将球谐系数置零以隔离语义属性与颜色信息，然后通过~$\alpha$-混合将语义类别~ID~投射到~2D~图像平面：
$$S = \sum_{c} C T_i \alpha_i s_{c_i}$$
语义分割通过交叉熵损失在~2D~平面进行监督，梯度通过修改后的光栅化过程反向传播到~3D~高斯基元。这一设计的优势在于：分配的类别~ID~允许直接、大规模地选择性移除或修改某一类~3D~高斯点，无需额外的凸包后处理或特征相似度计算：
$$G_{\text{new}}(x; \theta) = G(x; \theta) \cdot \mathbb{I}(s_{i_c} \neq c_{r=1..n})$$
其中~$\mathbb{I}$~为指示函数，仅保留不属于被移除类别的高斯点。

为验证方法，SCGS~构建了大规模户外~3D~资产数据集，使用~Insta360 X1/X2/X3~全景相机沿路径拍摄，通过组合前视、~$60°/30°$~侧视和~$10°$~仰俯视图像实现场景的全面覆盖，使用~DPT~获取~ADEK20~标签的分割掩码。实验表明，SCGS~在自建户外数据集上超越~Gaussian Grouping~等基线方法的渲染质量和分割精度，能够移除大型场景元素而不影响剩余部分。用户研究进一步表明，在~VR~环境中自由移动时，参与者明显偏好~SCGS~的交互体验，认为其在临场感和感知质量上显著优于普通~3DGS~场景。

占用网格和语义场景补全是自动驾驶和机器人感知中的核心任务，3DGS 为其带来了稀疏表示的新范式。

\paragraph{从密集体素到稀疏高斯：表示范式的转变}

传统~SSC~方法依赖密集体素网格，大量计算浪费在空白空间上。近期的~Gaussian-based~方法用~3D~高斯作为中间表示，通过高斯到体素的泼溅（splatting）操作生成最终的占用预测，在计算效率和表示能力之间取得更好的平衡。

VG3S（2026.03）~\cite{2603_06210}（Visual Geometry Grounded Gaussian Splatting）从特征提取层面解决上述瓶颈：现有高斯基方法的图像编码器仅靠有限的占用监督训练，无法习得强~3D~几何先验，导致跨视角结构不一致——表现为不完整的可行驶平面和人造结构碎片化。VG3S~引入预训练的几何感知~VFM（如~DVGT）提取视觉词元，并通过可插拔的层次化几何特征适配器（Hierarchical Geometric Feature Adapter, HGFA）将其转化为几何增强的表示。HGFA~由三个顺序模块组成：（1）分组自适应词元融合（GATF），按层级聚合~VFM~的多尺度嵌入，保留不同抽象层次的空间信息；（2）任务对齐词元精炼（TATR），以占用预测目标为导向校准通用~VFM~特征，使其从通用视觉语义转向面向~3D~几何的专用表示；（3）潜在空间特征金字塔（LSFP），对校准后的特征进行多尺度空间重构，输出与高斯基元编码器兼容的几何增强词元。这些词元随后通过~GaussianFormer~的编码器-解码器流水线生成高斯基元并泼溅为体素占用预测。在~nuScenes~基准上，VG3S~在~IoU~和~mIoU~上分别比基线提升~12.6\%和~7.5\%，且跨多个~VFM~骨干均有一致增益，验证了几何先验注入策略的通用性。

如果说~VG3S~从特征层面注入几何先验，GaussRender（2025.02）~\cite{2502_05040} 则从监督信号层面强化几何一致性。传统~3D~体素损失（交叉熵、Lovász、Dice 等）逐体素独立计算，不约束相邻体素间的空间一致性，导致悬浮体素、断裂表面等视觉伪影。GaussRender~引入可微高斯泼溅渲染模块，将预测的~3D~占用网格和真实~3D~占用分别通过高斯泼溅投影到~2D~虚拟相机视角（包括透视相机视角和鸟瞰视角），在~2D~空间中比较渲染结果。具体而言，每个体素被建模为一个~3D~高斯球（尺度参数~$s$~与体素边长关联，最优值约为~$s = c/2$，其中~$c$~为体素边长），通过~$\alpha$-混合渲染语义图像。总训练损失为~$L = L_{3D} + \lambda L_{2D}$，其中~$L_{2D}$~惩罚产生不一致~2D~投影的~3D~构型。高斯尺度的选择至关重要：过大则只有最近体素覆盖图像、梯度集中在少数元素上；过小则体素间出现间隙、渲染稀疏激活。GaussRender~作为即插即用模块，无缝集成到~TPVFormer、SurroundOcc、Symphonies~等现有架构中，仅增加约~10\%~的计算和内存开销，但由于~2D~损失加速了学习收敛，达到同等性能所需的总训练时间反而缩短约~17\%。在~SurroundOcc-nuScenes、Occ3D-nuScenes、SSCBench-KITTI360~等多个基准上，GaussRender~显著提升几何保真度，尤其在~RayIoU~等表面敏感指标上表现突出，同时在鸟瞰（BeV）指标上也展现出系统性增益。

\paragraph{SplatSSC：解耦深度引导高斯泼溅}

SplatSSC（2025.08）~\cite{2508_02261}（Decoupled Depth-Guided Gaussian Splatting for SSC）解决了高斯基~SSC~中的两个关键问题：随机初始化导致的基元冗余，以及离群基元引入的错误伪影。在初始化阶段，SplatSSC~利用预训练深度估计器（Depth-Anything-V2）生成深度特征，通过分组多尺度融合（Group-wise Multi-scale Fusion, GMF）模块融合~RGB~和深度特征流，生成精细化的深度图作为几何先验，进而通过深度引导的~lifting~模块生成稀疏但具代表性的初始高斯基元。为抑制离群基元噪声，SplatSSC~设计了解耦高斯聚合器（Decoupled Gaussian Aggregator, DGA），其核心思想是在高斯到体素泼溅过程中将体素占用概率分解为几何项~$\phi'(x)$（由高斯不透明度和空间位置决定）与条件语义项~$\psi^k(x)$（由~softmax~归一化的语义向量和高斯似然加权计算）的乘积。

这种解耦确保离群基元的几何影响不会污染语义预测。在贝叶斯框架下，语义项被建模为高斯混合模型，利用~softmax~归一化语义向量之和为~1~的关键性质，语义归一化项化简为纯几何似然之和，从而实现高效的闭式计算。配合概率尺度损失（Probability Scale Loss），SplatSSC~在~Occ-ScanNet~数据集上达到当时最优性能，IoU~和~mIoU~分别比此前方法高出~6.3\%和~4.1\%，同时延迟和内存开销降低~9.3\%以上。值得注意的是，SplatSSC~仅需~1200~个高斯基元即可达到最高语义精度，进一步增加基元数量反而因冗余而降低~mIoU。

\paragraph{TGSFormer：时序高斯具身 SSC} TGSFormer（2025.12）~\cite{2512_00300}（Temporal Gaussian Splatting for Embodied SSC）面向具身智能场景下的在线~SSC~问题。现有高斯基方法大多在预定义空间内随机初始化大量基元，导致冗余且难以扩展到无界场景；而现有具身~SSC~方法在扩展时面临高斯基元数量持续增长、重复基元冲突等问题。TGSFormer~维持一个持久化高斯记忆用于时序预测，其核心由两个模块组成：双时序编码器（Dual Temporal Encoder, DTE）和置信度感知体素融合（Confidence-Aware Voxel Fusion, CAVF）。DTE~以当前帧的高斯特征为查询、历史高斯特征为键值对，通过基于语义熵的置信度加权交叉注意力机制融合时序信息——语义熵越低（预测越确定），该基元的贡献权重越高。CAVF~随后将重叠的当前与历史基元合并为体素对齐表示：在体素网格中，对落入同一体素的基元进行密度调节和特征聚合，利用幂次变换（power transform）将语义不确定性映射为置信度权重，控制新旧信息的融合比例，从而保持全局记忆的紧凑性。实验表明，TGSFormer~在局部和具身~SSC~基准上均达到当时最优结果，使用的基元数量比此前方法减少~20\%~以上，同时单帧推理延迟也低于~SplatSSC~和~EmbodiedOcc。

\paragraph{SPHERE：语义-物理双表示融合} SPHERE（2025.09）~\cite{2509_11171}（Semantic-PHysical Engaged REpresentation）从表示融合的角度出发，整合体素和高斯双表示来联合利用语义和物理信息。其核心观察是：体素方法善于捕获语义上下文但难以建模物理规律性（如表面平滑性和结构连续性），而高斯方法在几何细节上灵活但空间冗余常导致语义混淆。语义引导的高斯初始化（SGI）模块利用双分支~3D~场景表示（体素分支提供语义上下文，高斯分支提供几何先验）定位具有判别性语义的焦点体素（focal voxel）作为锚点，仅在这些关键位置初始化高斯基元，避免在语义均匀的大面积区域（如路面、墙面）浪费表示容量。物理感知谐波增强（PHE）模块引入语义球谐函数（Semantic Spherical Harmonics），将语义信息编码到高斯基元的球谐系数中，使每个基元不仅承载标量语义标签，还能建模方向感知的物理上下文细节（如不同观察角度下的语义一致性）。PHE~还通过焦点分布对齐（Focal Distribution Alignment）约束体素和高斯两种表示的语义预测分布一致，促进语义-几何的协同优化。在~SemanticKITTI~和~SSCBench-KITTI-360~基准上，SPHERE~在保持语义精度的同时显著改善了几何真实感。

\paragraph{GS-Occ3D：可扩展纯视觉占用重建}

GS-Occ3D（2025.07）~\cite{2507_19451} 提出了可扩展的纯视觉占用重建框架，直接从图像重建占用而无需~LiDAR~标注——这是实现大规模众包数据自动标注的关键前提。GS-Occ3D~使用基于八叉树的高斯面元（Gaussian Surfel）表示，以检测器无关的~SfM~生成的稀疏点云作为场景骨架初始化八叉树。八叉树结构在训练过程中动态扩展或收缩以匹配场景密度和复杂度：每个稀疏体素可生成最多~$m$~个高斯基元，约束在小区域内作为局部面元表示，近似局部表面几何。不同八叉树层级具有不同的体素分辨率——粗糙层级空间分辨率低但覆盖范围大，精细层级分辨率高但仅用于密集区域。场景被分解为静态背景、地面和动态对象，分别采用定制化建模策略：地面作为主导结构元素被显式重建以提高大面积一致性；动态对象（如车辆）通过分割模型分离后被单独建模，以捕获运动相关的占用模式。此外，GS-Occ3D~设计了标签管选管线：通过逐帧分割和多帧聚合增加动态对象的点云密度，再利用光线投射（ray-casting）解决遮挡关系，生成准确的体素占用标签。在~Waymo~数据集上，GS-Occ3D~达到当时最优的几何重建结果（超越依赖单目线索或~LiDAR~输入的方法），其生成的占用标签还能在~Occ3D-Waymo~和~Occ3D-nuScenes~上有效监督下游占用预测模型，展现出良好的零样本泛化能力。

\paragraph{AutoOcc：开放端语义占用自动标注} AutoOcc（2025.02）~\cite{2502_04981} 则进一步实现了开放端语义占用的自动标注。与~GS-Occ3D~侧重几何重建不同，AutoOcc~的核心挑战在于如何从~2D~视觉语言模型的知识迁移到~3D~语义空间。其管线首先利用视觉语言模型（如~InternVL）对多视角图像生成开放端语义类别列表，并通过语义注意力图（Semantic Attention Map）提取~VLM~ Transformer~解码器中各层各头的注意力权重，聚合为~2D~空间注意力图，定位每个语义类别的空间区域。这些~2D~语义注意力图与视觉基础模型的特征结合，通过可微高斯泼溅优化语义感知高斯（Semantic-aware Gaussians）——每个高斯基元携带语义属性，作为中间几何描述符桥接~2D~语义和~3D~结构。为处理动态场景中时间变化的语义一致性，AutoOcc~引入时空一致性约束，对不同时间帧的同一位置进行语义对齐。最终，累积高斯到体素泼溅算法将优化后的语义高斯转换为密集体素占用标注。在~Occ3D-nuScenes~基准上，AutoOcc~无需人工标签即在静态和动态场景中取得鲁棒性能，并在~SemanticKITTI~上展现出跨数据集的零样本泛化能力。

\paragraph{GaussianFlowOcc：稀疏高斯时序占用估计} GaussianFlowOcc（2025.02）~\cite{2502_17288} 从表示和训练两个层面同时革新占用估计流程。在表示层面，它将传统密集体素网格替换为~$N$~个稀疏~3D~高斯基元（每个基元参数化为位置、不透明度、尺度、旋转四元数和语义~logits），消除对~3D~体素标签的依赖，转而通过高斯泼溅将估计结果渲染回~2D~相机空间进行监督。在模型架构层面，GaussianFlowOcc~设计了高斯变换器（Gaussian Transformer），由三个迭代块组成，每个块包含三类注意力：（1）与图像特征的交叉注意力（使用~$M$~个诱导点实现线性复杂度）；（2）高斯基元之间的自注意力；（3）与前帧高斯的时序注意力。高斯头部（Gaussian Heads）为每个基元估计尺度、旋转、不透明度和语义~logits。在时序建模方面，时序模块（Temporal Module）为每个高斯计算~3D~时序光流偏移，将其与前一帧的时序邻居关联，并在高斯泼溅过程中补偿动态对象的运动，从而实现弱监督下的时序一致性。整个模型仅依赖~2D~伪标签（Grounded-SAM~语义分割和~Metric3D~深度估计），在~Occ3D-nuScenes~基准上显著超越此前所有弱监督方法，推理速度比当时最优快~50~倍。

\begin{takeaway}
目标检测、场景图与场景补全将3DGS的语义理解推向结构化场景表示。GS-Occ3D 实现无需LiDAR标注的可扩展纯视觉占用重建，AutoOcc 用VLM实现开放端语义占用自动标注，GaussianFlowOcc 仅依赖2D伪标签在推理速度上比此前最优快50倍。核心趋势是从密集预测走向结构化理解，从封闭词汇走向开放词汇。
\end{takeaway}

%=====================================================================


\section{其他场景理解方法}
\label{sec:10_7}

本节收集 3DGS 场景理解中不属于前述主要类别的重要工作，涵盖变化检测与场景比较、异常检测与质量评估、3D 线段重建等方向。

\subsection{变化检测与场景分析}
\label{ss:10_5}

本节讨论3条技术路线：（1）变化检测与场景比较；（2）异常检测与质量评估；（3）其他3D线段重建方法。

\subsubsection{变化检测与场景比较}

本路线介绍利用 3DGS 的显式 3D 表示在基元空间直接比较以实现场景变化检测的方法。变化检测（change detection）要求系统比较同一场景在不同时间的表现，识别并定位发生变化的对象或区域。3DGS 为变化检测提供了天然优势：其显式的 3D 表示允许在基元空间直接比较，避免了 2D 方法中的视角不一致问题。

\paragraph{基元空间的变化检测}

GD-DIFF（2026.05）~\cite{2605_07203}（Gaussian Drift-DIFF）从根本上重新审视了~3DGS~变化检测的范式。此前方法遵循"渲染-比较"（render-then-compare）范式：将变化前场景渲染为~2D~图像，与变化后图像进行像素或特征残差比较。GD-DIFF~提出将问题从像素空间转移到基元空间——利用高斯基元的原生属性（位置、各向异性协方差、颜色）直接检测变化。核心挑战在于~3DGS~表示的欠约束性：即使场景未发生变化，独立优化产生的基元在数量、位置、形状和颜色上也可能不同。

GD-DIFF~的管线包含四个阶段：几何漂移建模、核评分（几何+外观）、聚合渲染和歧义消解。其核心创新是各向异性漂移模型，将几何漂移归因于两个来源并分别建模为加性各向异性协方差膨胀。第一，\textbf{表示歧义}（representation ambiguity）：3DGS~表示对给定场景不唯一——许多基元配置可以渲染出相同的观测图像，独立优化会收敛到不同的基元解。对于每个基元~$g_i^{(1)}$，找到其在~$G^{(2)}$~中的欧氏最近邻~$g_j^*$，将位移~$\delta_i = \mu_{j^*}^{(2)} - \mu_i^{(1)}$~分解为法向和切向分量：$d_n^{(i)} = |n_i^\top \delta_i|$，$d_t^{(i)} = \|\delta_i - (n_i^\top \delta_i)n_i\|_2$，然后通过上四分位数~$Q_{0.75}$~统计得到漂移尺度~$u_n^2$~和~$u_t^2$，组装出各向异性膨胀矩阵~$U_i = u_t^2 I_3 + (u_n^2 - u_t^2)n_i n_i^\top$~并叠加到基元自身协方差上得到~$\tilde{\Sigma}_i = \Sigma_i + U_i$。第二，\textbf{观测不确定性}（observation uncertainty）：通过累积每个基元在可见相机上的~Fisher~信息矩阵（FIM）~$H_i = \sum_{c \in C_i^{\text{vis}}} \frac{1}{d_{i,c}^2}(I_3 - v_{i,c}v_{i,c}^\top)$~来量化——多角度近距离观测的基元产生良态的~$H_i$，而单一视角观测的基元~$H_i$~则近乎奇异（投影~$I_3 - vv^\top$~编码了每个相机仅约束图像平面方向而非视线方向）。最终有效协方差~$\Sigma_i^{\text{eff}} = \tilde{\Sigma}_i + sH_i^+$~融合了两种漂移来源，其中~$s$~是数据驱动的缩放因子。

在此基础上，几何核通过各向异性~Mahalanobis RBF~评估基元间的几何相似性：$k_\text{geo}(g_i, g_j) = \exp(-\frac{1}{2}(\mu_i - \mu_j)^\top M_{ij}^{-1}(\mu_i - \mu_j))$，其中~$M_{ij} = \Sigma_i^{\text{eff}} + \Sigma_j^{\text{eff}}$。外观核则通过数据驱动的颜色带宽~$\sigma_c^2$~评估~DC~颜色差异：$k_\text{app}(g_i, g_j) = \exp(-\|c_i - c_j\|^2 / 2\sigma_{c,i}^2)$。每个基元的几何变化评分~$s_i^\text{geo} = 1 - \max_j k_\text{geo}(g_i, g_j)$~和外观变化评分~$s_i^\text{app}$~分别计算，最终通过可观测性加权的~alpha~混合渲染到任意视角的变化图。

基元空间操作带来两个独特优势：变化图天然多视角一致（无需额外优化目标），几何和外观变化被分别评分——可以无监督地区分结构性变化（几何不匹配）和表面级变化（仅外观改变）。在真实世界基准~PASLCD~上，GD-DIFF~在~mIoU~上比此前最优方法提升约~17\%。

\paragraph{多视角免位姿变化定位}

Galappaththige et al.（2024.12）~\cite{2412_03911} 提出了免标签、免位姿的多视角变化检测方法，将变化信息编码为~3DGS~的一个额外通道。该方法首先从参考图像构建~3DGS~参考场景表示，然后在变化后图像视角下渲染参考图像并与之比较，利用~DINOv2~特征得到特征感知变化掩码、通过结构差异得到结构感知变化掩码，两者互补融合为候选变化掩码。接着以参考场景的~3DGS~为初始化，学习一个变化感知~3DGS（Change-3DGS），其中每个高斯增加独立的变化不透明度因子~$\alpha_c$~和变化幅度~$\beta$——前者允许被移动/移除对象对应的高斯在~RGB~渲染时降低不透明度，同时保留变化掩码中的高变化幅度；后者直接编码变化强度并渲染为变化图。为避免变化高斯被标准不透明度剪枝误删，剪枝条件改为~$\min(\alpha, \alpha_c) < \tau$。仅需~5~张变化后场景的图像，即可学习~3DGS~中的变化通道并生成变化掩码，在复杂多对象场景中~mIoU~和~F1~分别比单视角基线提升~1.7~倍和~1.5~倍，同时支持未见视角的变化掩码生成。

\paragraph{物体级变化检测}

3DGS-CD（2024.11）~\cite{2411_03706} 是首个基于~3DGS~的物体级变化检测方法，通过比较两组未对齐图像估计~3D~对象级变化。其模块化管线包含六个阶段：（1）在变化前图像上训练~3DGS~模型；（2）利用~SuperGlue~将变化后图像定位到~SfM~稀疏点云坐标系中；（3）在每个变化后视角下渲染参考~RGB~图像，通过~EfficientSAM~编码器的余弦相似度图粗检测~2D~变化区域，再用包围框提示~EfficientSAM~获得高置信度的移出（move-out）和移入（move-in）~2D~对象掩码；（4）\textbf{跨视角对象关联}：对每个~2D~对象段，通过投影到~3D~并在其他视角反投影来建立视角间的对象对应关系，逐步聚合为包含密集点云和稀疏视觉特征点的~3D~对象模板；（5）基于视锥裁剪和投影的~3D~对象分割；（6）通过~ICP~配合~RANSAC~估计~6DoF~刚体变换，并用~3DGS~渲染的~RGB-D~图像对进行联合~ICP+光度~精化。该方法在杂乱环境中仅用一张变化后图像即可在~18~秒内准确识别变化，精度比此前最优辐射场方法高~14\%，速度快三个数量级——这一效率优势得益于~3DGS~基于光栅化的实时渲染（相比~NeRF~的射线投射）和~EfficientSAM~的轻量级推理。

\subsubsection{异常检测与质量评估}

本路线介绍利用 3DGS 高效新视角合成能力实现姿态无关的 3D 异常检测与工业质量评估的方法。与上一路线关注场景级时序差异不同，异常检测（anomaly detection）聚焦于单个对象表面相对于无缺陷参考的局部偏差。两者共享一个核心难题：查询视角未知且不与参考视角一致。3DGS 在工业质量检测领域的应用正在兴起，其高效的新视角合成能力为姿态无关（pose-agnostic）的 3D 异常检测提供了新的技术路径。

\paragraph{姿态无关的异常检测}

SplatPose（2024.04）~\cite{2404_06832} 首次将~3DGS~引入姿态无关的~3D~异常检测。给定一个~3D~对象的多视角无缺陷图像，SplatPose~先用~3DGS~建模对象的~3D~高斯点云（球谐阶数限制为零阶以保持旋转不变性并节省计算），然后通过\textbf{两阶段可微位姿估计管线}推断未知查询视角的相机位姿。粗估阶段采用~NeRF~的渲染结果检索最相似的训练视角；精化阶段则不直接优化相机参数，而是将~SE(3)~变换参数化为螺旋轴~$\mathcal{S} = (\omega, v) \in \mathbb{R}^6$~和平移距离~$\theta \in \mathbb{R}$，通过~Rodrigues~公式对整个~3D~高斯点云施加变换~$T = e^{[\mathcal{S}]\theta}$~后渲染图像，再用~$L = (1-\lambda)L_1 + \lambda(1-L_\text{SSIM})$~损失将梯度回传至位姿参数——由于梯度从损失函数流经渲染图像到~3D~点云再到变换参数，整个位姿估计过程完全可微。精化收敛后，在估计位姿下渲染无缺陷参考图像，与查询图像进行像素级特征比较完成异常检测。相比基于~NeRF~的方法（OmniPoseAD），SplatPose~在训练和推理速度上均达到当时最优，且检测性能更佳——即使训练数据更少。

\paragraph{实时推理的工业级异常检测}

SplatPose+（2024.10）~\cite{2410_12080} 在~SplatPose~基础上引入混合表示架构，将原本~3DGS~同时承担的定位和新视角合成两个职责分离，分别交给~SfM~模型和~3DGS~模型。

\textbf{SfM~定位管线}：采用~hloc~框架（NetVLAD + SuperPoint + LightGlue）构建参考数据库和稀疏~3D~模型，推理时通过图像检索找到~15~张最相似参考图像后进行特征匹配和三角定位，无需可微优化即可快速获得查询位姿。

\textbf{3DGS~新视角合成}：以~SfM~稀疏点云（而非随机初始化）初始化~3DGS~训练，点云分布在真实物体表面上提供了更好的起始几何；同时恢复了完整的~3~阶球谐函数（SplatPose~仅用~0~阶），获得更精确的视角相关色彩化；并降低致密化频率至每~1000~次迭代一次（因为~SfM~点云已充分覆盖物体表面）。虽然需要额外计算~SfM~模型，但推理速度比~SplatPose~快~147~倍，且在稀疏视角设置下优势更为显著。在姿态无关异常检测基准~MAD-SIM~数据集上达到新的最优性能。

\paragraph{高斯泼溅引导的视频异常检测}

GS-MoE（2025.08）~\cite{2508_06318} 将~3DGS~的思想迁移到视频异常检测（VAD）领域。弱监督~VAD~仅提供视频级标签，需在帧级进行预测。经典的多实例学习（MIL）范式仅依赖得分最高的~$k$~个片段（snippet）进行优化，导致模型过度关注最异常的片段而忽略边界信息。GS-MoE~的核心创新是\textbf{时间高斯泼溅（TGS）损失}，通过在时间轴上泼溅高斯核来增强弱监督信号。具体而言，首先在模型预测的异常得分时间序列上检测局部极大值（"峰值"），每个峰值~$P_i$~的宽度~$W_i$~由前后单调变化的片段数决定；然后以峰值为中心初始化高斯核~$\mathcal{G}_i$，对峰值位置置~1，对宽度范围内得分高于~$s_{P_i} - \sigma_i$~的片段也置~1（其中~$\sigma_i$~为以峰值为中心的正态分布标准差）；最后每个核通过~$f_i(t) = \sum_{\mathcal{G}_i,t} \cdot \exp(-\frac{(t-P_i)^2}{2\sigma_i^2})$~泼溅到整个时间轴上，使得异常事件的完整时间窗口都能获得监督信号，而非仅依赖~top-$k$~片段。GS-MoE~结合\textbf{混合专家（MoE）}架构：每个专家专精于特定异常类型（如盗窃、爆炸、交通事故），通过门控网络将粗粒度任务特征分配给对应专家，专家再输出类别特定的精细异常得分。在~UCF-Crime~数据集上达到~91.58\% AUC（比此前最优提升~3.56\%），在~XD-Violence~和~MSAD~上同样取得最优结果。

\subsubsection{其他：3D 线段重建}

本路线介绍利用 3DGS 训练后高斯点自然聚集于物体边缘的分布特性进行 3D 线段重建的方法。除了变化检测和异常检测这类判别性任务，3DGS 还可用于增强 3D 场景的几何抽象表达。线段是场景结构的高效抽象表达，但传统 3D 线段重建方法面临 2D 投影不稳定或 3D 点云噪声稀疏的挑战。3DGS 训练后高斯点自然聚集于物体边缘和色彩边界处，这一分布特性为线段重建提供了天然的几何先验。

\paragraph{三维高斯上的线段表示}

LineGS（2024.12）~\cite{2412_00477} 提出了~3DGS~场景上的~3D~线段表示方法，利用场景边缘处高密度分布的高斯点来优化几何引导方法（如~Line3D++、ELSR）生成的初始线段。其核心是\textbf{线段-高斯对齐机制}，通过针对性后处理消除初始线段的五类典型缺陷：（1）\textbf{位置偏移}——以线段为中心轴定义圆柱空间~$C(\ell, r)$，利用八叉树高效检索圆柱内的高斯点集~$X(\ell)$，通过线性回归将线段平移到高斯点最密集的位置；（2）\textbf{过度延伸}——通过二分搜索定位高斯密度骤降的边界点并裁剪端点，快速找到线段的真实终止位置；（3）\textbf{离群线段}——基于全局平均密度设置阈值，密度低于平均值的线段直接移除；（4）\textbf{重复线段}——计算线段间亲和矩阵，高相似度的线段聚类后在簇内合并；（5）\textbf{断裂}——同类簇中距离近但不连续的线段被桥接。整个后处理流程在保持几何有效性的同时，实现了线段与~3DGS~底层数据分布的对齐，在几何精度和模型紧凑性上均显著优于基线方法，同时提供了~3DGS~场景的一种简洁抽象表示。

\begin{takeaway}
其他场景理解方法展示了3DGS在变化检测、异常检测和3D线段重建等非分割任务中的价值。LineGS 利用3DGS训练后高斯点自然聚集于物体边缘的分布特性实现3D线段重建，这一工作揭示了3DGS的隐式几何先验可服务于高层场景抽象任务。
\end{takeaway}

%=====================================================================

\section{近期研究（2026 年 8--11 月）}
\paragraph{GaussianSelector：将用户笔迹提升为高斯超点图的前背景约束}
2026.08~\cite{2608_01492}。GaussianSelector 将已重建场景的高斯按外观、不透明度和空间邻近组织为超点图，用光栅时的透射与覆盖把前背景笔迹落到对该像素有贡献的原语，再在图上给其余超点求选择结果。方法直接操作已有高斯，无需为选择任务训练附加网络。

论文在 LLFF-NVOS 的八个任务和 3D-OVS 五个场景比较，并区分单轮与追加笔迹。结果必须结合用户输入视图数和交互轮数解读，遮挡区域能否选中仍取决于可见笔迹与图中连接关系。

\paragraph{Super-Gaussian：以分层分组支持体数据的虚拟现实交互}
2026.08~\cite{2608_04475}。Super-Gaussian 将体可视化的高斯基元按特征聚成可交互单元。在虚拟现实中，用户先用稀疏种子与图传播选取区域，再作组级修正和原语级细化；自然语言界面用于语义标注与查询。

论文将几何选择对比、头显渲染时间与测地 SLIC 消融分开评估。所用初始聚类由预处理得到，窗口参数控制分组粗细，在低对比或边界不明的区域可能产生不稳定分组；图传播也受交互种子质量影响。

\paragraph{OutLangSplat：融合区域结构与语义以理解无人机户外场景}
2026.08~\cite{2608_04560}。OutLangSplat 将二维视觉语言特征和三维高斯邻域结构在区域级对齐，随后以像素上的高斯渲染贡献及跨视图语义一致性加权聚合，减少遮挡和噪声视图的错误响应。

论文在 Buildings1、Buildings2、Polytech 和 Campus 四个无人机户外场景验证开放词汇语义分割与实例定位。它不重新优化场景重建，几何不完整或高斯位置错误会影响对应和聚合，作者将小对象语义与重建联合优化列为后续方向。

\paragraph{CDSeg：用可渲染高斯载体把二维标签回写到三维}
2026.08~\cite{2608_05482}。CDSeg 用已标定视图渲染产出的像素至原语关联，把外部二维掩码的标签融合到可见原语。模式 I 保留输入点索引，便于回写网格或激光雷达点云；模式 II 将标签留在已有、已优化的高斯场景中，可随视图重新渲染。

论文用 DesktopObjects-360 和 NeRDS-360 评价优化载体上的实例任务，用 ScanNet-v2 验证语义标签回写，并在 KITTI-360 的静态车辆窗口考察可见性限制。从未出现在输入视图的表面不会获得标签，已有高斯的空缺和失真也会传递到结果。

\paragraph{InstanceSplat：从无位姿图像前馈重建带实例身份的高斯场}
2026.08~\cite{2608_07144}。InstanceSplat 由无位姿多视图预测共享高斯表示，在同一场景中编码几何、外观、实例和语言语义。跨视图实例对齐为可渲染的实例特征提供监督，实例结构进一步参与联合重建与语义学习。

论文在 ScanNet 与 ScanNet++ 混合训练，在 50 个留出 ScanNet 场景比较 2、4、8 和 16 个输入视图。评估分辨率为 448×448；论文分别报告新视角渲染与分割指标，并指出 4 视图结果中的平均准确率会受采样目标的类别比例影响。

\paragraph{LEGO：把部件、对象与场景的语义按结构层编码}
2026.08~\cite{2608_10057}。LEGO 将多视图 SAM 分割重排成三维一致的结构层级，再分别学习高斯子特征、绑定 CLIP 语言特征并组织为分层场景图，以减少粗层和细层监督互相干扰。场景图还能供语言模型按层检索实体和上下文。

论文在 NVOS 和 SPIn-NeRF 评价提示分割，在 LERF-OVS 与补充的 3D-OVS 等设置检查开放词汇理解，报告 NVOS 的 mIoU 提升最高达 1.6 个百分点。方法依赖初始几何和多视图掩码；开放词汇结果对应论文使用的语义模型和评估场景。

\paragraph{CausalSplat：以高斯语义场与场景图支撑指代推理}
2026.08~\cite{2608_11150}。CausalSplat 先将二维分割提升为高斯语义场，再把实例的多模态属性及空间连接组织成图。视觉语言模型结合图解析复杂指代表达并定位目标，输出可分割的三维实体。

作者构建 Causal-LERF 和 Causal-ScanNet，分别在渲染像素和三维点上计算 mIoU。补充实验拆分空间、常识、可供性和预测或反事实问题。这里验证的是通过这些描述识别目标的表现，输出质量取决于感知和语言推断阶段所获得的证据。

\paragraph{Seed2GS：一次识别并传播对象种子以提取高斯}
2026.08~\cite{2608_11928}。Seed2GS 在冻结的三维场景中，用文本对一张带相机位姿的参考渲染识别目标，借虚拟视图跟踪扩大覆盖，然后通过固定光栅化为各高斯拟合临时前景分数。它可以不使用原始重建相机，但查询仍需给定一张参考视图和该视图的位姿。

论文在用户选择能看见目标的 LERF-MASK 参考视图条件下报告 92.1\% mIoU 和 9.3 秒，并在 3D-OVS 报告 95.7\% mIoU。前一数字未对任意参考视图平均；识别错误、严重遮挡、缺少几何和长距离跟踪漂移都可能导致失败。

\paragraph{QAGaussian：用关系推理选取自由指代表达指定的高斯}
2026.08~\cite{2608_16103}。QAGaussian 从语言查询产生多尺度候选槽，再通过候选关系图和几何约束组合证据，输出逐高斯的掩码概率。例如带参照物的表达中，空间关系可用于减轻目标与参照物被一起选中的问题。

作者在 Mosaic3D-5.6M 预训练并用该数据的验证集选定设置，随后测试 ScanRefer、ReferIt3D、Multi3DRefer、ReferSplat 和部件数据。不同标注被转换到共同评价设定，比较结果应按这些对应协议解释；预测仍依赖高斯重建和指代数据的覆盖。

\paragraph{GroupForward：以实例分组承载前馈场景语言语义}
2026.08~\cite{2608_17535}。GroupForward 从稀疏、无位姿且未标定图像前馈重建高斯，以紧凑实例嵌入建立跨视图分组。开放词汇特征按实例聚合，随后以实例构建场景图并结合多视图观察，供视觉语言模型搜索目标和输出指代分割。

作者以预训练 VGGT 初始化几何部分，并使用 8 维实例嵌入。论文将语义三维重建与场景指代推理分开评价；低维只描述实例分组嵌入，语言特征仍在实例层维护。图推理能否定位目标还取决于前馈阶段能否把目标分成有效实例。

\paragraph{LetOccVote：利用跨帧投票改进弱监督三维占据预测}
2026.09~\cite{2609_04846}。LetOccVote 的深度投票把多帧在共同空间里的支持用于提升和深度监督，再通过高斯将多帧信息传给三维占据表示。语义投票用同意和冲突区域筛选不可靠的二维伪语义监督。

在 Occ3D-nuScenes 的主表，较大模型报告 53.27 IoU、20.39 mIoU，另按论文标星类别集合报告 23.11 mIoU。两个语义数值使用不同的类别口径。方案借助跨观察共识减轻伪标签噪声，评价以多帧驾驶场景和二维伪标签监督条件为背景。

\paragraph{LangStreet：把语言证据绑定到锚点解码的持久地址}
2026.09~\cite{2609_11616}。LangStreet 将随视图生成的高斯原语用作路由，把多视图语义观察累加到持久槽位并向上汇成锚点摘要。对于观察稀疏的槽位，通过保持锚点方向的层级补全增强稳定性，再用相对锚点编码保存局部差异。

作者在 KITTI-360、Virtual KITTI 2 和 Waymo 上共享冻结的场景重建、解码器及二维语言引擎进行比较。论文显示紧凑字段可以接近完整字段并节省存储；结论对应同一预处理观察和持久结构，不能解释为重新训练重建或二维模型后的端到端提升。

\paragraph{MoQSplat：将高斯场景按空间和质量层送入渐进传输}
2026.09~\cite{2609_18624}。MoQSplat 在 Media over QUIC 的内容组织中，将场景空间区域映射到 Track，再以语义分组和由粗到细的质量层组织对象。接收端按关注区域及层级决定请求优先级，发布者和转发者无需为各接收方维护单独的选层状态。

论文先测内容准备质量，再测网络条件变化下的端到端接收。两组使用的共同参照是 Mip-NeRF 360 的 Bicycle，内容准备以固定内参、距离和高度的绕行轨迹计算 VMAF。因此不透明度排序和带宽适配的观察属于单场景初步实验。

\paragraph{NormLift：用跨视图支持校准高斯语言特征置信度}
2026.09~\cite{2609_18898}。NormLift 将单位球面上的逐高斯余弦对齐写成可闭式求解的提升问题，结果对应反投影特征的归一化。聚合前的向量范数反映观察的一致程度，再以有效视图支持校准，作为邻域模式投票修正低可信原语的依据。

作者以 OpenGaussian 协议在 ScanNet 测试开放词汇三维分割并比较 mIoU 和 mAcc。后处理以既有高斯场景为基础，置信度与投票利用观察和邻域的证据，不补回原场景缺失的表面。

\paragraph{ParticleSplat：由多视图自监督学习可编辑的对象级潜粒子}
2026.09~\cite{2609_19463}。ParticleSplat 从已知相机姿态和内参的观察学习结构化三维潜粒子，并将粒子前馈解码为显式高斯。粒子形成对象级控制接口，改变对象潜变量或位置后可观察对应的新视图变化。

论文把分解、新视图渲染和操作学习分为不同实验，其中真实桌面场景使用手持手机捕捉多视图，机器人决策收益另有仿真设置。真实捕捉无需固定多机阵列仍需获得带位姿的多视图，不能理解成相机位姿无关。

\paragraph{CoRef-GS：通过地图配准和实例融合做跨机器人语言指代}
2026.09~\cite{2609_20586}。CoRef-GS 先构建可按实例查询的局部语义高斯地图，通过跨机器人配准和实例融合将观察统一到共享坐标，再按查询视角使用空间关系选出语言指定的对象。

作者在含真实和模拟室内场景的 CoQuad-Ref 上区分地图配准、开放词汇查询和协作指代结果。主结果报告粗 Sim(3) 初始化已改善配准，完整方案在模拟设置中进一步降低误差。证据来自互补观察的室内四足机器人设置，不代表大规模动态车队的通信和时间同步表现。

\paragraph{建筑感知卫星高斯重建：经掩码质检和区域误差选择权重}
2026.09~\cite{2609_25578}。该工作用建筑掩码调节卫星高斯重建的区域损失，控制器同时参考掩码质检、场景元数据和验证误差决定是否采用建筑先验。系统记录评价结果，使数字表面模型的选择可追溯到场景条件。

论文用注册后的 DFC2019 风格 DSM 与语义标注计算建筑、植被和全场 MAE。JAX\_004 上的建筑权重 2.0 将建筑区域误差从 0.844 米降至 0.806 米，但分阶段方案在全场误差上的折中更好。改善以通过掩码质检的场景为条件，不能只用单一建筑误差决定部署。

\paragraph{PePESeg3D：在几何和特征两阶段注入感知先验}
2026.09~\cite{2609_28645}。PePESeg3D 首先以单目深度和分割掩码引导高斯几何贴近结构边界，再通过深度颜色线索与跨视图质心约束学习多尺度语义，以补充二维掩码在部分粒度上的缺漏。

论文在 SPIn-NeRF、LERF-Mask 等设置报告标签传播与提示分割结果，另用 PSNR、SSIM 和 LPIPS 评价重建，两个任务的指标不能互代。作者观察到共享特征空间的尺度门控会使不同粒度的掩码互相影响，尤其在邻近尺度检索时。

\paragraph{SplatLabel：以含生命周期的动态高斯生成三维语义伪标签}
2026.09~\cite{2609_29836}。SplatLabel 将连续运动、生命周期和存在边界纳入动态高斯属性，让二维语义线索与激光雷达几何在时间上共同决定三维语义和体积占用。高斯在此是时空伪标签生成器。

论文以 SemanticKITTI 验证序列 08 评价，使用 SAM 3 和类别名称直接生成语义特征。方法依赖精确相机和 LiDAR 同步且显存开销较大，恶劣光照、天气或传感器质量退化会影响结果。

\paragraph{EviSplat：把多视图语义证据保留到查询时再聚合}
2026.09~\cite{2609_34853}。EviSplat 保留实例的各视图特征，并为高斯学习反映不同可见外观的紧凑证据场。查询时先按相关性选取实例观察，再用原语局部的可靠证据细化分数，不需要为了新查询重训表示。

论文在 LERF-OVS 和 3D-OVS 评价渲染掩码，在 ScanNet 按 OpenGaussian 设定评价 10 个场景的三维点，并将渲染后查语言特征与在三维先选高斯的方案分列。结论需区分这两类输出及其协议；保留逐视图证据也改变了语义字段的存储结构。

\paragraph{Imagine3D-LLM：把紧凑高斯重建变成多模态问答的训练任务}
2026.09~\cite{2609_38177}。Imagine3D-LLM 在多视图视觉标记后加入可学习的高斯摘要标记，把它们解码为紧凑场景表示，并与语言预测共同训练。对多图问答而言，高斯在这里提供显式空间结构监督，促进跨视图统一物体表征。

论文比较多类空间推理与三维理解基准，还分析生成高斯和内部特征。作者指出不用紧凑高斯教师也能学习摘要标记，但达到相近下游结果需要更多训练步骤。问答收益与场景几何的细节准确度需要分别度量。

\paragraph{MEGA：从分割高斯的虚拟视图蒸馏对象级网格}
2026.10~\cite{2610_01707}。MEGA 从已分割的高斯对象渲染一组虚拟视图，对每个对象以光度和掩码监督蒸馏神经表面，输出可用于对象级操作的网格。分割和表面蒸馏承担不同环节，蒸馏不会绕过初始对象分离。

作者的定量设置使用三排、每排 72 个虚拟视点。在 LERF-Mask 掩码的额外实验中，MEGA 平均 mIoU 为 90.1，作为分割教师的 SAGO 为 91.0，不能把该指标描述成超过教师。网格提取质量与该分割和虚拟视图采样有关。

\paragraph{MaRO-GS：以跨视图掩码可靠性约束对象重建}
2026.10~\cite{2610_06472}。MaRO-GS 根据几何一致性过滤不可靠对象掩码，再限定高斯密度控制只支持具有对象证据的位置，以抑制背景高斯混入，轮廓训练为对象边界提供约束。

主渲染比较采用 LERF-Mask、Mip-NeRF 360 和 Tanks and Temples 上统一的掩码输入协议，测量与掩码目标的渲染一致性。作者指出严重错误的掩码仍会破坏结果；Hard DTU 子集较小，保守过滤的召回只有 4.6\%，不能将整体改善全归因于过滤。

\paragraph{OpenSplatGraph：从可靠性语义高斯场按需生长场景图}
2026.10~\cite{2610_07569}。OpenSplatGraph 为稠密语义高斯地图保存可靠性信息，按语言查询提取实例并将其关联到现有节点。持久的场景图逐步维护对象身份、属性和关系，为随后查询保留可复用的结构。

实物实验使用搭载 RGB-D 相机的 Stretch 3 移动机械臂，在探索中融合语义地图，另有三维分割和场景图评价。该框架以查询触发实例提取，不维护独立于查询的完整对象清单；拥挤和重遮挡仍会使轻量关联变得不可靠。

\paragraph{训练后语义提升：分别估计检测器、提升和表示误差}
2026.10~\cite{2610_08756}。该方法不重训既有高斯。给定单个目标类别，它累积各视图的目标和非目标证据，以高阈值选择种子，再以较低阈值接受含种子的连通区域。评估将预测映射到相机可见且有标注的网格顶点，从而排除不可见表面带来的混淆。

论文分别比较标注掩码和检测器掩码，在 Replica 选定设置并在 ScanNet++ 验证迁移。类若未出现在任何输入掩码就没有可提升的证据，且逐类处理可能让同一高斯落入多个类别，不能将该方法直接视为单次多类别分割。

\paragraph{可测量语义高斯：面向建筑巡检的区域和中心深度监督}
2026.10~\cite{2610_13143}。该工作在高斯优化中联合使用目标区域约束与逐点深度监督，直接约束三维高斯中心的位置，并构建语义场，使建筑高斯既可观察部件类别，也能从显式结构中取点测量。

作者分别在合成 QuakeCity 和手机采集的 HKUPB 上评价区域 RGB、中心几何与语义表现。中心深度监督改善的是总体对齐，局部仍可能因原语稀疏影响精确取点。作者指出细小裂缝和局部结构很难重建，同时输入位姿和深度仍依赖摄影测量。

\paragraph{DIGS：在高斯上保留开放词汇标签分布}
2026.10~\cite{2610_16177}。DIGS 用逐高斯的前若干候选标签分布替代单一语义向量，并以减温 Dirichlet 更新限制单帧观察的影响，减轻来自多实体边界的语义混合。实例级视觉和文本表征帮助在不同粒度下给高斯赋分。

论文以四个开放词汇基准测试，其中 LERF-OVS 的四个场景按 OpenGaussian 三维原生协议评价实例定位与分割。映射时间、显存和单次查询另行报告，方法依赖由基础分割和语言模型得到的输入证据。

\paragraph{目标感知高斯语言地图：最小特征注入与中心化定位}
2026.10~\cite{2610_21619}。该方法将指定目标在各视图中的视觉特征向文本嵌入靠拢，以闭式求出的最小强度注入，再将更新后的特征提升至高斯地图。查询时的中心化读出消除注入产生的共模分量，用以返回目标三维位置。

论文在三个基准将开放词汇分割与语言目标定位分开评价，并展示真实四足机器人。未做中心化的消融在三个数据上均可能因查询分数高度相关而退化；例如 ScanNet 的首选定位从基线 72.9 降至 30.0。论文的导航位置比较固定底层地图，不能等同于不同系统完整导航管线的比较。

\paragraph{O-SCD++：以递归统计在固定高斯地图检测变化}
2026.10~\cite{2610_24462}。O-SCD++ 利用高斯线性合成，将由新观察获得的变化证据累积到固定参考场景的逐高斯状态中。递推更新和闭式决策替代反复缓存、重放观测并迭代优化的处理方式。

作者以真实多视图场景变化的 PASLCD 为主要实验，场景包含多个几何和表面变化，并有不同光照的重访数据。系统无法恢复输入线索未识别到的变化；例如只用语义线索时，同类表面的颜色改变可能被忽略。

\paragraph{PanoGraphGS：统一高斯全景分割和楼层至对象的层级场景图}
2026.10~\cite{2610_27696}。PanoGraphGS 在共享的高斯几何上训练类别原型和持久实例关联，从同一表示解码新视角全景掩码与原语级语义。语义及空间结构进一步组织成楼层、房间、背景实体与对象的层次图。

作者使用 ScanNet++、Replica 和自采 NaviGS 验证，论文将二维全景分割与三维场景图的楼层、房间和对象节点准确率分开报告。开放式厨房一类缺少显式边界的空间难以划分房间；图构建在当前版本中离线运行。

\paragraph{自动编辑参考高斯：为在线贝叶斯变化检测合成证据}
2026.10~\cite{2610_28490}。该工作用自动编辑后的参考高斯生成可监督视图，预先拟合场景专用的变化证据模型。上线后以局部贝叶斯因子和递推的历史证据估计变化，无需对在线数据反向传播或重放观测。

实验把新占据与腾空区域都纳入变化掩码，以真实 PASLCD 和合成 CCS3D 为比较环境。作者报告 PASLCD 的前景 mIoU 达 67.64\% 并以 17.2 帧每秒处理。简单对象编辑与重光照覆盖不了所有真实变化；证据累积在固定参考几何上，限制了新出现对象的位置表达。

\paragraph{R4DGS：为动态高斯场构建可持续指代的实体库}
2026.10~\cite{2610_30326}。该工作将动态场景中的语言指代表述为目标实体选择。ReferGaussian 从预计算的四维高斯及多帧视觉掩码构建持久实体，存入 EntityBank，以多模态模型解释短语与时序约束，并在实体层输出结果，避免另训逐场景语义场。

评价使用 R4D-Bench-QA 和四个场景的 4D LangSplat 协议，把时序状态、多目标与无目标指代纳入测试，并对比语言场和分割先行路线。方法依赖预计算的动态场景与掩码轨迹，无法单靠实体推理恢复底层场景未观察或未重建的内容。

\paragraph{RecallGS：用当前几何核实可回用的历史语义观察}
2026.10~\cite{2610_31293}。RecallGS 从历史图像检索与当前查询相似的帧，再通过透明度合成中的像素贡献检查历史帧是否真正支持当前高斯，并只融合满足几何支持的线索。它将记忆用于语义特征补充，而不需再训练新的语义网络。

论文沿用 EmbodiedSplat 的预处理测试场景，在 ScanNet 分 10、15、19 类评估，ScanNet++ 采用 20 类设置。结果限于这些室内地图和已有历史视图；检索得到视觉相似帧不等于看到了同一表面，几何验证的用途正是过滤此类误配。

\paragraph{GEMS：保存高斯变化量并按时间重渲染场景记忆}
2026.10~\cite{2610_33158}。GEMS 针对反复访问的半静态环境，先建高斯地图，再只记录检测到变化的区域，以节省各时点的场景存储。问答代理按问题选择历史时间和相机位姿渲染场景，逐步收集支持答案的视觉证据。

GEMS-Bench 由 3 个真实室内场景与 12 个模拟 HSSD 场景构成，分别含 244 和 229 对问答。问答用模型裁判进行分级。实验基于离散重访和半静态布置，不检验连续快速运动的重建或时间跟踪。

\paragraph{RelateGS：筛选高斯语义观察并回查关系证据}
2026.10~\cite{2610_33868}。RelateGS 在既有高斯重建和带位姿 RGB-D 观察上工作。它在提升语义时过滤弱证据，并以持久对象图汇总多视图信息。当支撑某一关系的观察不足时，系统从既有视图补查相关线索，减少因单视图误判造成的指代错误。

指代实验使用 ReferIt3D 的 Sr3D+ 和 Nr3D，在沿用先前协议的 8 个 ScanNet 室内场景上比较。该机制依赖能正确关联对象和获取支持关系的帧；回用已有图像不能补足从未被观察的线索。

\paragraph{LegoOcc++：共享前馈高斯的占据和开放词汇语义推理}
2026.10~\cite{2610_34093}。LegoOcc++ 从输入图像前馈预测带语言特征的高斯。基元位置、协方差和不透明度经 Poisson 式聚合估计体素占据，并由原语嵌入预测占据体素的类别。实例一致性与共视支持用于减少同一物体内部的语义破碎。

论文分别讨论室内单目和室外驾驶设置。室外方案可以不使用占据标注，室内实验仍需要二值占据监督；不能将二者统称为无需三维几何监督。作者指出视觉与几何差异大的概括类别如家具仍较难统一校准。

\paragraph{ClusterSplat：将高斯指代转为场景适应的离散簇选择}
2026.10~\cite{2610_36462}。ClusterSplat 学习带实例信息的高斯特征，以描述长度准则生成场景适应的簇，再按语言查询为各簇排序。最终可渲染被选簇的分数产生二维掩码，从而把按单个原语过滤改为以候选对象为单位选择。

论文在配对 ScanNet v2 与 ScanRefer 的场景以及 Ref-LERF 的二维协议比较。作者指出主要瓶颈是候选簇：相邻同类粘连或把一个对象拆成数簇都可能让首选验证通过但空间 IoU 仍不足。

\paragraph{按查询分割高斯：从一个二维枢轴传播并投票到三维}
2026.11~\cite{2611_19429}。该方法以密集视觉语言特征给文本查询生成相似度图，选择二维枢轴区域，随后以视频分割器在有限视图预算内跟踪并细化掩码。根据每个高斯的可见证据投票，形成此次查询对应的三维实例掩码。

论文在 LERF-OVS 将首次查询所需时间、分割 mIoU 与定位准确率分别报告，并另列多个室内基准。该方案使用预训练模型和图像特征处理；视点变化大时跟踪可能漂移，严重遮挡也会使最终掩码不完整。

\paragraph{VODA：从屏幕掩码监督扩展到深度簇内的全部高斯}
2026.11~\cite{2611_21383}。VODA 将同一像素参与渲染的高斯按局部深度分组，在目标深度簇中直接监督每个高斯的特征，避免训练只惩罚对渲染掩码贡献较大的表面原语。该目标可以作为现有语言和实例分割方案的特征学习组件。

论文以 Removal-3D 的移除前后深度变化度量检验体积选择完整性，并在 SPIn-NeRF 和 360-USID 的成对实拍中验证遮罩质量。实拍深度差参考由独立拟合的有对象与无对象高斯场构造，作为代理指标使用，不能视为真值三维几何。


\section*{本章小结}
\addcontentsline{toc}{section}{本章小结}

\begin{futurebox}{本章总结与展望}{ch10_summary}

本章综述了~3DGS~分割与场景理解领域的相关论文，按语义分割与开放词汇理解、特征与检测、场景理解与建图三大主题展开，涵盖开放词汇场景理解、3D~分割、特征蒸馏与表示、异常检测与质量评估、3D~目标检测、6D~位姿估计、占用网格与语义场景补全、变化检测与场景比较、场景图与语义建图等~9~个子方向。核心脉络可以概括为：

\begin{itemize}
\item \textbf{从渲染到理解}：3DGS~最初为新视角合成而设计，本章展示了如何将高斯基元从纯粹的颜色载体扩展为语义、语言和实例信息的载体，实现从"看得见"到"看得懂"的跃迁。
\item \textbf{2D-3D 语义桥接}：几乎所有场景理解方法的核心挑战都是如何将~2D~基础模型（SAM、CLIP、DINOv2）的语义能力可靠地提升到~3D~空间，同时解决多视角不一致、边界模糊和特征模糊等固有难题。
\item \textbf{实时性与通用性的张力}：部分方法追求逐场景优化的高精度（如~LangSplat~系列），另一些则探索前馈式通用化方案（如~OVGaussian、SLGaussian），两种路线各有优劣。
\item \textbf{从像素到基元}：占用网格、变化检测和场景图等工作表明，场景理解不一定要在像素层面完成。在高斯基元空间、体素空间或图结构上直接操作，往往能获得天然的多视角一致性和更高的计算效率。
\end{itemize}

当前该领域最突出的开放问题包括：

\begin{enumerate}
\item \textbf{动态场景的实时语义理解}：现有方法主要面向静态场景，如何在动态~4D~场景中实现在线、实时、开放词汇的语义理解仍是未解之题
\item \textbf{跨场景泛化}：大多数高斯基语义方法需要逐场景优化，无法直接迁移到新场景。如何学习真正的通用~3D~语义表示是通向大规模部署的关键
\item \textbf{细粒度语义与层次化理解}：当前方法主要关注对象级理解，对部件级（part-level）、材质级和属性级的细粒度语义支持仍显不足
\item \textbf{语义与几何的耦合优化}：语义信息能否反哺几何重建质量（而非仅作为下游任务）是一个值得探索的方向，OmniMap~和~SPHERE~已展示了初步尝试
\item \textbf{标准化的~3D~评估协议}：当前~3D~语义理解领域缺乏统一的~3D~空间评估协议，多数方法仍在~2D~渲染空间评估，SceneSplat++~和~Holi-Spatial~等数据集正在推动这一方向的改进
\end{enumerate}

\end{futurebox}
